% %%%%%%%%%%%%%%%%%%%%%%%%%%%%%%%%%%%%%%%%%%%%%%%%%%%%%%%%%%%%%%%%%%%%%%% %
%                                                                         %
% The Project Gutenberg EBook of Elementary Illustrations of the          %
% Differential and Integral Calculus, by Augustus De Morgan               %
%                                                                         %
% This eBook is for the use of anyone anywhere at no cost and with        %
% almost no restrictions whatsoever.  You may copy it, give it away or    %
% re-use it under the terms of the Project Gutenberg License included     %
% with this eBook or online at www.gutenberg.net                          %
%                                                                         %
%                                                                         %
% Title: Elementary Illustrations of the Differential and Integral Calculus
%                                                                         %
% Author: Augustus De Morgan                                              %
%                                                                         %
% Release Date: March 3, 2012 [EBook #39041]                              %
% Most recently updated: June 11, 2021                         %
%                                                                         %
% Language: English                                                       %
%                                                                         %
% Character set encoding: UTF-8                                           %
%                                                                         %
% *** START OF THIS PROJECT GUTENBERG EBOOK DIFFERENTIAL, INTEGRAL CALCULUS ***
%                                                                         %
% %%%%%%%%%%%%%%%%%%%%%%%%%%%%%%%%%%%%%%%%%%%%%%%%%%%%%%%%%%%%%%%%%%%%%%% %

\def\ebook{39041}
%%%%%%%%%%%%%%%%%%%%%%%%%%%%%%%%%%%%%%%%%%%%%%%%%%%%%%%%%%%%%%%%%%%%%%
%%                                                                  %%
%% Packages and substitutions:                                      %%
%%                                                                  %%
%% book:     Required.                                              %%
%% inputenc: Latin-1 text encoding. Required.                       %%
%%                                                                  %%
%% ifthen:   Logical conditionals. Required.                        %%
%%                                                                  %%
%% amsmath:  AMS mathematics enhancements. Required.                %%
%% amssymb:  Additional mathematical symbols. Required.             %%
%%                                                                  %%
%% alltt:    Fixed-width font environment. Required.                %%
%%                                                                  %%
%% indentfirst: Indent first paragraph of each section. Optional.   %%
%% yfonts:   Gothic font on title page. Optional.                   %%
%%                                                                  %%
%% footmisc: Start footnote numbering on each page. Required.       %%
%% icomma:   No space after commas in math mode. Optional.          %%
%%                                                                  %%
%% multicol: Multicolumn environment for index. Required.           %%
%% makeidx:  Indexing. Required.                                    %%
%%                                                                  %%
%% array:    Extended array facilities. Required.                   %%
%%                                                                  %%
%% graphicx: Standard interface for graphics inclusion. Required.   %%
%%                                                                  %%
%% calc:     Length calculations. Required.                         %%
%%                                                                  %%
%% fancyhdr: Enhanced running headers and footers. Required.        %%
%%                                                                  %%
%% geometry: Enhanced page layout package. Required.                %%
%% hyperref: Hypertext embellishments for pdf output. Required.     %%
%%                                                                  %%
%%                                                                  %%
%% Producer's Comments:                                             %%
%%                                                                  %%
%%   OCR text for this ebook was obtained on January 6, 2012, from  %%
%%   http://www.archive.org/details/elementaryillust00demouoft.     %%
%%                                                                  %%
%%   Minor changes to the original are noted in this file in three  %%
%%   ways:                                                          %%
%%     1. \Typo{}{} for typographical corrections, showing original %%
%%        and replacement text side-by-side.                        %%
%%     2. \Chg{}{} and \Add{}, for inconsistent/missing punctuation,%%
%%        italicization, and capitalization.                        %%
%%     3. [** TN: Note]s for lengthier or stylistic comments.       %%
%%                                                                  %%
%%                                                                  %%
%% Compilation Flags:                                               %%
%%                                                                  %%
%%   The following behavior may be controlled by boolean flags.     %%
%%                                                                  %%
%%   ForPrinting (false by default):                                %%
%%   If false, compile a screen optimized file (one-sided layout,   %%
%%   blue hyperlinks). If true, print-optimized PDF file: Larger    %%
%%   text block, two-sided layout, black hyperlinks.                %%
%%                                                                  %%
%%                                                                  %%
%% PDF pages: 194 (if ForPrinting set to false)                     %%
%% PDF page size: 4.5 x 6.5" (non-standard)                         %%
%%                                                                  %%
%% Summary of log file:                                             %%
%% * Five overfull hboxes, three underfull hboxes.                  %%
%%                                                                  %%
%% Compile History:                                                 %%
%%                                                                  %%
%% February, 2012: (Andrew D. Hwang)                                %%
%%             texlive2007, GNU/Linux                               %%
%%                                                                  %%
%% Command block:                                                   %%
%%                                                                  %%
%%     pdflatex x4                                                  %%
%%     makeindex                                                    %%
%%     pdflatex x3                                                  %%
%%                                                                  %%
%%                                                                  %%
%% March 2012: pglatex.                                             %%
%%   Compile this project with:                                     %%
%%   pdflatex 39041-t.tex ..... FOUR times                          %%
%%   makeindex 39041-t.idx                                          %%
%%   pdflatex 39041-t.tex ..... THREE times                         %%
%%                                                                  %%
%%   pdfTeX, Version 3.1415926-1.40.10 (TeX Live 2009/Debian)       %%
%%                                                                  %%
%%%%%%%%%%%%%%%%%%%%%%%%%%%%%%%%%%%%%%%%%%%%%%%%%%%%%%%%%%%%%%%%%%%%%%
\listfiles
\documentclass[12pt]{book}[2005/09/16]

%%%%%%%%%%%%%%%%%%%%%%%%%%%%% PACKAGES %%%%%%%%%%%%%%%%%%%%%%%%%%%%%%%
\usepackage[utf8]{inputenc}[2006/05/05]

\usepackage{ifthen}[2001/05/26]  %% Logical conditionals

\usepackage{amsmath}[2000/07/18] %% Displayed equations
\usepackage{amssymb}[2002/01/22] %% and additional symbols

\usepackage{alltt}[1997/06/16]   %% boilerplate, credits, license

\IfFileExists{indentfirst.sty}{%
  \usepackage{indentfirst}[1995/11/23]
}{}

\IfFileExists{yfonts.sty}{%
  \usepackage{yfonts}[2003/01/08]
}{%
  \providecommand{\textgoth}[1]{#1}%
}
\usepackage[perpage,symbol]{footmisc}[2005/03/17]
\IfFileExists{icomma.sty}{%
  \usepackage{icomma}
}{}

\usepackage{multicol}[2006/05/18]
\usepackage{makeidx}[2000/03/29]

\usepackage{array}[2005/08/23]

\usepackage{graphicx}[1999/02/16]%% For diagrams

\usepackage{calc}[2005/08/06]

\usepackage{fancyhdr} %% For running heads

%%%%%%%%%%%%%%%%%%%%%%%%%%%%%%%%%%%%%%%%%%%%%%%%%%%%%%%%%%%%%%%%%
%%%% Interlude:  Set up PRINTING (default) or SCREEN VIEWING %%%%
%%%%%%%%%%%%%%%%%%%%%%%%%%%%%%%%%%%%%%%%%%%%%%%%%%%%%%%%%%%%%%%%%

% ForPrinting=true                     false (default)
% Asymmetric margins                   Symmetric margins
% 1 : 1.6 text block aspect ratio      3 : 4 text block aspect ratio
% Black hyperlinks                     Blue hyperlinks
% Start major marker pages recto       No blank verso pages
%
\newboolean{ForPrinting}

%% UNCOMMENT the next line for a PRINT-OPTIMIZED VERSION of the text %%
%\setboolean{ForPrinting}{true}

%% Initialize values to ForPrinting=false
\newcommand{\Margins}{hmarginratio=1:1}     % Symmetric margins
\newcommand{\HLinkColor}{blue}              % Hyperlink color
\newcommand{\PDFPageLayout}{SinglePage}
\newcommand{\TransNote}{Transcriber's Note}
\newcommand{\TransNoteCommon}{%
  The camera-quality files for this public-domain ebook may be
  downloaded \textit{gratis} at
  \begin{center}
    \texttt{www.gutenberg.org/ebooks/\ebook}.
  \end{center}

  This ebook was produced using scanned images and OCR text generously
  provided by the University of Toronto Gerstein Library through the
  Internet Archive.
  \bigskip

  Punctuation in displayed equations has been regularized, and clear
  typographical errors have been changed. Aside from this, every
  effort has been made to preserve the phrasing and punctuation of the
  original.
  \bigskip
}

\newcommand{\TransNoteText}{%
  \TransNoteCommon

  This PDF file is optimized for screen viewing, but may be recompiled
  for printing. Please consult the preamble of the \LaTeX\ source file
  for instructions and other particulars.
}
%% Re-set if ForPrinting=true
\ifthenelse{\boolean{ForPrinting}}{%
  \renewcommand{\Margins}{hmarginratio=2:3} % Asymmetric margins
  \renewcommand{\HLinkColor}{black}         % Hyperlink color
  \renewcommand{\PDFPageLayout}{TwoPageRight}
  \renewcommand{\TransNote}{Transcriber's Note}
  \renewcommand{\TransNoteText}{%
    \TransNoteCommon

    This PDF file is optimized for printing, but may be recompiled for
    screen viewing. Please consult the preamble of the \LaTeX\ source
    file for instructions and other particulars.
  }
}{% If ForPrinting=false, don't skip to recto
  \renewcommand{\cleardoublepage}{\clearpage}
}
%%%%%%%%%%%%%%%%%%%%%%%%%%%%%%%%%%%%%%%%%%%%%%%%%%%%%%%%%%%%%%%%%
%%%%  End of PRINTING/SCREEN VIEWING code; back to packages  %%%%
%%%%%%%%%%%%%%%%%%%%%%%%%%%%%%%%%%%%%%%%%%%%%%%%%%%%%%%%%%%%%%%%%

\ifthenelse{\boolean{ForPrinting}}{%
  \setlength{\paperwidth}{8.5in}%
  \setlength{\paperheight}{11in}%
% 1:1.6
  \usepackage[body={5in,8in},\Margins]{geometry}[2002/07/08]
}{%
  \setlength{\paperwidth}{4.5in}%
  \setlength{\paperheight}{6.5in}%
  \raggedbottom
% 3:4
  \usepackage[body={4.25in,5.6in},\Margins,includeheadfoot]{geometry}[2002/07/08]
}

\providecommand{\ebook}{00000}    % Overridden during white-washing
\usepackage[pdftex,
  hyperref,
  hyperfootnotes=false,
  pdftitle={The Project Gutenberg eBook \#\ebook: Elementary illustrations of the differential and integral calculus.},
  pdfauthor={Augustus De Morgan},
  pdfkeywords={University of Toronto, The Internet Archive, Andrew D. Hwang},
  pdfstartview=Fit,    % default value
  pdfstartpage=1,      % default value
  pdfpagemode=UseNone, % default value
  bookmarks=true,      % default value
  linktocpage=false,   % default value
  pdfpagelayout=\PDFPageLayout,
  pdfdisplaydoctitle,
  pdfpagelabels=true,
  bookmarksopen=true,
  bookmarksopenlevel=0,
  colorlinks=true,
  linkcolor=\HLinkColor]{hyperref}[2007/02/07]


%% Fixed-width environment to format PG boilerplate %%
\newenvironment{PGtext}{%
\begin{alltt}
\fontsize{8.1}{10}\ttfamily\selectfont}%
{\end{alltt}}

% Errors found during digitization
\newcommand{\Typo}[2]{#2}

% Stylistic changes made for consistency
\newcommand{\Chg}[2]{#2}
%\newcommand{\Chg}[2]{#1} % Use this to revert inconsistencies in the original
\newcommand{\Add}[1]{\Chg{}{#1}}


%% Miscellaneous global parameters %%
% No hrule in page header
\renewcommand{\headrulewidth}{0pt}

% Loosen spacing
\setlength{\emergencystretch}{1.5em}
\newcommand{\Loosen}{\spaceskip 0.375em plus 0.75em minus 0.25em}

% Scratch pad for length calculations
\newlength{\TmpLen}

\setlength{\extrarowheight}{12pt}
\renewcommand{\arraystretch}{1.2}
\newcommand{\ArrayCompress}[1][0.6]{\renewcommand{\arraystretch}{#1}}

%% Running heads %%
\newcommand{\FlushRunningHeads}{\clearpage\fancyhf{}}
\newcommand{\InitRunningHeads}{%
  \setlength{\headheight}{15pt}
  \pagestyle{fancy}
  \thispagestyle{empty}
  \ifthenelse{\boolean{ForPrinting}}
             {\fancyhead[RO,LE]{\thepage}}
             {\fancyhead[R]{\thepage}}
}

% Uniform style for running heads
\newcommand{\RHeads}[1]{\textsc{\MakeLowercase{#1}}}

\newcommand{\SetRunningHeads}[2][]{%
  \ifthenelse{\equal{#1}{}}{%
    \fancyhead[C]{\RHeads{#2}}%
  }{%
    \fancyhead[CE]{\RHeads{#1}}%
    \fancyhead[CO]{\RHeads{#2}}%
  }%
}

\newcommand{\BookMark}[2]{\phantomsection\pdfbookmark[#1]{#2}{#2}}

%% Major document divisions %%
\newcommand{\PGBoilerPlate}{%
  \pagenumbering{Alph}
  \pagestyle{empty}
  \BookMark{0}{PG Boilerplate.}
}
\newcommand{\FrontMatter}{%
  \cleardoublepage
  \frontmatter
  \BookMark{-1}{Front Matter.}
}
\newcommand{\MainMatter}{%
  \FlushRunningHeads
  \InitRunningHeads
  \SetRunningHeads[Elementary Illustrations of]{The Differential and Integral Calculus.}
  \mainmatter
  \BookMark{-1}{Main Matter.}
}
\newcommand{\BackMatter}{%
  \FlushRunningHeads
  \InitRunningHeads
  \backmatter
  \BookMark{-1}{Back Matter.}
}
\newcommand{\PGLicense}{%
  \FlushRunningHeads
  \pagenumbering{Roman}
  \InitRunningHeads
  \BookMark{-1}{PG License.}
  \SetRunningHeads{License.}
}

%% ToC formatting %%
\AtBeginDocument{\renewcommand{\contentsname}%
  {\protect\thispagestyle{empty}%
    \vspace*{-3.5\baselineskip}
    \protect\Section{Contents.}
    \vspace*{-1.5\baselineskip}}
}

\newcommand{\TableofContents}{%
  \FlushRunningHeads
  \InitRunningHeads
  \SetRunningHeads{Contents.}
  \BookMark{0}{Contents.}
  \tableofcontents
}

% For internal use, to determine if we need the Sect./Page line
\newcommand{\ToCAnchor}{}

% \ToCLine{SecNo}{Title}
\newcommand{\ToCLine}[2]{%
  \label{toc:#1}%
  \ifthenelse{\not\equal{\pageref{toc:#1}}{\ToCAnchor}}{%
    \renewcommand{\ToCAnchor}{\pageref{toc:#1}}%
    \noindent\makebox[\textwidth][r]{\scriptsize \hfill PAGE}\\
  }{}%
  \settowidth{\TmpLen}{999}%
  \noindent\strut\parbox[b]{\textwidth-\TmpLen}{\small%
     \hangindent2em#2\dotfill}%
  \makebox[\TmpLen][r]{\pageref{section:#1}}%
  \smallskip
}

%% Sectional units %%
% For internal use by \Tag and \Eq
\newcounter{SecNo}
\newcommand{\Section}[2][]{
  \subsection*{\centering\large\normalfont\MakeUppercase{#2}}
  \refstepcounter{SecNo}
  \label{section:\theSecNo}%
  \ifthenelse{\not\equal{#1}{}}{%
    \addtocontents{toc}{\protect\ToCLine{\theSecNo}{#1}}%
  }{}
}

\newcommand{\SubSectHead}[1]{
  \subsubsection*{\centering\footnotesize\normalfont\MakeUppercase{#1}}
}

\newcommand{\Subsection}[2][]{
  \SubSectHead{#2}
  \refstepcounter{SecNo}
  \label{section:\theSecNo}%
  \ifthenelse{\equal{#1}{}}{%
    \addtocontents{toc}{\protect\ToCLine{\theSecNo}{#2}}%
  }{%
    \addtocontents{toc}{\protect\ToCLine{\theSecNo}{#1}}%
  }
}

\newcommand{\Signature}[2]{%
  \medskip
  \null\hfill#1\hspace{\parindent} \\
  \hspace*{\parindent}#2
}

%% Index formatting
\makeindex
\makeatletter
\renewcommand{\@idxitem}{\par\hangindent 30\p@\global\let\idxbrk\nobreak}
\renewcommand\subitem{\idxbrk\@idxitem \hspace*{12\p@}\let\idxbrk\relax}
\renewcommand{\indexspace}{\par\penalty-3000 \vskip 10pt plus5pt minus3pt\relax}

\renewenvironment{theindex}{%
  \setlength\columnseprule{0.5pt}\setlength\columnsep{18pt}%
  \FlushRunningHeads
  \InitRunningHeads
  \SetRunningHeads{Index}
  \BookMark{0}{Index}
  \begin{multicols}{2}[{\Section[Index]{Index.}\small}]%
    \setlength\parindent{0pt}\setlength\parskip{0pt plus 0.3pt}%
    \thispagestyle{empty}\let\item\@idxitem\raggedright%
  }{%
  \end{multicols}\FlushRunningHeads
}
\makeatother

\newcommand{\EtSeq}[1]{\hyperpage{#1}~et~seq.}

%% Bibliography %%
\newlength{\ParIndent}
\settowidth{\ParIndent}{2em}
% Template for definitions, theorems, examples
\newenvironment{Book}{\medskip\par\noindent\hangindent\ParIndent\ignorespaces}{\par}

\newenvironment{Descrip}{%
  \begin{list}{}{%
      \setlength{\leftmargin}{1.5\ParIndent}%
      \setlength{\itemindent}{\ParIndent}%
      \setlength{\listparindent}{\ParIndent}%
      \setlength{\topsep}{4pt}%
      \item\small\ignorespaces%
  }%
}{%
\end{list}}

\newcommand{\BibSect}[1]{%
  \subsubsection*{\centering\normalsize\normalfont\MakeUppercase{#1}}
}
\newcommand{\BibSubsect}[1]{%
  \subsubsection*{\centering\footnotesize\normalfont\MakeUppercase{#1}}
}

%% Diagrams %%
\newcommand{\Graphic}[2]{%
%  \phantomsection\label{fig:#2}%
  \includegraphics[width=#1]{./images/fig#2.png}%
}
% \Figure[width]{figure number}{file}
\newcommand{\DefWidth}{4.25in}% Default figure width
\newcommand{\Figure}[2][]{%
  \begin{figure}[hbt!]
    \Graphic{\textwidth}{#2}%
    \ifthenelse{\equal{#1}{}}{%
      \phantomsection\label{fig:#2}%
    }{}%
  \end{figure}%
}

% Figure labels
\newcommand{\Fig}[2][Fig.]{\hyperref[fig:#2]{#1~{\upshape #2}}}

\newcommand{\Pagelabel}[1]{\phantomsection\label{page:#1}}

% Page separators
\newcommand{\PageSep}[1]{\Pagelabel{#1}\ignorespaces}
\newcommand{\PageNo}[1]{\pageref{page:#1}}
\newcommand{\PageRef}[2][page]{\hyperref[page:#2]{#1~\pageref*{page:#2}}}
\newcommand{\PageRefs}[3][pages]{%
  \ifthenelse{\equal{\pageref*{page:#2}}{\pageref*{page:#3}}}{%
    \hyperref[page:#2]{page~\pageref*{page:#2}}%
  }{% Else
    #1 \hyperref[page:#2]{\pageref*{page:#2}}--\hyperref[page:#3]{\pageref*{page:#3}}%
  }%
}

% Miscellaneous textual conveniences (N.B. \emph, not \textit)
\newcommand{\eg}{\emph{e.g.}}
\newcommand{\ie}{\emph{i.e.}}
\newcommand{\QED}{\quad\textsc{q.e.d.}}
\renewcommand{\etc}{\text{etc}}
\newcommand{\Ed}{\textit{Ed}}

\newcommand{\ia}{\textit{a}}
\newcommand{\ib}{\textit{b}}

\renewcommand{\th}{\textsuperscript{th}}

\renewcommand{\(}{{\upshape(}}
\renewcommand{\)}{{\upshape)}}

\newcommand{\First}[1]{\textsc{#1}}

\newcommand{\Title}[1]{\textit{#1}}
\newcommand{\Author}[1]{\textsc{#1}}
\newcommand{\Translator}[1]{\textit{#1}}

%% \Price macro accepts numerals, letters, ".", and " ".
%% Letters are italicized, spaces made non-breaking.
\makeatletter
\newif\ifins@rtsp@ce
\def\Price#1{{\ins@rtsp@cefalse\Pr@ce #1\null}}
\def\Pr@ce#1{%
  \ifx#1\null\let\next\empty
  \else\let\next\Pr@ce
  \ifins@rtsp@ce~\ins@rtsp@cefalse\fi
  \if.#1.\ins@rtsp@cetrue
  \else\ifcat#1a\textit{#1}%
  \else{#1}%
  \fi\fi\fi\next}
\makeatother

%% Miscellaneous mathematical formatting %%
\newcommand{\dd}{\partial}

\DeclareMathOperator{\arc}{arc}
\DeclareMathOperator{\chord}{chord}

%[** TN: Original uses something like these]
%\newcommand{\ux}{u'}
%\newcommand{\uy}{u\raisebox{-1.25ex}{$'$}}
\newcommand{\ux}{u_{x}'}
\newcommand{\uy}{u_{y}'}

\renewcommand{\phi}{\varphi}

\newcommand{\efrac}[2]{\frac{#1}{#2}}

\DeclareInputMath{176}{{}^{\circ}}
\DeclareInputMath{183}{\cdot}

\newcommand{\Z}{\phantom{0}}
\newcommand{\DotRow}[1]{\multispan{#1}{\dotfill}}
\newcommand{\Ibar}{\mathop{\vphantom{|}}}% ** Integral bar

\newcommand{\THeadSize}{\scriptsize}
\newcommand{\ColHead}[2][]{%
  \ifthenelse{\equal{#1}{}}{%
    \settowidth{\TmpLen}{\THeadSize #2}%
  }{%
    \settowidth{\TmpLen}{\THeadSize #1}%
  }
  \parbox[c]{\TmpLen}{\centering\THeadSize #2\medskip}%
}

\newcommand{\Low}[1]{%
  \setlength{\TmpLen}{-1.5ex}
  \smash[b]{\raisebox{\TmpLen}{\ensuremath{#1}}}
}

% \PadTo[alignment]{width text}{visible text}
\newcommand{\PadTo}[3][c]{%
  \settowidth{\TmpLen}{$#2$}%
  \makebox[\TmpLen][#1]{$#3$}%
}
\newcommand{\PadTxt}[3][c]{%
  \settowidth{\TmpLen}{\text{#2}}%
  \makebox[\TmpLen][#1]{#3}%
}

% Cross-ref-able equation tags
\newcommand{\Tag}[1]{%
  \phantomsection\label{eqn:\theSecNo.#1}\tag*{\ensuremath{#1}}%
}
\newcommand{\Eq}[2][\theSecNo]{%
  \hyperref[eqn:#1.#2]{\ensuremath{#2}}%
}

%% Upright capital letters in math mode
\DeclareMathSymbol{A}{\mathalpha}{operators}{`A}
\DeclareMathSymbol{B}{\mathalpha}{operators}{`B}
\DeclareMathSymbol{C}{\mathalpha}{operators}{`C}
\DeclareMathSymbol{D}{\mathalpha}{operators}{`D}
\DeclareMathSymbol{E}{\mathalpha}{operators}{`E}
\DeclareMathSymbol{F}{\mathalpha}{operators}{`F}
\DeclareMathSymbol{G}{\mathalpha}{operators}{`G}
\DeclareMathSymbol{H}{\mathalpha}{operators}{`H}
\DeclareMathSymbol{I}{\mathalpha}{operators}{`I}
\DeclareMathSymbol{J}{\mathalpha}{operators}{`J}
\DeclareMathSymbol{K}{\mathalpha}{operators}{`K}
\DeclareMathSymbol{L}{\mathalpha}{operators}{`L}
\DeclareMathSymbol{M}{\mathalpha}{operators}{`M}
\DeclareMathSymbol{N}{\mathalpha}{operators}{`N}
\DeclareMathSymbol{O}{\mathalpha}{operators}{`O}
\DeclareMathSymbol{P}{\mathalpha}{operators}{`P}
\DeclareMathSymbol{Q}{\mathalpha}{operators}{`Q}
\DeclareMathSymbol{R}{\mathalpha}{operators}{`R}
\DeclareMathSymbol{S}{\mathalpha}{operators}{`S}
\DeclareMathSymbol{T}{\mathalpha}{operators}{`T}
\DeclareMathSymbol{U}{\mathalpha}{operators}{`U}
\DeclareMathSymbol{V}{\mathalpha}{operators}{`V}
\DeclareMathSymbol{W}{\mathalpha}{operators}{`W}
\DeclareMathSymbol{X}{\mathalpha}{operators}{`X}
\DeclareMathSymbol{Y}{\mathalpha}{operators}{`Y}
\DeclareMathSymbol{Z}{\mathalpha}{operators}{`Z}

%%%%%%%%%%%%%%%%%%%%%%%% START OF DOCUMENT %%%%%%%%%%%%%%%%%%%%%%%%%%
\begin{document}
%% PG BOILERPLATE %%
\PGBoilerPlate
\begin{center}
\begin{minipage}{\textwidth}
\small
\begin{PGtext}
The Project Gutenberg EBook of Elementary Illustrations of the
Differential and Integral Calculus, by Augustus De Morgan

This eBook is for the use of anyone anywhere at no cost and with
almost no restrictions whatsoever.  You may copy it, give it away or
re-use it under the terms of the Project Gutenberg License included
with this eBook or online at www.gutenberg.net


Title: Elementary Illustrations of the Differential and Integral Calculus

Author: Augustus De Morgan

Release Date: March 3, 2012 [EBook #39041]
Most recently updated: June 11, 2021

Language: English

Character set encoding: UTF-8     

*** START OF THIS PROJECT GUTENBERG EBOOK DIFFERENTIAL, INTEGRAL CALCULUS ***
\end{PGtext}
\end{minipage}
\end{center}
\newpage
%% Credits and transcriber's note %%
\begin{center}
\begin{minipage}{\textwidth}
\begin{PGtext}
Produced by Andrew D. Hwang.
\end{PGtext}
\end{minipage}
\vfill
\end{center}

\begin{minipage}{0.85\textwidth}
\small
\BookMark{0}{Transcriber's Note.}
\subsection*{\centering\normalfont\scshape%
\normalsize\MakeLowercase{\TransNote}}%

\raggedright
\TransNoteText
\end{minipage}
%%%%%%%%%%%%%%%%%%%%%%%%%%% FRONT MATTER %%%%%%%%%%%%%%%%%%%%%%%%%%
\PageSep{ii}
\FrontMatter


\Section{In the Same Series.}

\Book{ON THE STUDY AND DIFFICULTIES OF MATHEMATICS.}
By \Author{Augustus De~Morgan}. Entirely new edition,
with portrait of the author, index, and annotations,
bibliographies of modern works on algebra, the philosophy
of mathematics, pan-geometry,~etc. Pp.,~288. Cloth, \$1.25
net~(\Price{5s.}).

\Book{LECTURES ON ELEMENTARY MATHEMATICS.} By
\Author{Joseph Louis Lagrange}. Translated from the French by
\Translator{Thomas~J. McCormack}. With photogravure portrait of
Lagrange, notes, biography, marginal analyses,~etc. Only
separate edition in French or English, Pages,~172. Cloth,
\$1.00 net~(\Price{5s.}).

\Book{ELEMENTARY ILLUSTRATIONS OF THE DIFFERENTIAL
AND INTEGRAL CALCULUS.} By \Author{Augustus De~Morgan}.
New reprint edition. With sub-headings, and
a brief bibliography of English, French, and German text-books
of the Calculus. Pp.,~144. Price, \$1.00 net~(\Price{5s.}).

\Book{MATHEMATICAL ESSAYS AND RECREATIONS.} By
\Author{Hermann Schubert}, Professor of Mathematics in the
Johanneum, Hamburg, Germany. Translated from the
German by \Translator{Thomas~J. McCormack}. Containing essays on
the Notion and Definition of Number, Monism in Arithmetic,
On the Nature of Mathematical Knowledge, The
Magic Square, The Fourth Dimension, The Squaring of
the Circle. Pages,~149. Cuts,~37. Price, Cloth,~75c net~(\Price{3s. 6d.}).

\Book{HISTORY OF ELEMENTARY MATHEMATICS.} By \Author{Dr.\
Karl Fink}, late Professor in Tübingen. Translated from
the German by Prof.\ \Translator{Wooster Woodruff Beman} and Prof.\
\Translator{David Eugene Smith}. (Nearly Ready.)
\vfill
\begin{center}
\small
THE OPEN COURT PUBLISHING CO. \\
\footnotesize
324 DEARBORN ST., CHICAGO.
\end{center}
\PageSep{iii}
\cleardoublepage
%[** Title page]
\begin{center}
ELEMENTARY ILLUSTRATIONS
\vfil
\footnotesize OF THE
\vfil
\textsc{\LARGE Differential and Integral \\[4pt]
Calculus}
\vfil\vfil\vfil

\footnotesize BY \\
\normalsize AUGUSTUS DE MORGAN
\vfil\vfil\vfil

\textit{\small NEW EDITION}
\vfil\vfil\vfil\vfil

\footnotesize CHICAGO \\
THE OPEN COURT PUBLISHING COMPANY \\
\scriptsize FOR SALE BY \\
\footnotesize\textsc{Kegan Paul, Trench, Trübner \&~Co., Ltd., London} \\
1899
\end{center}
\PageSep{iv}
% [Blank page]
\PageSep{v}


\Section{Editor's Preface.}

\First{The} publication of the present reprint of De~Morgan's \Title{Elementary
Illustrations of the Differential and Integral Calculus}
forms, quite independently of its interest to professional
students of mathematics, an integral portion of the general educational
plan which the Open Court Publishing Company has been
systematically pursuing since its inception,---which is the dissemination
among the public at large of sound views of science and of
an adequate and correct appreciation of the methods by which
truth generally is reached. Of these methods, mathematics, by
its simplicity, has always formed the type and ideal, and it is
nothing less than imperative that its ways of procedure, both in
the discovery of new truth and in the demonstration of the necessity
and universality of old truth, should be laid at the foundation
of every philosophical education. The greatest achievements in
the history of thought---Plato, Descartes, Kant---are associated
with the recognition of this principle.

But it is precisely mathematics, and the pure sciences generally,
from which the general educated public and independent
students have been debarred, and into which they have only rarely
attained more than a very meagre insight. The reason of this is
twofold. In the first place, the ascendant and consecutive character
of mathematical knowledge renders its results absolutely unsusceptible
of presentation to persons who are unacquainted with
what has gone before, and so necessitates on the part of its devotees
a thorough and patient exploration of the field from the very
beginning, as distinguished from those sciences which may, so to
speak, be begun at the end, and which are consequently cultivated
with the greatest zeal. The second reason is that, partly through
the exigencies of academic instruction, but mainly through the
martinet traditions of antiquity and the influence of mediæval
\PageSep{vi}
logic-mongers, the great bulk of the elementary text-books of
mathematics have unconsciously assumed a very repellent form,---something
similar to what is termed in the theory of protective
mimicry in biology ``the terrifying form.'' And it is mainly to
this formidableness and touch-me-not character of exterior, concealing
withal a harmless body, that the undue neglect of typical
mathematical studies is to be attributed.

To this class of books the present work forms a notable exception.
It was originally issued as numbers 135 and 140 of the
Library of Useful Knowledge (1832), and is usually bound up with
De~Morgan's large \Title{Treatise on the Differential and Integral
Calculus} (1842). Its style is fluent and familiar; the treatment
continuous and undogmatic. The main difficulties which encompass
the early study of the Calculus are analysed and discussed in
connexion with practical and historical illustrations which in point
of simplicity and clearness leave little to be desired. No one who
will read the book through, pencil in hand, will rise from its perusal
without a clear perception of the aim and the simpler fundamental
principles of the Calculus, or without finding that the profounder
study of the science in the more advanced and more
methodical treatises has been greatly facilitated.

The book has been reprinted substantially as it stood in its
original form; but the typography has been greatly improved, and
in order to render the subject-matter more synoptic in form and
more capable of survey, the text has been re-paragraphed and a
great number of descriptive sub-headings have been introduced, a
list of which will be found in the Contents of the book. An index
also has been added.

Persons desirous of continuing their studies in this branch of
mathematics, will find at the end of the text a bibliography of the
principal English, French, and German works on the subject, as
well as of the main Collections of Examples. From the information
there given, they may be able to select what will suit their
special needs.

\Signature{Thomas J. McCormack.}
{\textsc{La Salle}, Ill., August, 1899.}
\PageSep{vii}


\TableofContents
\iffalse
CONTENTS:

PAGE

On the Ratio or Proportion of Two Magnitudes 2
On the Ratio of Magnitudes that Vanish Together.... 4
On the Ratios of Continuously Increasing or Decreasing Quantities 7
The Notion of Infinitely Small Quantities 11
On Functions 14
Infinite Series 15
Convergent and Divergent Series 17
Taylor's Theorem. Derived Functions 19
Differential Coefficients 22
The Notation of the Differential Calculus 25
Algebraical Geometry.... 29
On the Connexion of the Signs of Algebraical and the Directions
  of Geometrical Magnitudes 31
The Drawing of a Tangent to a Curve 36
Rational Explanation of the Language of Leibnitz.... 38
Orders of Infinity 42
A Geometrical Illustration: Limit of the Intersections of Two
  Coinciding Straight Lines 45
The Same Problem Solved by the Principles of Leibnitz. . 48
An Illustration from Dynamics; Velocity, Acceleration, etc.. 52
Simple Harmonic Motion 57
The Method of Fluxions 60
Accelerated Motion 60
Limiting Ratios of Magnitudes that Increase Without Limit. 65
Recapitulation of Results Reached in the Theory of Functions. 74
Approximations by the Differential Calculus 74
Solution of Equations by the Differential Calculus.... 77
Partial and Total Differentials 78
\PageSep{viii}
Application of the Theorem for Total Differentials to the
  Determination of Total Resultant Errors 84
Rules for Differentiation 85
Illustration of the Rules for Differentiation 86
Differential Coefficients of Differential Coefficients .... 88
Calculus of Finite Differences. Successive Differentiation . 88
Total and Partial Differential Coefficients. Implicit Differentiation 94
Applications of the Theorem for Implicit Differentiation .. 101
Inverse Functions 102
Implicit Functions 106
Fluxions, and the Idea of Time 110
The Differential Coefficient Considered with Respect to Its
  Magnitude 112
The Integral Calculus 115
Connexion of the Integral with the Differential Calculus.. 120
Nature of Integration 122
Determination of Curvilinear Areas. The Parabola... 124
Method of Indivisibles 125
Concluding Remarks on the Study of the Calculus.... 132
Bibliography of Standard Text-books and Works of Reference
  on the Calculus 133
Index 143
\fi
\PageSep{1}
\MainMatter


\Section{Differential and Integral Calculus.}

\SubSectHead{Elementary Illustrations.}

\First{The} Differential and Integral Calculus, or, as it
was formerly called in this country [England],
the Doctrine of Fluxions, has always been supposed
to present remarkable obstacles to the beginner. It
is matter of common observation, that any one who
commences this study, even with the best elementary
works, finds himself in the dark as to the real meaning
of the processes which he learns, until, at a certain
stage of his progress, depending upon his capacity,
some accidental combination of his own ideas throws
light upon the subject. The reason of this may be, that
it is usual to introduce him at the same time to new
principles, processes, and symbols, thus preventing
his attention from being exclusively directed to one
new thing at a time. It is our belief that this should
be avoided; and we propose, therefore, to try the experiment,
whether by undertaking the solution of
some problems by common algebraical methods, without
calling for the reception of more than one new
symbol at once, or lessening the immediate evidence
of each investigation by reference to general rules, the
study of more methodical treatises may not be somewhat
\PageSep{2}
facilitated. We would not, nevertheless, that
the student should imagine we can remove all obstacles;
we must introduce notions, the consideration
of which has not hitherto occupied his mind; and
shall therefore consider our object as gained, if we
can succeed in so placing the subject before him, that
two independent difficulties shall never occupy his
mind at once.


\Subsection{On the Ratio or Proportion of Two Magnitudes.}

The ratio or proportion of two magnitudes is best
\index{Proportion|EtSeq}%
\index{Ratio!defined|EtSeq}%
conceived by expressing them in numbers of some
unit when they are commensurable; or, when this is
not the case, the same may still be done as nearly as
we please by means of numbers. Thus, the ratio of
the diagonal of a square to its side is that of $\sqrt{2}$ to~$1$,
which is very nearly that of $14142$ to~$10000$, and is
certainly between this and that of $14143$ to~$10000$.
Again, any ratio, whatever numbers express it, may
be the ratio of two magnitudes, each of which is as
small as we please; by which we mean, that if we
take any given magnitude, however small, such as the
line~$A$, we may find two other lines $B$~and~$C$, each
less than~$A$, whose ratio shall be whatever we please.
Let the given ratio be that of the numbers $m$~and~$n$.
Then, $P$~being a line, $mP$~and~$nP$ are in the proportion
of $m$ to~$n$; and it is evident, that let $m$,~$n$, and~$A$
be what they may, $P$~can be so taken that $mP$~shall be
less than~$A$. This is only saying that $P$~can be taken
less than the $m$\th~part of~$A$, which is obvious, since~$A$,
however small it may be, has its tenth, its hundredth,
its thousandth part,~etc., as certainly as if it were
larger. We are not, therefore, entitled to say that
because two magnitudes are diminished, their ratio is
\PageSep{3}
diminished; it is possible that~$B$, which we will suppose
to be at first a hundredth part of~$C$, may, after
a diminution of both, be its tenth or thousandth, or
may still remain its hundredth, as the following example
will show:
\begin{alignat*}{5}
&C && 3600 && 1800 && 36 && 90 \\
&B && 36   && 1\tfrac{8}{10} && \tfrac{36}{100} && 9 \\
&B &{}={}& \frac{1}{100} C\qquad
 B &{}={}& \frac{1}{1000} C\qquad
 B &{}={}& \frac{1}{100} C\qquad
 B &{}={}& \frac{1}{10} C.
\end{alignat*}
Here the values of $B$~and~$C$ in the second, third, and
fourth column are less than those in the first; nevertheless,
the ratio of $B$ to~$C$ is less in the second column
than it was in the first, remains the same in the
third, and is greater in the fourth.

In estimating the approach to, or departure from
equality, which two magnitudes undergo in consequence
of a change in their values, we must not look
at their differences, but at the proportions which those
differences bear to the whole magnitudes. For example,
if a geometrical figure, two of whose sides are
$3$~and $4$~inches now, be altered in dimensions, so that
the corresponding sides are $100$~and $101$~inches, they
are nearer to equality in the second case than in the
first; because, though the difference is the same in
both, namely one inch, it is one third of the least side
in the first case, and only one hundredth in the second.
This corresponds to the common usage, which
rejects quantities, not merely because they are small,
but because they are small in proportion to those of
which they are considered as parts. Thus, twenty
miles would be a material error in talking of a day's
journey, but would not be considered worth mentioning
in one of three months, and would be called totally
\PageSep{4}
insensible in stating the distance between the
earth and sun. More generally, if in the two quantities
$x$~and~$x + a$, an increase of~$m$ be given to~$x$,
the two resulting quantities $x + m$~and $x + m +  a$ are
nearer to equality as to their \emph{ratio} than $x$~and~$x + a$,
\index{Equality}%
though they continue the same as to their \emph{difference}; for
\index{Differences!arithmetical}%
$\dfrac{x + a}{x} = 1 + \dfrac{a}{x}$ and $\dfrac{x + m + a}{x + m} = 1 + \dfrac{a}{x + m}$ of which
$\dfrac{a}{x + m}$~is less than~$\dfrac{a}{x}$, and therefore $1 + \dfrac{a}{x + m}$ is nearer
to unity than $1 + \dfrac{a}{x}$. In future, when we talk of an
approach towards equality, we mean that the ratio is
made more nearly equal to unity, not that the difference
is more nearly equal to nothing. The second
may follow from the first, but not necessarily; still
less does the first follow from the second.


\Subsection{On the Ratio of Magnitudes that Vanish Together.}

It is conceivable that two magnitudes should decrease
\index{Time, idea of}%
simultaneously,\footnote
  {In introducing the notion of time, we consult only simplicity. It would
  do equally well to write any number of successive values of the two quantities,
  and place them in two columns.}
so as to vanish or become
nothing, together. For example, let a point~$A$ move
on a circle towards a fixed point~$B$. The arc~$AB$ will
then diminish, as also the chord~$AB$, and by bringing
the point~$A$ sufficiently near to~$B$, we may obtain an
arc and its chord, both of which shall be smaller than
a given line, however small this last may be. But
while the magnitudes diminish, we may not assume
either that their ratio increases, diminishes, or remains
the same, for we have shown that a diminution
of two magnitudes is consistent with either of these.
\PageSep{5}
\index{Increase without limit|EtSeq}%
We must, therefore, look to each particular case for
the change, if any, which is made in the ratio by the
diminution of its terms.

Now two suppositions are possible in every increase
or diminution of the ratio, as follows: Let $M$~and~$N$
be two quantities which we suppose in a state
of decrease. The first possible case is that the ratio
of $M$ to~$N$ may decrease without limit, that is, $M$~may
be a smaller fraction of~$N$ after a decrease than it was
before, and a still smaller after a further decrease,
and so on; in such a way, that there is no fraction so
small, to which $\dfrac{M}{N}$~shall not be equal or inferior, if the
decrease of $M$~and~$N$ be carried sufficiently far. As
an instance, form two sets of numbers as in the adjoining
table:
\[
\begin{array}{*{7}{c}}
\PadTxt{Ratio of }{} M \PadTxt{ to~$N$}{}
  & 1 & \dfrac{1}{20} & \dfrac{1}{400} & \dfrac{1}{8000} & \dfrac{1}{160000} & \etc. \\
\PadTxt{Ratio of }{} N \PadTxt{ to~$N$}{}
  & 1 & \dfrac{1}{2} & \dfrac{1}{4} & \dfrac{1}{8} & \dfrac{1}{16} & \etc. \\
\text{Ratio of~$M$ to~$N$}
  & 1 & \dfrac{1}{10} & \dfrac{1}{100} & \dfrac{1}{1000} & \dfrac{1}{10000} & \etc.
\end{array}
\]
Here both $M$~and~$N$ decrease at every step, but $M$~loses
at each step a larger fraction of itself than~$N$,
and their ratio continually diminishes. To show that
this decrease is without limit, observe that $M$~is at
first equal to~$N$, next it is one tenth, then one hundredth,
then one thousandth of~$N$, and so on; by continuing
the values of $M$ and~$N$ according to the same
law, we should arrive at a value of~$M$ which is a
smaller part of~$N$ than any which we choose to name;
for example,~$.000003$. The second value of~$M$ beyond
our table is only one millionth of the corresponding
value of~$N$; the ratio is therefore expressed by~$.000001$
\PageSep{6}
which is less than~$.000003$. In the same law of formation,
the ratio of $N$ to~$M$ is also \emph{increased} without limit.

The second possible case is that in which the ratio
of $M$ to~$N$, though it increases or decreases, does not
increase or decrease without limit, that is, continually
approaches to some ratio, which it never will exactly
reach, however far the diminution of $M$ and~$N$ may
be carried. The following is an example:
\[
\begin{array}{*{9}{c}}
\PadTxt{Ratio of }{} M \PadTxt{ to~$N$}{}
  & 1 & \dfrac{1}{3} & \dfrac{1}{6} & \dfrac{1}{10} & \dfrac{1}{15} & \dfrac{1}{21} & \dfrac{1}{28} & \etc. \\
\PadTxt{Ratio of }{} N \PadTxt{ to~$N$}{}
  & 1 & \dfrac{1}{4} & \dfrac{1}{9} & \dfrac{1}{16} & \dfrac{1}{25} & \dfrac{1}{36} & \dfrac{1}{49} & \etc. \\
\text{Ratio of~$M$ to~$N$}
  & 1 & \dfrac{4}{3} & \dfrac{9}{6} & \dfrac{16}{10} & \dfrac{25}{15} & \dfrac{36}{21} & \dfrac{49}{28} & \etc.
\end{array}
\]
The ratio here increases at each step, for $\dfrac{4}{3}$~is greater
than~$1$, $\dfrac{9}{6}$~than~$\dfrac{4}{3}$, and so on. The difference between
this case and the last is, that the ratio of $M$ to~$N$,
though perpetually increasing, does not increase without
limit; it is never so great as~$2$, though it may be
brought as near to~$2$ as we please.

To show this, observe that in the successive values
of~$M$, the denominator of the second is~$1 + 2$, that of
the third $1 + 2 + 3$, and so on; whence the denominator
of the $x$\th~value of~$M$ is
\[
1 + 2 + 3 + \dots + x,\quad\text{or}\quad \frac{x(x + 1)}{2}\Add{.}
\]
Therefore the $x$\th~value of~$M$ is~$\dfrac{2}{x(x + 1)}$, and it is
evident that the $x$\th~value of~$N$ is~$\dfrac{1}{x^{2}}$, which gives the
$x$\th~value of the ratio $\dfrac{M}{N} = \dfrac{2x^{2}}{x(x + 1)}$, or~$\dfrac{2x}{x + 1}$, or
\PageSep{7}
$\dfrac{x}{x + 1} × 2$. If $x$~be made sufficiently great, $\dfrac{x}{x + 1}$~may
be brought as near as we please to~$1$, since, being
$1 - \dfrac{1}{x + 1}$, it differs from~$1$ by~$\dfrac{1}{x + 1}$, which may be
made as small as we please. But as $\dfrac{x}{x + 1}$, however
great $x$~may be, is always less than~$1$, $\dfrac{2x}{x + 1}$~is always
less than~$2$. Therefore (1)~$\dfrac{M}{N}$~continually increases;
(2)~may be brought as near to~$2$ as we please; (3)~can
never be greater than~$2$. This is what we mean by
saying that $\dfrac{M}{N}$~is an increasing ratio, the limit of
which is~$2$. Similarly of~$\dfrac{N}{M}$, which is the reciprocal
of~$\dfrac{M}{N}$, we may show (1)~that it continually decreases;
(2)~that it can be brought as near as we please to~$\frac{1}{2}$;
(3)~that it can never be less than~$\frac{1}{2}$. This we express
by saying that $\dfrac{N}{M}$~is a decreasing ratio, whose limit
is~$\frac{1}{2}$.


\Subsection{On the Ratios of Continuously Increasing or
Decreasing Quantities.}

To the fractions here introduced, there are intermediate
\index{Continuous quantities|EtSeq}%
\index{Quantities, continuous|EtSeq}%
fractions, which we have not considered.
Thus, in the last instance, $M$~passed from $1$ to~$\frac{1}{2}$ without
any intermediate change. In geometry and mechanics,
it is necessary to consider quantities as
increasing or decreasing \emph{continuously}; that is, a magnitude
does not pass from one value to another without
passing through every intermediate value. Thus
if one point move towards another on a circle, both
the arc and its chord decrease continuously. Let $AB$
\index{Arc and its chord, a continuously decreasing|EtSeq}%
(\Fig{1}) be an arc of a circle, the centre of which is~$O$.
\PageSep{8}
Let $A$ remain fixed, but let $B$, and with it the radius~$OB$,
move towards~$A$, the point~$B$ always remaining
on the circle. At every position of~$B$, suppose
the following figure. Draw $AT$ touching the circle at~$A$,
produce $OB$ to meet~$AT$ in~$T$, draw $BM$~and~$BN$
perpendicular and parallel to~$OA$, and join~$BA$. Bisect
the arc~$AB$ in~$C$, and draw~$OC$ meeting the chord in~$D$
and bisecting it. The right-angled triangles $ODA$
and $BMA$ having a common angle, and also right
angles, are similar, as are also $BOM$ and~$TBN$. If
now we suppose $B$ to move towards~$A$, before $B$
\Figure{1}
reaches~$A$, we shall have the following results: The
arc and chord~$BA$, the lines $BM$,~$MA$, $BT$,~$TN$, the
angles $BOA$,~$COA$,~$MBA$, and~$TBN$, will diminish
without limit; that is, assign a line and an angle,
however small, $B$~can be placed so near to~$A$ that the
lines and angles above alluded to shall be severally
less than the assigned line and angle. Again, $OT$~diminishes
and $OM$~increases, but neither without limit,
for the first is never less, nor the second greater, than
the radius. The angles $OBM$,~$MAB$, and~$BTN$, increase,
but not without limit, each being always less
than the right angle, but capable of being made as
\PageSep{9}
near to it as we please, by bringing~$B$ sufficiently near
to~$A$.

So much for the magnitudes which compose the
figure: we proceed to consider their ratios, premising
that the arc~$AB$ is greater than the chord~$AB$, and
less than $BN + NA$. The triangle~$BMA$ being always
similar to~$ODA$, their sides change always in the same
proportion; and the sides of the first decrease without
limit, which is the case with only one side of the
second. And since $OA$~and~$OD$ differ by~$DC$, which
diminishes without limit as compared with~$OA$, the
ratio $OD ÷ OA$ is an increasing ratio whose limit is~$1$.
But $OD ÷ OA = BM ÷ BA$. We can therefore bring~$B$
so near to~$A$ that $BM$~and~$BA$ shall differ by as
small a fraction of either of them as we please.

To illustrate this result from the trigonometrical
tables, observe that if the radius~$OA$ be the linear
unit, and $\angle BOA = \theta$, $BM$~and~$BA$ are respectively
$\sin\theta$ and $2\sin\frac{1}{2}\theta$. Let $\theta = 1°$; then $\sin\theta = .0174524$
and $2\sin\frac{1}{2}\theta = .0174530$; whence $2\sin\frac{1}{2}\theta ÷ \sin\theta = 1.00003$ very nearly, so that $BM$~differs from~$BA$ by
less than four of its own hundred-thousandth parts.
If $\angle BOA = 4'$, the same ratio is~$1.0000002$, differing
from unity by less than the hundredth part of the
difference in the last example.

Again, since $DA$~diminishes continually and without
limit, which is not the case either with $OD$ or~$OA$,
the ratios $OD ÷ DA$ and $OA ÷ DA$ increase without
limit. These are respectively equal to $BM ÷ MA$
and $BA ÷ MA$; whence it appears that, let a number
be ever so great, $B$~can be brought so near to~$A$, that
$BM$ and $BA$ shall each contain~$MA$ more times than
there are units in that number. Thus if $\angle BOA = 1°$,
$BM ÷ MA = 114.589$ and $BA ÷ MA = 114.593$ very
\PageSep{10}
nearly; that is, $BM$ and $BA$ both contain~$MA$
more than $114$~times. If $\angle BOA = 4'$, $BM ÷ MA = 1718.8732$,
and $BA ÷ MA = 1718.8375$ very nearly;
or $BM$ and $BA$ both contain~$MA$ more than $1718$~times.

No difficulty can arise in conceiving this result, if
the student recollect that the degree of greatness or
smallness of two magnitudes determines nothing as
to their ratio; since every quantity~$N$, however small,
can be divided into as many parts as we please, and
has therefore another small quantity which is its millionth
\Figure[nolabel]{1}
or hundred-millionth part, as certainly as if it
had been greater. There is another instance in the
line~$TN$, which, since $TBN$~is similar to~$BOM$, decreases
continually with respect to~$TB$, in the same
manner as does $BM$ with respect to~$OB$.

The arc~$BA$ always lies between $BA$ and $BN + NA$,
or $BM + MA$; hence $\dfrac{\arc BA}{\chord BA}$ lies between $1$ and
$\dfrac{BM}{BA} + \dfrac{MA}{BA}$. But $\dfrac{BM}{BA}$~has been shown to approach
continually towards~$1$, and $\dfrac{MA}{BA}$~to decrease without
limit; hence $\dfrac{\arc BA}{\chord BA}$ continually approaches towards~$1$.
\PageSep{11}
If $\angle BOA = 1°$,  $\dfrac{\arc BA}{\chord BA} = .0174533 ÷ .0174530 = 1.00002$,
very nearly. If $\angle BOA = 4'$, it is less than
$1.0000001$.

We now proceed to illustrate the various phrases
which have been used in enunciating these and similar
propositions.


\Subsection{The Notion of Infinitely Small Quantities.}

It appears that it is possible for two quantities $m$
and $m +  n$ to decrease together in such a way, that $n$~continually
decreases with respect to~$m$, that is, becomes
a less and less part of~$m$, so that $\dfrac{n}{m}$~also decreases
when $n$~and~$m$ decrease. Leibnitz,\footnote
  {Leibnitz was a native of Leipsic, and died in 1716, aged~70. His dispute
\index{Leibnitz}%
  with Newton, or rather with the English mathematicians in general, about
\index{Newton}%
  the invention of Fluxions, and the virulence with which it was carried on,
\index{Fluxions}%
  are well known. The decision of modern times appears to be that both Newton
  and Leibnitz were independent inventors of this method. It has, perhaps,
  not been sufficiently remarked how nearly several of their predecessors approached
  the same ground; and it is a question worthy of discussion, whether
  either Newton or Leibnitz might not have found broader hints in writings
  accessible to both, than the latter was ever asserted to have received from
  the former.}
in introducing
the Differential Calculus, presumed that in
such a case, $n$~might be taken so small as to be utterly
inconsiderable when compared with~$m$, so that $m + n$
might be put for~$m$, or \textit{vice versa}, without any error at
all. In this case he used the phrase that $n$~is \emph{infinitely}
small with respect to~$m$.

The following example will illustrate this term.
Since $(a  +  h)^{2} = a^{2} + 2ah + h^{2}$, it appears that if $a$~be
increased by~$h$, $a^{2}$~is increased by~$2ah + h^{2}$. But if $h$~be
taken very small, $h^{2}$~is very small with respect to~$h$,
for since $1:h :: h:h^{2}$, as many times as $1$~contains~$h$,
so many times does $h$~contain~$h^{2}$; so that by taking
\PageSep{12}
$h$~sufficiently small, $h$~may be made to be as many
times~$h^{2}$ as we please. Hence, in the words of Leibnitz,
if $h$~be taken \emph{infinitely} small, $h^{2}$~is \emph{infinitely} small
\index{Infinitely small, the notion of}%
with respect to~$h$, and therefore $2ah + h^{2}$ is the same
as~$2ah$; or if $a$~be increased by an infinitely small
quantity~$h$, $a^{2}$~is increased by another infinitely small
quantity~$2ah$, which is to~$h$ in the proportion of $2a$
to~$1$.

In this reasoning there is evidently an absolute
error; for it is impossible that $h$~can be so small, that
$2ah + h^{2}$ and $2ah$ shall be the same. The word \emph{small}
itself has no precise meaning; though the word \emph{smaller},
\index{Small, has no precise meaning}%
or \emph{less}, as applied in comparing one of two magnitudes
with another, is perfectly intelligible. Nothing is
either small or great in itself, these terms only implying
a relation to some other magnitude of the same
kind, and even then varying their meaning with the
subject in talking of which the magnitude occurs, so
that both terms may be applied to the same magnitude:
thus a large field is a very small part of the
earth. Even in such cases there is no natural point
at which smallness or greatness commences. The
thousandth part of an inch may be called a small distance,
a mile moderate, and a thousand leagues great,
but no one can fix, even for himself, the precise mean
between any of these two, at which the one quality
ceases and the other begins. These terms are not
therefore a fit subject for mathematical discussion,
until some more precise sense can be given to them,
which shall prevent the danger of carrying away with
the words, some of the confusion attending their use
in ordinary language. It has been usual to say that
when $h$~decreases from any given value towards nothing,
$h^{2}$~will become \emph{small} as compared with~$h$, because,
\PageSep{13}
let a number be ever so great, $h$~will, before it becomes
nothing, contain $h^{2}$~more than that number of
times. Here all dispute about a standard of smallness
is avoided, because, be the standard whatever it may,
the proportion of~$h^{2}$ to~$h$ may be brought under it. It
is indifferent whether the thousandth, ten-thousandth,
or hundred-millionth part of a quantity is to be considered
small enough to be rejected by the side of the
whole, for let $h$~be $\dfrac{1}{1000}$, $\dfrac{1}{10,000}$, or $\dfrac{1}{100,000,000}$ of the
unit, and $h$~will contain~$h^{2}$, $1000$, $10,000$, or $100,000,000$
of times.

The proposition, therefore, that $h$~can be taken so
small that $2ah + h^{2}$ and~$2ah$ are rigorously equal,
though not true, and therefore entailing error upon
all its subsequent consequences, yet is of this character,
that, by taking $h$ sufficiently small, all errors may
be made as small as we please. The desire of combining
simplicity with the appearance of rigorous
demonstration, probably introduced the notion of infinitely
small quantities; which was further established
by observing that their careful use never led to
any error. The method of stating the above-mentioned
proposition in strict and rational terms is as follows:
If $a$~be increased by~$h$, $a^{2}$~is increased by $2ah + h^{2}$,
which, whatever may be the value of~$h$, is to~$h$ in the
proportion of $2a + h$ to~$1$. The smaller $h$~is made,
the more near does this proportion diminish towards
that of $2a$ to~$1$, to which it may be made to approach
within any quantity, if it be allowable to take $h$ as
small as we please. Hence the ratio, $\emph{increment of } a^{2} ÷ \emph{increment of } a$, is a decreasing ratio, whose limit is~$2a$.

In further illustration of the language of Leibnitz,
\index{Leibnitz}%
we observe, that according to his phraseology, if $AB$~be
\PageSep{14}
an \emph{infinitely} small arc, the chord and arc~$AB$ are
equal, or the circle is a polygon of an \emph{infinite} number
of \emph{infinitely} small rectilinear sides. This should
be considered as an abbreviation of the proposition
proved (\PageRef{10}), and of the following: If a polygon
be inscribed in a circle, the greater the number of its
sides, and the smaller their lengths, the more nearly
will the perimeters of the polygon and circle be equal
to one another; and further, if any straight line be
given, however small, the difference between the perimeters
of the polygon and circle may be made less
than that line, by sufficient increase of the number of
sides and diminution of their lengths. Again, it would
be said (\Fig{1}) that if $AB$~be infinitely small, $MA$~is
infinitely less than~$BM$. What we have proved is,
that $MA$ may be made as small a part of~$BM$ as we
please, by sufficiently diminishing the arc~$BA$.


\Subsection{On Functions.}

An algebraical expression which contains~$x$ in any
\index{Functions!definition of|EtSeq}%
way, is called a \emph{function} of~$x$. Such are $x^{2} + a^{2}$,
$\dfrac{a + x}{a - x}$, $\log(x + y)$, $\sin 2x$. An expression may be a
function of more quantities than one, but it is usual
only to name those quantities of which it is necessary
to consider a change in the value. Thus if in $x^{2} + a^{2}$
$x$~only is considered as changing its value, this is
called a function of~$x$; if $x$~and~$a$ both change, it is
called a function of $x$~and~$a$. Quantities which change
their values during a process, are called \emph{variables}, and
\index{Variables!independent and dependent}%
those which remain the same, \emph{constants}; and variables
\index{Constants}%
which we change at pleasure are called \emph{independent},
while those whose changes necessarily follow from
\PageSep{15}
\index{Variables!independent and dependent}%
the changes of others are called \emph{dependent}. Thus in
\Fig{1}, the length of the radius~$OB$ is a constant, the
arc~$AB$ is the independent variable, while $BM$,~$MA$,
the chord~$AB$,~etc., are dependent. And, as in algebra
we reason on numbers by means of general symbols,
each of which may afterwards be particularised
as standing for any number we please, unless specially
prevented by the conditions of the problem, so, in
treating of functions, we use general symbols, which
may, under the restrictions of the problem, stand for
any function whatever. The symbols used are the letters
$F$,~$f$, $\Phi$,~$\phi$,~$\psi$; $\phi(x)$~and~$\psi(x)$, or $\phi x$~and~$\psi x$, may
represent any functions of~$x$, just as $x$~may represent
any number. Here it must be borne in mind that $\phi$~and~$\psi$
do not represent numbers which multiply~$x$, but
are \emph{the abbreviated directions to perform certain operations
with $x$ and constant quantities}. Thus, if $\phi x = x + x^{2}$,
$\phi$~is equivalent to a direction to add~$x$ to its
square, and the whole~$\phi x$ stands for the result of this
operation. Thus, in this case, $\phi(1) = 2$; $\phi(2) = 6$;
$\phi a = a +  a^{2}$; $\phi(x +  h) = x + h + (x +  h)^{2}$; $\phi \sin x = \sin x + (\sin x)^{2}$.
It may be easily conceived that this
notion is useless, unless there are propositions which
are generally true of all functions, and which may be
made the foundation of general reasoning.


\Subsection{Infinite Series.}

To exercise the student in this notation, we proceed
\index{Series|EtSeq}%
\index{Taylor's Theorem|EtSeq}%
to explain one of these functions which is of
most extensive application and is known by the name
of \emph{Taylor's Theorem}. If in~$\phi x$, any function of~$x$, the
value of~$x$ be increased by~$h$, or $x +  h$~be substituted
instead of~$x$, the result is denoted by~$\phi(x + h)$. It
\PageSep{16}
will generally\footnote
  {This word is used in making assertions which are for the most part
\index{Generally@\emph{Generally}, the word}%
  true, but admit of exceptions, few in number when compared with the other
  cases. Thus it generally happens that $x^{2} - 10x + 40$ is greater than~$15$, with
  the exception only of the case where $x = 5$. It is generally true that a line
  which meets a circle in a given point meets it again, with the exception only
  of the tangent.}
happen that this is either greater or
less than~$\phi x$, and $h$~is called the \emph{increment} of~$x$, and
\index{Increment}%
$\phi(x +  h) - \phi x$ is called the \emph{increment} of~$\phi x$, which is
negative when $\phi(x + h) < \phi x$. It may be proved
that $\phi(x + h)$ can generally be expanded in a series
of the form
\[
\phi x +  ph + qh^{2} + rh^{3} + \etc.,\quad \textit{ad infinitum},
\]
which contains none but whole and positive powers
of~$h$. It will happen, however, in many functions,
that one or more values can be given to~$x$ for which
it is impossible to expand $f(x + h)$ without introducing
negative or fractional powers. These cases are
considered by themselves, and the values of~$x$ which
produce them are called \emph{singular} values.
\index{Singular values}%

As the notion of a series which has no end of its
terms, may be new to the student, we will now proceed
to show that there may be series so constructed,
that the addition of any number of their terms, however
great, will always give a result less than some
determinate quantity. Take the series
\[
1 + x + x^{2} + x^{3} + x^{4} + \etc.,
\]
in which $x$~is supposed to be less than unity. The
first two terms of this series may be obtained by dividing
$1 - x^{2}$ by $1 - x$; the first three by dividing
$1 - x^{3}$ by $1 - x$; and the first $n$~terms by dividing
$1 - x^{n}$ by $1 - x$. If $x$~be less than unity, its successive
powers decrease without limit;\footnote
  {This may be proved by means of the proposition established in \Title{Study
  of Mathematics} (Chicago: The Open Court Publishing~Co., Reprint Edition),
  page~247. For $\dfrac{m}{n} × \dfrac{n}{m}$ is formed (if $m$~be less than~$n$) by dividing $\dfrac{m}{n}$ into $n$~parts,
  and taking away $n - m$ of them.}
that is, there is
\PageSep{17}
no quantity so small, that a power of~$x$ cannot be
found which shall be smaller. Hence by taking $n$~sufficiently
great, $\dfrac{1 - x^{n}}{1 - x}$ or $\dfrac{1}{1 - x} - \dfrac{x^{n}}{1 - x}$ may be
brought as near to~$\dfrac{1}{1 - x}$ as we please, than which,
however, it must always be less, since $\dfrac{x^{n}}{1 - x}$ can never
entirely vanish, whatever value $n$~may have, and therefore
there is always something subtracted from $\dfrac{1}{1 - x}$.
It follows, nevertheless, that $1 + x + x^{2} + \etc.$, if we
are at liberty to take as many terms as we please, can
be brought as near as we please to~$\dfrac{1}{1 - x}$, and in this
sense we say that
\[
\frac{1}{1 - x} = 1 + x + x^{2} + x^{3} + \etc.,\quad\textit{ad infinitum}.
\]


\Subsection{Convergent and Divergent Series.}

A series is said to be \emph{convergent} when the sum of
its terms tends towards some limit; that is, when, by
taking any number of terms, however great, we shall
never exceed some certain quantity. On the other
hand, a series is said to be \emph{divergent} when the sum of
a number of terms may be made to surpass any quantity,
however great. Thus of the two series,
\[
1 + \frac{1}{2} + \frac{1}{4} + \frac{1}{8} + \etc.
\]
and
\[
1 + 2 + 4 + 8 + \etc.\Add{,}
\]
the first is convergent, by what has been shown, and
the second is evidently divergent. A series cannot be
convergent, unless its separate terms decrease, so as,
\PageSep{18}
at last, to become less than any given quantity. And
the terms of a series may at first increase and afterwards
decrease, being apparently divergent for a finite
number of terms, and convergent afterwards. It will
only be necessary to consider the latter part of the
series.

Let the following series consist of terms decreasing
without limit:
\[
a + b + c + d + \dots + k + l + m + \dots,
\]
which may be put under the form
\[
%[** TN: Small parentheses in the original here and below, as noted]
a\left(1 + \frac{b}{a}
    + \frac{c}{b}\, \frac{b}{a}
    + \frac{d}{c}\, \frac{c}{b}\, \frac{b}{a} + \etc.\right);
\]
the same change of form may be made, beginning
from any term of the series, thus:
\[
%[** TN: Small ()]
k + l + m + \etc.
  = k\left(1 + \frac{l}{k} + \frac{m}{l}\, \frac{l}{k} + \etc.\right).
\]
We have introduced the new terms, $\dfrac{b}{a}$,~$\dfrac{c}{b}$,~etc., or the
ratios which the several terms of the original series
bear to those immediately preceding. It may be shown
(1)~that if the terms of the series $\dfrac{b}{a}$,~$\dfrac{c}{b}$,~$\dfrac{d}{c}$, etc., come
at last to be less than unity, and afterwards either
continue to approximate to a limit which is less than
unity, or decrease without limit, the series $a + b + c + \etc.$,
is convergent; (2)~if the limit of the terms
$\dfrac{b}{a}$,~$\dfrac{c}{b}$,~etc., is either greater than unity, or if they increase
without limit, the series is divergent.

(1\ia). Let $\dfrac{l}{k}$~be the first which is less than unity,
and let the succeeding ratios $\dfrac{m}{l}$,~etc., decrease, either
with or without limit, so that $\dfrac{l}{k} > \dfrac{m}{l} > \dfrac{n}{m}$, etc.;
whence it follows, that of the two series,
\PageSep{19}
\begin{align*}
%[** TN: Small ()]
k\biggl(1 &+ \frac{l}{k}
  + \frac{l}{k}\, \PadTo{\dfrac{m}{l}}{\dfrac{l}{k}}
  + \frac{l}{k}\,
      \PadTo{\dfrac{m}{l}}{\dfrac{l}{k}}\,
      \PadTo{\dfrac{n}{m}}{\dfrac{l}{k}}
  + \etc.\biggr), \\
k\biggl(1 &+ \frac{l}{k}
  + \frac{l}{k}\, \frac{m}{l}
  + \frac{l}{k}\, \frac{m}{l}\, \frac{n}{m}
  + \etc.\biggr),
\end{align*}
the first is greater than the second. But since $\dfrac{l}{k}$~is
less than unity, the first can never surpass $k × \dfrac{1}{1 - \dfrac{l}{k}}$,
or~$\dfrac{k^{2}}{k - l}$, and is convergent; the second is therefore
convergent. But the second is no other than $k + l + m + \etc.$;
therefore the series $a +  b + c + \etc.$, is convergent
from the term~$k$.

\Chg{(1\ib.)}{(1\ib).} Let $\dfrac{l}{k}$~be less than unity, and let the successive
ratios $\dfrac{l}{k}$,~$\dfrac{m}{l}$,~etc., increase, never surpassing a
limit~$A$, which is less than unity. Hence of the two
series,
\begin{align*}
%[** TN: Small ()]
k(1 &+ \PadTo{\dfrac{l}{k}}{A}
         + \PadTo{\dfrac{l}{k}}{A}\,\PadTo{\dfrac{m}{l}}{A}
         + \PadTo{\dfrac{l}{k}}{A}\,\PadTo{\dfrac{m}{l}}{A}\,\PadTo{\dfrac{n}{m}}{A}
         + \etc.), \\
k\biggl(1 &+ \frac{l}{k}
   + \frac{l}{k}\, \frac{m}{l}
   + \frac{l}{k}\, \frac{m}{l}\, \frac{n}{m}
   + \etc.\biggr),
\end{align*}
the first is the greater. But since $A$~is less than unity,
the first is convergent; whence, as before, $a + b + c + \etc.$,
converges from the term~$k$.

(2) The second theorem on the divergence of series
we leave to the student's consideration, as it is not
immediately connected with our object.


\Subsection{Taylor's Theorem. Derived Functions.}

We now proceed to the series
\index{Derivatives}%
\index{Derived Functions|EtSeq}%
\index{Functions!derived|EtSeq}%
\index{Taylor's Theorem|EtSeq}%
\[
ph + qh^{2} + rh^{3} + sh^{4} + \etc.,
\]
in which we are at liberty to suppose $h$ as small as
we please. The successive ratios of the terms to those
\PageSep{20}
immediately preceding are $\dfrac{qh^{2}}{ph}$ or~$\dfrac{q}{p}h$, $\dfrac{rh^{3}}{qh^{2}}$ or~$\dfrac{r}{q}h$,
$\dfrac{sh^{4}}{rh^{3}}$ or $\dfrac{s}{r}h$,~etc. If, then, the terms $\dfrac{q}{p}$,~$\dfrac{r}{q}$,~$\dfrac{s}{r}$, etc.,
are always less than a finite limit~$A$, or become so after
a definite number of terms, $\dfrac{q}{p}h$,~$\dfrac{r}{q}h$,~etc., will always
be, or will at length become, less than~$Ah$. And since $h$~may
be what we please, it may be so chosen that $Ah$~shall
be less than unity, for which $h$~must be less than~$\dfrac{1}{A}$.
In this case, by theorem~(1\ib), the series is convergent;
it follows, therefore, that a value of~$h$ can
always be found so small that $ph + qh^{2} + rh^{3} + \etc.$,
shall be convergent, at least unless the coefficients
$p$,~$q$,~$r$,~etc., be such that the ratio of any one to the
preceding increases without limit, as we take more
distant terms of the series. This never happens in
the developments which we shall be required to consider
in the Differential Calculus.

We now return to $\phi(x +h)$, which we have asserted
(\PageRef{16}) can be expanded (with the exception
of some particular values of~$x$) in a series of the form
$\phi x + ph + qh^{2} + \etc$. The following are some instances
of this development derived from the Differential
Calculus, most of which are also to be found in
treatises on algebra:
\index{Logarithms}%

{\scriptsize
\begin{alignat*}{4}
(x + h)^{n} &= x^{n}
    &+ nx^{n-1}h
   &&+ n(n - 1)x^{n-2} \frac{h^{2}}{2}
   &&+ n(n - 1)(n - 2)x^{n-3} \frac{h^{3}}{2·3} &\ \etc.\Add{,} \\
%
a^{x + h} &= a^{x}
    &+ ka^{x} h\rlap{\normalsize\footnotemark[1]}
   &&+ k^{2} a^{x} \frac{h^{2}}{2}
   &&+ k^{3} a^{x} \frac{h^{3}}{2·3} &\ \etc.\Add{,} \\
%
\log(x + h) &= \log x
    &+ \frac{1}{x}\, h
   &&- \frac{1}{x^{2}}\, \frac{h^{2}}{2}
   &&+ \frac{2}{x^{3}}\, \frac{h^{3}}{2·3} &\ \etc.\Add{,} \\
%
\sin(x + h) &= \sin x
    &+ \cos x\, h
   &&- \sin x\, \frac{h^{2}}{2}\rlap{\normalsize\footnotemark[2]}
   &&- \cos x\, \frac{h^{3}}{2·3} &\ \etc.\Add{,} \\
%
%[** TN: Moved up from top of page21]
\cos(x + h) &= \cos x
    &- \sin x\, h
   &&- \cos x\, \frac{h^{2}}{2}
   &&+ \sin x\, \frac{h^{3}}{2·3} &\ \etc.
\end{alignat*}}%
\footnotetext[1]{Here $k$~is the Naperian or hyperbolic logarithm of~$a$; that is, the common
  logarithm of~$a$ divided by~$.434294482$.}%
\footnotetext[2]{In the last two series the terms are positive and negative in pairs.}
\PageSep{21}

It appears, then, that the development of~$\phi(x + h)$
consists of certain functions of~$x$, the first of which is
$\phi x$~itself, and the remainder of which are multiplied
by $h$,~$\dfrac{h^{2}}{2}$, $\dfrac{h^{3}}{2·3}$, $\dfrac{h^{4}}{2·3·4}$, and so on. It is usual to denote
the coefficients of these divided powers of~$h$ by $\phi' x$,
$\phi'' x$, $\phi''' x$,\footnote
  {Called \emph{derived functions} or \emph{derivatives}.---\Ed.}
\index{Derivatives}%
\index{Derived Functions}%
\index{Functions!derived}%
etc., where $\phi'$,~$\phi''$,~etc., are merely functional
symbols, as is $\phi$~itself; but it must be recollected
that $\phi' x$,~$\phi'' x$,~etc., are rarely, if ever, employed
to signify anything except the coefficients of~$h$, $\dfrac{h^{2}}{2}$,~etc.,
in the development of~$\phi(x + h)$. Hence this development
is usually expressed as follows:
\[
\phi(x + h)
  = \phi x + \phi' x\, h + \phi''x\, \frac{h^{2}}{2} + \phi''' x\, \frac{h^{3}}{2·3} + \etc.
\]

{\Loosen Thus, when $\phi x = x^{n}$, $\phi' x = nx^{n-1}$, $\phi'' x = n(n - 1)x^{n-2}$, etc.;
when $\phi x = \sin x$, $\phi' x = \cos x$, $\phi'' x = -\sin x$,~etc.
In the first case $\phi'(x + h) = n(x + h)^{n-1}$,
$\phi''(x + h) = n(n - 1)(x + h)^{n-2}$; and in the second
$\phi'(x + h) = \cos (x + h)$, $\phi''(x + h) = -\sin(x + h)$.}

The following relation exists between $\phi x$,~$\phi' x$,
$\phi'' x$,~etc. In the same manner as $\phi' x$~is the coefficient
of~$h$ in the development of~$\phi(x + h)$, so $\phi'' x$~is the coefficient
of~$h$ in the development of~$\phi'(x + h)$, and
$\phi''' x$~is the coefficient of~$h$ in the development of~$\phi''(x + h)$;
$\phi^{\text{iv}} x$~is the coefficient of~$h$ in the development
of $\phi'''(x + h)$, and so on.

The proof of this is equivalent to \emph{Taylor's Theorem}
already alluded to (\PageRef{15}); and the fact may be
verified in the examples already given. When $\phi x = a^{x}$,
$\phi' x = ka^{x}$, and $\phi'(x + h) = ka^{x+h} = k(a^{x} + ka^{x}\, h + \etc.)$.
The coefficient of~$h$ is here~$k^{2} a^{x}$, which is the
\PageSep{22}
same as~$\phi'' x$. (See the second example of the preceding
table.) Again, $\phi''(x + h) = k^{2}a^{x+h} = k^{2}(a^{x} + ka^{x}\, h + \etc.)$,
in which the coefficient of~$h$ is~$k^{3}a^{x}$, the
same as~$\phi''' x$. Again, if $\phi x = \log x$, $\phi' x = \dfrac{1}{x}$, and
$\phi'(x + h) = \dfrac{1}{x + h} = \dfrac{1}{x} - \dfrac{h}{x^{2}} + \etc.$, as appears by
common division. Here the coefficient of~$h$ is~$-\dfrac{1}{x^{2}}$,
which is the same as $\phi'' x$~in the third example. Also
$\phi''(x + h) = -\dfrac{1}{(x + h)^{2}} = -(x + h)^{-2}$, which by the
Binomial Theorem is $-(x^{-2} - 2x^{-3}\, h + \etc.)$. The
coefficient of~$h$ is~$2x^{-3}$ or~$\dfrac{2}{x^{3}}$, which is~$\phi''' x$ in the
same example.


\Subsection{Differential Coefficients.}

It appears, then, that if we are able to obtain the
\index{Coefficients, differential|EtSeq}%
\index{Differential coefficients|EtSeq}%
coefficient of~$h$ in the development of \emph{any} function
whatever of~$x + h$, we can obtain all the other coefficients,
since we can thus deduce $\phi' x$ from~$\phi x$, $\phi'' x$
from~$\phi' x$, and so on. It is usual to call~$\phi' x$ the first
differential coefficient of~$\phi x$, $\phi'' x$~the second differential
coefficient of~$\phi x$, or the first differential coefficient
of~$\phi' x$; $\phi''' x$~the third differential coefficient of~$\phi x$,
or the second of~$\phi' x$, or the first of~$\phi'' x$; and so on.\footnote
  {The first, second, third, etc., differential coefficients, as thus obtained,
  are also called the first, second, third, etc., \emph{derivatives}.---\Ed.}
\index{Derivatives}%
The name is derived from a method of obtaining~$\phi' x$,
etc., which we now proceed to explain.

Let there be any function of~$x$, which we call~$\phi x$,
in which $x$~is increased by an increment~$h$; the function
then becomes
\[
\phi x + \phi' x\, h
       + \phi'' x\, \frac{h^{2}}{2}
       + \phi''' x\, \frac{h^{3}}{2·3} + \etc.
\]
\PageSep{23}
The original value~$\phi x$ is increased by the increment
\[
\phi' x\, h + \phi'' x\, \frac{h^{2}}{2} x
            + \phi''' x\, \frac{h^{3}}{2·3} + \etc.;
\]
whence ($h$~being the increment of~$x$)
\[
\frac{\emph{increment of } \phi x}{\emph{increment of } x}
  = \phi' x + \phi'' x\, \frac{h}{2} x
            + \phi''' x\, \frac{h^{2}}{2·3} + \etc.,
\]
which is an expression for the ratio which the increment
of a function bears to the increment of its variable.
It consists of two parts. The one,~$\phi' x$, into
which $h$~does not enter, depends on $x$~only; the remainder
is a series, every term of which is multiplied
by some power of~$h$, and which therefore diminishes
as $h$~diminishes, and may be made as small as we
please by making $h$~sufficiently small.

To make this last assertion clear, observe that all
the ratio, except its first term~$\phi' x$, may be written as
follows:
\[
%[** TN: Small () in the original]
h\left(\phi'' x\, \frac{1}{2} + \phi''' x\, \frac{h}{2·3} + \etc.\right);
\]
the second factor of which (\PageRef{19}) is a convergent
series whenever $h$~is taken less than~$\dfrac{1}{A}$, where $A$~is
the limit towards which we approximate by taking
the coefficients $\phi'' x × \dfrac{1}{2}$, $\phi''' x × \dfrac{1}{2·3}$,~etc., and forming
the ratio of each to the one immediately preceding.
This limit, as has been observed, is finite in
every series which we have occasion to use; and
therefore a value for~$h$ can be chosen so small, that
for it the series in the last-named formula is convergent;
still more will it be so for every smaller value
of~$h$. Let the series be called~$P$. If $P$~be a finite quantity,
which decreases when $h$~decreases, $Ph$~can be
made as small as we please by sufficiently diminishing~$h$;
\PageSep{24}
whence $\phi' x + Ph$ can be brought as near as we
please to~$\phi' x$. Hence the ratio of the increments of
$\phi x$ and~$x$, produced by changing $x$ into~$x + h$, though
never equal to~$\phi' x$, approaches towards it as $h$~is diminished,
and may be brought as near as we please
to it, by sufficiently diminishing~$h$. Therefore to find
the coefficient of~$h$ in the development of~$\phi(x + h)$,
find $\phi(x + h) - \phi x$, divide it by~$h$, and find the limit
towards which it tends as $h$~is diminished.

In any series such as
\index{Series|EtSeq}%
\[
a + bh + ch^{2} + \dots + kh^{n} + lh^{n+1} + mh^{n+2} + \etc.
\]
which is such that some given value of~$h$ will make it
convergent, it may be shown that $h$~can be taken so
small that any one term shall contain all the succeeding
ones as often as we please. Take any one term,
as~$kh^{n}$. It is evident that, be $h$ what it may,
\[
kh^{n} : lh^{n+1} + mh^{n+2} + \etc.,\ ::\ k : lh + mh^{2} + \etc.,
\]
the last term of which is $h(l + mh + \etc.)$. By reasoning
similar to that in the last paragraph, we can
show that this may be made as small as we please,
since one factor is a series which is always finite when
$h$~is less than~$\dfrac{1}{A}$, and the other factor~$h$ can be made
as small as we please. Hence, since $k$~is a given
quantity, independent of~$h$, and which therefore remains
the same during all the changes of~$h$, the series
$h(l + mh + \etc.)$ can be made as small a part of~$k$ as
we please, since the first diminishes without limit,
and the second remains the same. By the proportion
above established, it follows then that $lh^{n+1} + mh^{n+2} + \etc.$,
can be made as small a part as we please of~$kh^{n}$.
It follows, therefore, that if, instead of the full
development of~$\phi(x + h)$, we use only its two first
\PageSep{25}
terms $\phi x +  \phi' x\, h$, the error thereby introduced may,
by taking $h$ sufficiently small, be made as small a portion
as we please of the small term~$\phi' x\, h$.


\Subsection{The Notation of the Differential Calculus.}

The first step usually made in the Differential Calculus
\index{Calculus, notation of}%
\index{Notation!of the Differential Calculus}%
is the determination of~$\phi' x$ for all possible values
of~$\phi x$, and the construction of general rules for
that purpose. Without entering into these we proceed
to explain the notation which is used, and to apply
the principles already established to the solution
of some of those problems which are the peculiar
province of the Differential Calculus.

When any quantity is increased by an increment,
which, consistently with the conditions of the problem,
may be supposed as small as we please, this increment
is denoted, not by a separate letter, but by
prefixing the letter~$d$, either followed by a full stop or
not, to that already used to signify the quantity. For
example, the increment of~$x$ is denoted under these
circumstances by~$dx$; that of~$\phi x$ by~$d.\phi x$; that of~$x^{n}$
by~$d.x^{n}$. If instead of an increment a decrement
be used, the sign of~$dx$, etc., must be changed in all
expressions which have been obtained on the supposition
of an increment; and if an increment obtained
by calculation proves to be negative, it is a sign that
a quantity which we imagined was increased by our
previous changes, was in fact diminished. Thus, if
$x$~becomes $x + dx$, $x^{2}$~becomes $x^{2} + d.x^{2}$. But this is
also $(x + dx)^{2}$ or $x^{2} + 2x\, dx + (dx)^{2}$; whence $d.x^{2} = 2x\, dx + (dx)^{2}$.
Care must be taken not to confound
$d.x^{2}$, the increment of~$x^{2}$, with~$(dx)^{2}$, or, as it is often
written,~$dx^{2}$, the square of the increment of~$x$. Again,
\PageSep{26}
if $x$~becomes $x + dx$, $\dfrac{1}{x}$ becomes $\dfrac{1}{x} + d.\dfrac{1}{x}$ and the
change of~$\dfrac{1}{x}$ is $\dfrac{1}{x + dx} - \dfrac{1}{x}$ or $-\dfrac{dx}{x^{2} + x\, dx}$; showing
that an increment of~$x$ produces a decrement in~$\dfrac{1}{x}$.

It must not be imagined that because $x$~occurs in
the symbol~$dx$, the value of the latter in any way depends
upon that of the former: both the first value of~$x$,
and the quantity by which it is made to differ from
its first value, are at our pleasure, and the letter~$d$ must
merely be regarded as an abbreviation of the words
``\emph{difference of}.'' In the first example, if we divide
\index{Differences!of increments}%
both sides of the resulting equation by~$dx$, we have
$\dfrac{d.x^{2}}{dx} = 2x + dx$. The smaller $dx$~is supposed to be,
the more nearly will this equation assume the form
$\dfrac{d.x^{2}}{dx} = 2x$, and the ratio of $2x$ to~$1$ is the limit of the
\index{Limits|EtSeq}%
ratio of the increment of~$x^{2}$ to that of~$x$; to which
this ratio may be made to approximate as nearly as
we please, but which it can never actually reach. In
the Differential Calculus, the limit of the ratio only is
retained, to the exclusion of the rest, which may be
explained in either of the two following ways:

(1) The fraction $\dfrac{d.x^{2}}{dx}$ may be considered as standing,
not for any value which it can actually have as
long as $dx$~has a real value, but for the limit of all
those values which it assumes while $dx$~diminishes.
In this sense the equation $\dfrac{d.x^{2}}{dx} = 2x$ is strictly true.
But here it must be observed that the algebraical
meaning of the sign of division is altered, in such a
way that it is no longer allowable to use the numerator
and denominator separately, or even at all to consider
\PageSep{27}
them as quantities. If $\dfrac{dy}{dx}$~stands, not for the
ratio of two quantities, but for the limit of that ratio,
which cannot be obtained by taking any real value of~$dx$,
however small, the whole $\dfrac{dy}{dx}$ may, by convention,
have a meaning, but the separate parts $dy$ and~$dx$
have none, and can no more be considered as separate
quantities whose ratio is~$\dfrac{dy}{dx}$, than the two loops
of the figure~$8$ can be considered as separate numbers
whose sum is eight. This would be productive of no
great inconvenience if it were never required to separate
the two; but since all books on the Differential
Calculus and its applications are full of examples in
which deductions equivalent to assuming $dy = 2x\, dx$
are drawn from such an equation as $\dfrac{dy}{dx} = 2x$, it becomes
necessary that the first should be explained, independently
of the meaning first given to the second.
It may be said, indeed, that if $y = x^{2}$, it follows that
$\dfrac{dy}{dx} = 2x + dx$, in which, \emph{if we make $dx = 0$}, the result
is $\dfrac{dy}{dx} = 2x$. But if $dx = 0$, $dy$~also~$= 0$, and this
equation should be written $\dfrac{0}{0} = 2x$, as is actually done
in some treatises on the Differential Calculus,\footnote
  {This practice was far more common in the early part of the century
  than now, and was due to the precedent of Euler (1755). For the sense in
\index{Euler}%
  which Euler's view was correct, see the \Title{Encyclopedia Britannica}, art.\ \Title{Infinitesimal
  Calculus}, Vol.~XII, p.~14, 2nd~column.---\Ed.}
to the
great confusion of the learner. Passing over the difficulties\footnote
  {See \Title{Study of Mathematics} (Reprint Edition, Chicago: The Open Court
  Publishing~Co., 1898), page~126.}
of the fraction~$\dfrac{0}{0}$, still the former objection
recurs, that the equation $dy = 2x\, dx$ cannot be used
\PageSep{28}
(and it \emph{is} used even by those who adopt this explanation)
without supposing that~$0$, which merely implies
an absence of all magnitude, can be used in different
senses, so that one~$0$ may be contained in another a
certain number of times. This, even if it can be considered
as intelligible, is a notion of much too refined
a nature for a beginner.

(2) The presence of the letter~$d$ is an indication,
not only of an increment, but of an increment which
we are at liberty to suppose as small as we please.
The processes of the Differential Calculus are intended
to deduce relations, not between the ratios of different
increments, but between the limits to which those ratios
approximate, when the increments are decreased.
And it may be true of some parts of an equation, that
though the taking of them away would alter the relation
between $dy$ and~$dx$, it would not alter the limit
towards which their ratio approximates, when $dx$
and~$dy$ are diminished. For example, $dy = 2x\, dx + (dx)^{2}$.
If $x = 4$ and $dx = .01$, then $dy = .0801$ and
$\dfrac{dy}{dx} = 8.01$. If $dx = .0001$, $dy = .00080001$ and $\dfrac{dy}{dx} = 8.0001$.
The limit of this ratio, to which we shall
come still nearer by making $dx$ still smaller, is~$8$. The
term~$(dx)^{2}$, though its presence affects the value of~$dy$
and the ratio~$\dfrac{dy}{dx}$, does not affect the limit of the latter,
for in $\dfrac{dy}{dx}$ or $2x + dx$, the latter term~$dx$, which arose
from the term~$(dx)^{2}$, diminishes continually and without
limit. If, then, we throw away the term~$(dx)^{2}$,
the consequence is that, make $dx$ what we may, we
never obtain~$dy$ as it would be if correctly deduced
from the equation $y = x^{2}$, but we obtain the limit of
the ratio of~$dy$ to~$dx$. If we throw away all powers of~$dx$
\PageSep{29}
above the first, and use the equations so obtained,
all ratios formed from these last, or their consequences,
are themselves the limiting ratios of which we are in
search. \emph{The equations which we thus use are not absolutely
true in any case, but may be brought as near as we
please to the truth}, by making $dy$~and~$dx$ sufficiently
small. If the student at first, instead of using $dy = 2x\, dx$,
were to write it thus, $dy = 2x\, dx + \etc.$, the \emph{etc.}\
would remind him that there are other terms; \emph{necessary},
if the value of~$dy$ corresponding to any value of~$dx$
is to be obtained; \emph{unnecessary}, if the \emph{limit} of the
ratio of $dy$ to~$dx$ is all that is required.

We must adopt the first of these explanations when
$dy$ and $dx$ appear in a fraction, and the second when
they are on opposite sides of an equation.


\Subsection{Algebraical Geometry.}

If two straight lines be drawn at right angles to
each other, dividing the whole of their plane into four
parts, one lying in each right angle, the situation
of any point is determined when we know, (1)~in
which angle it lies, and (2)~its perpendicular distances
from the two right lines. Thus (\Fig{2}) the point~$P$
lying in the angle~$AOB$, is known when $PM$~and~$PN$,
or when $OM$~and~$PM$ are known; for, though there
is an infinite number of points whose distance from~$OA$
only is the same as that of~$P$, and an infinite number
of others, whose distance from~$OB$ is the same as
that of~$P$, there is no other point whose distances
from both lines are the same as those of~$P$. The line~$OA$
is called the axis of~$x$, because it is usual to denote
any variable distance measured on or parallel to~$OA$
by the letter~$x$. For a similar reason, $OB$~is called
\PageSep{30}
\index{Co-ordinates}%
the axis of~$y$. The \emph{co-ordinates}\footnote
  {The distances $OM$ and~$MP$ are called the \emph{co-ordinates} of the point~$P$. It
  is moreover usual to call the co-ordinate~$OM$, the \emph{abscissa}, and $MP$, the \emph{ordinate},
  of the point~$P$.}
or perpendicular distances
of a point~$P$ which is supposed to vary its position,
are thus denoted by $x$~and~$y$; hence $OM$ or~$PN$
is~$x$, and $PM$ or~$ON$ is~$y$. Let a linear unit be chosen,
so that any number may be represented by a straight
line. Let the point~$M$, setting out from~$O$, move in
the direction~$OA$, always carrying with it the indefinitely
extended line~$MP$ perpendicular to~$OA$. While
this goes on, let $P$~move upon the line~$MP$ in such a
way, that $MP$ or~$y$ is always equal to a given function
of~$OM$ or~$x$; for example, let $y = x^{2}$, or let the number
\index{Parabola, the}%
\Figure{2}
of units in~$PM$ be the square of the number of
units in~$OM$. As $O$~moves towards~$A$, the point~$P$
will, by its motion on~$MP$, compounded with the motion
of the line $MP$ itself, describe a curve~$OP$, in
which $PM$~is less than, equal to, or greater than,~$OM$,
according as $OM$~is less than, equal to, or greater
than the linear unit. It only remains to show how
the other branch of this curve is deduced from the
equation $y = x^{2}$. And to this end we shall first have
to interpolate a few remarks.
\PageSep{31}


\Subsection{On the Connexion of the Signs of Algebraical and
the Directions of Geometrical Magnitudes.}

It is shown in algebra, that if, through misapprehension
\index{Signs|EtSeq}%
of a problem, we measure in one direction, a
line which ought to lie in the exactly opposite direction,
or if such a mistake be a consequence of some
previous misconstruction of the figure, any attempt
to deduce the length of that line by algebraical reasoning,
will give a negative quantity as the result.
And conversely it may be proved by any number of
examples, that when an equation in which $a$~occurs
has been deduced strictly on the supposition that $a$~is
a line measured in one direction, a change of sign in~$a$
will turn the equation into that which would have
been deduced by the same reasoning, had we begun
by measuring the line~$a$ in the contrary direction.
Hence the change of~$+a$ into~$-a$, or of~$-a$ into~$+a$,
corresponds in geometry to a change of direction of
the line represented by~$a$, and \textit{vice versa}.

In illustration of this general fact, the following
\index{Circle, equation of|EtSeq}%
\index{Circle cut by straight line, investigated|EtSeq}%
problem may be useful. Having a circle of given radius,
whose centre is in the intersection of the axes
of $x$~and~$y$, and also a straight line cutting the axes in
two given points, required the co-ordinates of the
points (if any) in which the straight line cuts the circle.
Let $OA$, the radius of the circle~$= r$, $OE = a$,
$OF = b$, and let the co-ordinates of~$P$, one of the
points of intersection required, be $OM = x$, $MP = y$.
(\Fig{3}.) The point~$P$ being in the circle whose radius
is~$r$, we have from the right-angled triangle~$OMP$,
$x^{2} + y^{2} = r^{2}$, which equation belongs to the co-ordinates
of every point in the circle, and is called
\PageSep{32}
the equation of the circle. Again, $EM : MP :: EO : OF$
by similar triangles; or $a - x : y :: a : b$, whence $ay + bx = ab$,
which is true, by similar reasoning, for every
point of the line~$EF$. But for a point~$P'$ lying in~$EF$
produced, we have $EM' : M'P' :: EO : OF$, or $x + a : y :: a : b$,
whence $ay - bx = ab$, an equation which may
be obtained from the former by changing the sign of~$x$;
and it is evident that the direction of~$x$, in the
\Figure{3}
second case, is opposite to that in the first. Again,
for a point~$P''$ in $FE$ produced, we have $EM'' : M''P'' :: EO : OF$,
or $x - a : y :: a : b$, whence $bx - ay = ab$, which
may be deduced from the first by changing the sign
of~$y$; and it is evident that $y$~is measured in different
directions in the first and third cases. Hence the
equation $ay + bx = ab$ belongs to all parts of the
straight line~$EF$, if we agree to consider $M''P''$ as
negative, when $MP$~is positive, and $OM'$~as negative
\PageSep{33}
when $OM$~is positive. Thus, if $OE = 4$, and $OF = 5$,
and $OM = 1$, we can determine~$MP$ from the equation
$ay + bx = ab$, or $4y + 5 = 20$, which gives $y$~or $MP = 3\frac{3}{4}$.
But if $OM'$~be $1$ in length, we can determine~$M'P'$
either by calling $OM'$,~$1$, and using the equation
$ay - bx = ab$, or calling $OM'$,~$-1$, and using the equation
$ay + bx = ab$, as before. Either gives $M'P' = 6\frac{1}{4}$.
The latter method is preferable, inasmuch as it enables
us to contain, in one investigation, all the different
cases of a problem.

We shall proceed to show that this may be done
in the present instance. We have to determine the
co-ordinates of the point~$P$, from the following equations:
\begin{align*}
ay + bx &= ab, \\
x^{2} + y^{2} &= r^{2}.
\end{align*}
Substituting in the second the value of~$y$ derived from
the first, or $b\left(\dfrac{a - x}{b}\right)$, we have
\[
x^{2} + b^{2}\, \frac{(a - x)^{2}}{a^{2}} = r^{2},
\]
or
\[
(a^{2} + b^{2}) x^{2} - 2ab^{2}x + a^{2}(b^{2} - r^{2}) = 0;
\]
and proceeding in a similar manner to find~$y$, we have
\[
(a^{2} + b^{2}) y^{2} - 2a^{2}by + b^{2}(a^{2} - r^{2}) = 0,
\]
which \Typo{give}{gives}
\begin{align*}
x &= a\, \frac{b^{2} ± \sqrt{(a^{2} + b^{2})r^{2} - a^{2}b^{2}}}{a^{2} + b^{2}}, \\
y &= b\, \frac{a^{2} \mp \sqrt{(a^{2} + b^{2})r^{2} - a^{2}b^{2}}}{a^{2} + b^{2}};
\end{align*}
the upper or the lower sign to be taken in both.
Hence when $(a^{2} + b^{2})r^{2} > a^{2}b^{2}$, that is, when $r$~is greater
than the perpendicular let fall from~$O$ upon~$EF$, which
perpendicular is
\PageSep{34}
\[
\frac{ab}{\sqrt{a^{2} + b^{2}}},
\]
there are two points of intersection. When $(a^{2} + b^{2})r^{2} = a^{2}b^{2}$,
the two values of~$x$ become equal, and also
those of~$y$, and there is only one point in which the
straight line meets the circle; in this case $EF$~is a
tangent to the circle. And if $(a^{2} + b^{2})r^{2} < a^{2}b^{2}$, the
values of $x$~and~$y$ are impossible, and the straight line
does not meet the circle.

Of these three cases, we confine ourselves to the
first, in which there are two points of intersection.
The product of the values of~$x$, with their proper
sign, is\footnote
  {See \Title{Study of Mathematics} (Chicago: The Open Court Pub.~Co.), page~136.}
\[
a^{2}\, \frac{b^{2} - r^{2}}{a^{2} + b^{2}},
\]
and of~$y$,
\[
b^{2}\, \frac{a^{2} - r^{2}}{a^{2} + b^{2}},
\]
the signs of which are the same as those of~$b^{2} - r^{2}$,
and $a^{2} - r^{2}$. If $b$~and~$a$ be both $> r$, the two values
of~$x$ have the same sign; and it will appear from the
figure, that the lines they represent are measured in
the same direction. And this whether $b$~and~$a$ be positive
or negative, since $b^{2} - r^{2}$ and $a^{2} - r^{2}$ are both
positive when $a$~and~$b$ are numerically greater than~$r$,
whatever their signs may be. That is, if our rule,
connecting the signs of algebraical and the directions
of geometrical magnitudes, be true, let the directions
of $OE$ and $OF$ be altered in any way, so long as $OE$
and $OF$ are both greater than~$OA$, the two values of~$OM$
will have the same direction, and also those of~$MP$.
This result may easily be verified from the
figure.
\PageSep{35}

Again, the values of $x$~and~$y$ having the same sign,
that sign will be (see the equations) the same as that
of $2ab^{2}$ for~$x$, and of $2a^{2}b$ for~$y$, or the same as that of
$a$~for~$x$ and of $b$~for~$y$. That is, when $OE$~and~$OF$ are
both greater than~$OA$, the direction of each set of co-ordinates
will be the same as those of $OE$ and~$OF$,
which may also be readily verified from the figure.

Many other verifications might thus be obtained of
the same principle, viz., that any equation which corresponds
to, and is true for, all points in the angle~$AOB$,
may be used without error for all points lying
in the other three angles, by substituting the proper
numerical values, with a negative sign, for those co-ordinates
whose directions are opposite to those of
the co-ordinates in the angle~$AOB$. In this manner,
if four points be taken similarly situated in the four
angles, the numerical values of whose co-ordinates
are $x = 4$ and  $y = 6$, and if the co-ordinates of that
point which lies in the angle~$AOB$, are called $+4$ and~$+6$;
those of the points lying in the angle~$BOC$ will
be $-4$~and~$+6$; in the angle~$COD$ $-4$~and~$-6$;
and in the angle~$DOE$ $+4$~and~$-6$.

To return to \Fig{2}, if, after having completed the
branch of the curve which lies on the right of~$BC$,
and whose equation is $y = x^{2}$, we seek that which lies
on the left of~$BC$, we must, by the principles established,
substitute $-x$ instead of~$x$, when the numerical
value obtained for~$(-x)^{2}$ will be that of~$y$, and the
sign will show whether $y$~is to be measured in a similar
or contrary direction to that of~$MP$. Since $(-x)^{2} = x^{2}$,
the direction and value of~$y$, for a given value
of~$x$, remains the same as on the right of~$BC$; whence
the remaining branch of the curve is similar and equal
in all respects to~$OP$, only lying in the angle~$BOD$.
\PageSep{36}
And thus, if $y$ be any function of~$x$, we can obtain a
geometrical representation of the same, by making $y$
the ordinate, and $x$~the abscissa of a curve, every ordinate
of which shall be the linear representation of
the numerical value of the given function corresponding
to the numerical value of the abscissa, the linear
unit being a given line.


\Subsection{The Drawing of a Tangent to a Curve.}

If the point~$P$ (\Fig{2}), setting out from~$O$, move
along the branch~$OP$, it will continually change the
\Figure[nolabel]{2}
\emph{direction} of its motion, never moving, at one point, in
\index{Direction}%
the direction which it had at any previous point. Let
the moving point have reached~$P$, and let $OM = x$,
$MP = y$. Let $x$~receive the increment $MM' = dx$, in
consequence of which $y$ or $MP$ becomes~$M'P'$, and
receives the increment $QP' = dy$; so that $x + dx$ and
$y + dy$ are the co-ordinates of the moving point~$P$,
when it arrives at~$P'$. Join~$PP'$, which makes, with
$PQ$ or~$OM$, an angle, whose tangent is $\dfrac{P'Q}{PQ}$ or~$\dfrac{dy}{dx}$.
Since the relation $y = x^{2}$ is true for the co-ordinates of
every point in the curve, we have $y + dy = (x +  dx)^{2}$,
\PageSep{37}
the subtraction of the former equation from which
gives $dy = 2x\, dx + (dx)^{2}$, or $\dfrac{dy}{dx} = 2x + dx$. If the
point~$P'$ be now supposed to move backwards towards~$P$,
the chord~$PP'$ will diminish without limit, and the
inclination of $PP'$ to $PQ$ will also diminish, but not
without limit, since the tangent of the angle~$P'PQ$, or~$\dfrac{dy}{dx}$,
\index{Tangent}%
is always greater than the limit~$2x$. If, therefore,
a line~$PV$ be drawn through~$P$, making with~$PQ$ an
angle whose tangent is~$2x$, the chord~$PP'$ will, as $P'$~approaches
towards~$P$, or as $dx$~is diminished, continually
approximate towards~$PV$, so that the angle~$P'PV$
may be made smaller than any given angle, by
sufficiently diminishing~$dx$. And the line~$PV$ cannot
again meet the curve on the side of~$PP'$, nor can any
straight line be drawn between it and the curve, the
proof of which we leave to the student.

Again, if $P'$~be placed on the other side of~$P$, so that
its co-ordinates are $x - dx$ and $y - dy$, we have $y - dy = (x - dx)^{2}$,
which, subtracted from $y = x^{2}$, gives $dy = 2x\, dx - (dx)^{2}$,
or $\dfrac{dy}{dx} = 2x - dx$. By similar reasoning,
if the straight line~$PT$ be drawn in continuation
of~$PV$, making with~$PN$ an angle, whose tangent is~$2x$,
the chord~$PP'$ will continually approach to this
line, as before.

The line~$TPV$ indicates the direction in which the
point~$P$ is proceeding, and is called the \emph{tangent} of the
curve at the point~$P$. If the curve were the interior
of a small solid tube, in which an atom of matter were
made to move, being projected into it at~$O$, and if all
the tube above~$P$ were removed, the line~$PV$ is in the
direction which the atom would take on emerging at~$P$,
and is the line which it would describe. The angle
\PageSep{38}
which the tangent makes with the axis of~$x$ in any
\index{Tangent}%
curve, may be found by giving $x$ an increment, finding
the ratio which the corresponding increment of~$y$
bears to that of~$x$, and determining the limit of that
ratio, or the \emph{differential coefficient}. This limit is the
\index{Coefficients, differential}%
\index{Differential coefficients}%
trigonometrical tangent\footnote
  {There is some confusion between these different uses of the word tangent.
  The geometrical tangent is, as already defined, the line between which
  and a curve no straight line can be drawn; the trigonometrical tangent has
  reference to an angle, and is the ratio which, in any right-angled triangle,
  the side opposite the angle bears to that which is adjacent.}
of the angle which the geometrical
tangent makes with the axis of~$x$. If $y = \phi x$,
$\phi' x$~is this trigonometrical tangent. Thus, if the curve
be such that the ordinates are the Naperian logarithms\footnote
  {It may be well to notice that in analysis the Naperian logarithms are
\index{Logarithms}%
  the only ones used; while in practice the common, or Briggs's logarithms,
  are always preferred.}
of the abscissæ, or $y = \log x$, and $y + dy =
\log x + \dfrac{1}{x}\, dx - \dfrac{1}{2x^{2}}\, dx^{2}$, etc., the geometrical tangent
of any point whose abscissa is~$x$, makes with the axis
an angle whose trigonometrical tangent is~$\dfrac{1}{x}$.

This problem, of drawing a tangent to any curve,
was one, the consideration of which gave rise to the
methods of the Differential Calculus.


\Subsection{Rational Explanation of the Language of Leibnitz.}

As the peculiar language of the theory of infinitely
\index{Infinitely small, the notion of|EtSeq}%
\index{Leibnitz}%
small quantities is extensively used, especially in
works of natural philosophy, it has appeared right to
us to introduce it, in order to show how the terms
which are used may be made to refer to some natural
and rational mode of explanation. In applying this
language to \Fig{2}, it would be said that the curve~$OP$
is a polygon consisting of an infinite number of
\index{Polygon}%
\PageSep{39}
infinitely small sides, each of which produced is a
tangent to the curve; also that if $MM'$ be taken infinitely
small, the chord and arc~$PP'$ coincide with
\index{Arc and its chord, a continuously decreasing|EtSeq}%
one of these rectilinear elements; and that an infinitely
small arc coincides with its chord. All which
must be interpreted to mean that, the chord and arc
being diminished, approach more and more nearly to
a ratio of equality as to their lengths; and also that
the greatest separation between an arc and its chord
may be made as small a part as we please of the whole
chord or arc, by sufficiently diminishing the chord.

We shall proceed to a strict proof of this; but in
the meanwhile, as a familiar illustration, imagine a
small arc to be cut off from a curve, and its extremities
joined by a chord, thus forming an arch, of which
the chord is the base. From the middle point of the
chord, erect a perpendicular to it, meeting the arc,
which will thus represent the height of the arch.
Imagine this figure to be magnified, without distortion
or alteration of its proportions, so that the larger figure
may be, as it is expressed, a true picture of the
smaller one. However the original arc may be diminished,
let the magnified base continue of a given
length. This is possible, since on any line a figure
may be constructed similar to a given figure. If the
original curve could be such that the height of the
arch could never be reduced below a certain part of
the chord, say one thousandth, the height of the magnified
arch could never be reduced below one thousandth
of the magnified chord, since the proportions
of the two figures are the same. But if, in the original
curve, an arc can be taken so small that the height
of the arch is as small a part as we please of the
chord, it will follow that in the magnified figure where
\PageSep{40}
the chord is always of one length, the height of the
arch can be made as small as we please, seeing that
it can be made as small a part as we please of a given
line. It is possible in this way to conceive a whole
curve so magnified, that a given arc, however small,
shall be represented by an arc of any given length,
however great; and the proposition amounts to this,
that let the dimensions of the magnified curve be any
\index{Curve, magnified}%
\index{Magnified curve}%
given number of times the original, however great, an
arch can be taken upon the original curve so small,
that the height of the corresponding arch in the magnified
figure shall be as small as we please.
\Figure{4}

Let $PP'$ (\Fig{4}) be a part of a curve, whose equation
is $y = \phi(x)$, that is, $PM$~may always be found by
substituting the numerical value of~$OM$ in a given
function of~$x$. Let $OM = x$ receive the increment
$MM' = dx$, which we may afterwards suppose as small
as we please, but which, in order to render the figure
more distinct, is here considerable. The value of $PM$
or~$y$ is~$\phi x$, and that of $P'M'$ or $y + dy$ is~$\phi(x +  dx)$.

Draw $PV$, the tangent at~$P$, which, as has been
\index{Tangent}%
shown, makes, with~$PQ$, an angle, whose trigonometrical
tangent is the limit of the ratio~$\dfrac{dy}{dx}$, when $x$~is decreased,
or~$\phi' x$. Draw the chord~$PP'$, and from any
\PageSep{41}
point in it, for example, its middle point~$p$, draw~$pv$
parallel to~$PM$, cutting the curve in~$a$. The value of~$P'Q$,
or~$dy$, or $\phi(x + dx) - \phi x$ is %[** TN: This line displayed in the orig]
\[
P'Q = \phi' x\, dx
    + \phi'' x\, \frac{(dx)^{2}}{2}
    + \phi''' x\, \frac{(dx)^{3}}{2·3} + \etc.
\]
But $\phi' x\, dx$~is $\tan VPQ·PQ = VQ$. Hence $VQ$~is the
first term of this series, and $P'V$~the aggregate of the
rest. But it has been shown that $dx$~can be taken so
small, that any one term of the above series shall contain
the rest, as often as we please. Hence $PQ$~can
be taken so small that $VQ$~shall contain~$VP'$ as often
as we please, or the ratio of $VQ$ to~$VP'$ shall be as
great as we please. And the ratio $VQ$ to~$PQ$ continues
finite, being always~$\phi' x$; hence $P'V$~also decreases
without limit as compared with~$PQ$.

Next, the chord~$PP'$ or $\sqrt{(dx)^{2} + (dy)^{2}}$, or
\[
dx \sqrt{1 + \left(\frac{dy}{dx}\right)^{2}}
\]
is to~$PQ$ or~$dx$ in the ratio of $\sqrt{1 + \left(\dfrac{dy}{dx}\right)^{2}} : 1$, which,
as $PQ$~is diminished, continually approximates to that
of $\sqrt{1 + (\phi' x)^{2}} : 1$, which is the ratio of~$PV : PQ$.
Hence the ratio of~$PP' : PV$ continually approaches to
unity, or $PQ$~may be taken so small that the difference
of $PP'$~and~$PV$ shall be as small a part of either
of them as we please.

Finally, the arc~$PP'$ is greater than the chord~$PP'$
and less than $PV + VP'$. Hence $\dfrac{\arc PP'}{\chord PP'}$ lies between
$1$~and $\dfrac{PV}{PP'} + \dfrac{VP'}{PP'}$, the former of which two
fractions can be brought as near as we please to unity,
and the latter can be made as small as we please; for
\PageSep{42}
since $P'V$~can be made as small a part of~$PQ$ as we
please, still more can it be made as small a part as we
please of~$PP'$, which is greater than~$PQ$. Therefore
the arc and chord~$PP'$ may be made to have a ratio as
nearly equal to unity as we please. And because $pa$~is
less than~$pv$, and therefore less than~$P'V$, it follows
that $pa$~may be made as small a part as we please of~$PQ$,
and still more of~$PP'$.

In these propositions is contained the rational explanation
of the proposition of Leibnitz, that ``an infinitely
\index{Leibnitz}%
small arc is equal to, and coincides with, its
chord.''


\Subsection{Orders of Infinity.}

Let there be any number of series, arranged in
\index{Infinity, orders of|EtSeq}%
\index{Orders of infinity|EtSeq}%
powers of~$h$, so that the lowest power is first; let
them contain none but whole powers, and let them all
be such, that each will be convergent, on giving to~$h$
a sufficiently small value: as follows,
\begin{alignat*}{4}
Ah + Bh^{2} &{}+{}& Ch^{3} &{}+{}& Dh^{4} &{}+{}& Eh^{5} &+ \etc.
\Tag{(1)} \\
     B'h^{2} &{}+{}& C'h^{3} &{}+{}& D'h^{4} &{}+{}& E'h^{5} &+ \etc.
\Tag{(2)} \\
           &&  C''h^{3} &{}+{}& D''h^{4} &{}+{}& E''h^{5} &+ \etc.
\Tag{(3)} \\
           &&&& D'''h^{4} &{}+{}& E'''h^{5} &+ \etc.
\Tag{(4)} \\
           &&&&&& \etc. &
\tag*{\etc.}
\end{alignat*}

As $h$~is diminished, all these expressions decrease
without limit; but the first \emph{increases} with respect to
the second, that is, contains it more times after a decrease
of~$h$ than it did before. For the ratio of \Eq{(1)}
to~\Eq{(2)} is that of $A + Bh + Ch^{2} + \etc.$ to $B'h + C'h^{2} + \etc.$,
the ratio of the two not being changed by dividing
both by~$h$. The first term of the latter ratio
approximates continually to~$A$, as $h$~is diminished,
and the second can be made as small as we please,
and therefore can be contained in the first as often as
\PageSep{43}
we please. Hence the ratio \Eq{(1)}~to~\Eq{(2)} can be made
as great as we please. By similar reasoning, the ratio
\Eq{(2)}~to~\Eq{(3)}, of \Eq{(3)}~to~\Eq{(4)}, etc., can be made as great as
we please. We have, then, a series of quantities,
each of which, by making $h$ sufficiently small, can be
made as small as we please. Nevertheless this decrease
increases the ratio of the first to the second, of
the second to the third, and so on, and the increase is
without limit.

Again, if we take \Eq{(1)}~and~$h$, the ratio of \Eq{(1)}~to~$h$ is
that of $A + Bh + Ch^{2} + \etc.$ to~$1$, which, by a sufficient
decrease of~$h$, may be brought as near as we
please to that of $A$~to~$1$. But if we take \Eq{(1)}~and~$h^{2}$,
the ratio of \Eq{(1)}~to~$h^{2}$ is that of $A + Bh + \etc.$ to~$h$,
which, by previous reasoning, may be increased without
limit; and the same for any higher power of~$h$.
Hence \Eq{(1)}~is said to be \emph{comparable} to the first power
of~$h$, or \emph{of the first order}, since this is the only power
of~$h$ whose ratio to~\Eq{(1)} tends towards a finite limit.
By the same reasoning, the ratio of \Eq{(2)}~to~$h^{2}$, which is
that of $B' + C'h + \etc.$ to~$1$, continually approaches
that of $B'$~to~$1$; but the ratio \Eq{(2)}~to~$h$, which is that
of $B'h + C'h^{2} + \etc.$ to~$1$, diminishes without limit, as
$h$~is decreased, while the ratio of \Eq{(2)}~to~$h^{2}$, or of $B' + C'h + \etc.$
to~$h$, increases without limit. Hence \Eq{(2)}~is
said to be \emph{comparable} to the second power of~$h$, or \emph{of
the second order}, since this is the only power of~$h$ whose
ratio to~\Eq{(2)} tends towards a finite limit. In the language
of Leibnitz if $h$~be an infinitely small quantity,
\Eq{(1)}~is an infinitely small quantity of the first order,
\Eq{(2)}~is an infinitely small quantity of the second
order, and so on.

We may also add that the ratio of two series of
the same order continually approximates to the ratio
\PageSep{44}
of their lowest terms. For example, the ratio of $Ah^{3} + Bh^{4} + \etc.$
to $A'h^{3} + B'h^{4} + \etc.$ is that of $A + Bh + \etc.$
to $A' + B'h + \etc.$, which, as $h$~is diminished,
continually approximates to the ratio of $A$ to~$A'$, which
is also that of $Ah^{3}$ to~$A'h^{3}$, or the ratio of the lowest
terms. In \Fig{4}, $PQ$~or $dx$ being put in place of~$h$,
$QP'$, or $\phi' x\, dx + \phi'' x\, \dfrac{(dx)^{2}}{2}$, etc., is of the first order,
as are~$PV$, and the chord~$PP'$; while $P'V$, or
$\phi'' x\, \dfrac{(dx)^{2}}{2} + \etc.$, is of the second order.

The converse proposition is readily shown, that if
the ratio of two series arranged in powers of~$h$ continually
approaches to some finite limit as $h$~is diminished,
the two series are of the same order, or the exponent
of the lowest power of~$h$ is the same in both.
Let $Ah^{a}$ and $Bh^{b}$ be the lowest powers of~$h$, whose ratio,
as has just been shown, continually approximates
to the actual ratio of the two series, as $h$~is diminished.
The hypothesis is that the ratio of the two series, and
therefore that of $Ah^{a}$ to~$Bh^{b}$, has a finite limit. This
cannot be if $a > b$, for then the ratio of $Ah^{a}$ to $Bh^{b}$ is
that of $Ah^{a-b}$ to~$B$, which diminishes without limit;
neither can it be when $a < b$, for then the same ratio
is that of $A$ to~$Bh^{b-a}$, which increases without limit;
hence $a$~must be equal to~$b$.

We leave it to the student to prove strictly a proposition
assumed in the preceding; viz., that if the
ratio of $P$~to~$Q$ has unity for its limit, when $h$~is diminished,
the limiting ratio of $P$~to~$R$ will be the same
as the limiting ratio of $Q$~to~$R$. We proceed further
to illustrate the Differential Calculus as applied to
Geometry.
\PageSep{45}


\Subsection[A Geometrical Illustration: Limit of the Intersections of Two Coinciding Straight Lines.]
{A Geometrical Illustration.}

Let $OC$ and~$OD$ (\Fig{5}) be two axes at right angles
to one another, and let a line~$AB$ of given length
be placed with one extremity in each axis. Let this
line move from its first position into that of~$A'B'$ on
one side, and afterwards into that of~$A''B''$ on the
other side, always preserving its first length. The
motion of a ladder, one end of which is against a wall,
and the other on the ground, is an instance.

Let $A'B'$ and $A''B''$ intersect~$AB$ in $P'$~and~$P''$. If
\index{Ladder against wall|EtSeq}%
$A''B''$~were gradually moved from its present position
into that of~$A'B'$, the point~$P''$ would also move gradually
\Figure{5}
from its present position into that of~$P'$, passing,
in its course, through every point in the line~$P'P''$.
But here it is necessary to remark that $AB$~is itself
one of the positions intermediate between $A'B'$ and~$A''B''$,
and when two lines are, by the motion of one
of them, brought into one and the same straight line,
they intersect one another (if this phrase can be here
applied at all) in every point, and all idea of one distinct
point of intersection is lost. Nevertheless $P''$~describes
one part of~$P''P'$ before $A''B''$~has come into
the position~$AB$, and the rest afterwards, when it is
between $AB$ and~$A'B'$.
\PageSep{46}

Let $P$~be the point of separation; then every point
of~$P'P''$, except~$P$, is a real point of intersection of~$AB$,
with one of the positions of~$A''B''$, and when
$A''B''$~has moved very near to~$AB$, the point~$P''$ will
be very near to~$P$; and there is no point so near to~$P$,
that it may not be made the intersection of $A''B''$ and~$AB$,
by bringing the former sufficiently near to the
latter. This point~$P$ is, therefore, the \emph{limit} of the intersections
\index{Intersections, limit of|EtSeq}%
\index{Limit of intersections|EtSeq}%
of $A''B''$~and~$AB$, and cannot be found by
the ordinary application of algebra to geometry, but
may be made the subject of an inquiry similar to those
\Figure[nolabel]{5}
which have hitherto occupied us, in the following
manner:

Let $OA = a$, $OB = b$, $AB = A'B' = A''B'' = l$. Let
$AA' = da$, $BB' = db$, whence $OA' = a + da$, $OB' = b - db$.
We have then $a^{2} + b^{2} = l^{2}$, and $(a + da)^{2} + (b - db)^{2} = l^{2}$;
subtracting the former of which from
the development of the latter, we have
\[
2a\, da + (da)^{2} - 2b\, db + (db)^{2} = 0\Add{,}
\]
or
\[
\frac{db}{da} = \frac{2a + da}{2b - db}\Add{.}
\Tag{(1)}
\]
As $A'B'$ moves towards~$AB$, $da$~and~$db$ are diminished
without limit, $a$~and~$b$ remaining the same; hence the
limit of the ratio~$\dfrac{db}{da}$ is $\dfrac{2a}{2b}$ or~$\dfrac{a}{b}$.
\PageSep{47}

Let the co-ordinates\footnote
  {The lines $OM'$ and $M'P'$ are omitted, to avoid crowding the figure.}
of~$P'$ be $OM' = x$ and $M'P = y$.
Then (\PageRef{32}) the co-ordinates of any point in~$AB$
have the equation
\[
ay + bx = ab\Add{.}
\Tag{(2)}
\]
The point~$P'$ is in this line, and also in the one which
cuts off $a + da$ and $b - db$ from the axes, whence
\[
(a + da)y + (b - db)x = (a + da)(b - db)\Add{;}
\Tag{(3)}
\]
subtract \Eq{(2)} from~\Eq{(3)} after developing the latter, which
gives
\[
y\, da - x\, db = b\, da - a\, db - da\, db\Add{.}
\Tag{(4)}
\]
If we now suppose $A'B'$~to move towards~$AB$, equation~\Eq{(4)}
gives no result, since each of its terms diminishes
without limit. If, however, we divide~\Eq{(4)} by~$da$,
and substitute in the result the value of~$\dfrac{db}{da}$ obtained
from~\Eq{(1)} we have
\[
y - x\, \frac{2a + da}{2b - db}
 = b - a\, \frac{2a + da}{2b - db} - db\Add{.}
\Tag{(5)}
\]
From this and~\Eq{(2)} we might deduce the values of $y$
and~$x$, for the point~$P'$, as the figure actually stands.
Then by diminishing $db$~and $da$ without limit, and
observing the limit towards which $x$~and~$y$ tend, we
might deduce the co-ordinates of~$P$, the limit of the
intersections.

The same result may be more simply obtained, by
diminishing $da$~and~$db$ in equation~\Eq{(5)}, before obtaining
the values of $y$~and~$x$. This gives
\[
y - \frac{a}{b}\, x = b - \frac{a^{2}}{b} \quad\text{or}\quad
by - ax = b^{2} - a^{2}\Add{.}
\Tag{(6)}
\]
From \Eq{(6)}~and~\Eq{(2)} we find (\Fig{6})
\[
x = OM = \frac{a^{3}}{a^{2} + b^{2}} = \frac{a^{3}}{l^{2}} \quad\text{and}\quad
y = MP = \frac{b^{3}}{a^{2} + b^{2}} = \frac{b^{3}}{l^{2}}.
\]
\PageSep{48}

This limit of the intersections is different for every
different position of the line~$AB$, but may be determined,
in every case, by the following simple construction.

Since (\Fig{6}) $BP: PN$, or $OM :: BA : AO$, we
have $BP = OM\, \dfrac{BA}{AO} = \dfrac{a^{3}}{l^{2}}\, \dfrac{l}{a} = \dfrac{a^{2}}{l}$; and, similarly,
$PA = \dfrac{b^{2}}{l}$. Let $OQ$~be drawn perpendicular to~$BA$;
then since $OA$~is a mean proportional between $AQ$
and~$AB$, we have $AQ = \dfrac{a^{2}}{l}$, and similarly $BQ = \dfrac{b^{2}}{l}$.
Hence $BP = AQ$ and $AP = BQ$, or the point~$P$ is
as far from either extremity of~$AB$ as $Q$~is from the
other.
\Figure{6}


\Subsection{The Same Problem Solved by the Principles of
Leibnitz.}

We proceed to solve the same problem, using the
\index{Leibnitz}%
principles of Leibnitz, that is, supposing magnitudes
can be taken so small, that those proportions may be
regarded as absolutely correct, which are not so in
reality, but which only approach more nearly to the
truth, the smaller the magnitudes are taken. The inaccuracy
of this supposition has been already pointed
out; yet it must be confessed that this once got over,
\PageSep{49}
the results are deduced with a degree of simplicity
and consequent clearness, not to be found in any other
method. The following cannot be regarded as a demonstration,
except by a mind so accustomed to the
subject that it can readily convert the various inaccuracies
into their corresponding truths, and see, at one
glance, how far any proposition will affect the final
result. The beginner will be struck with the extraordinary
assertions which follow, given in their most
naked form, without any attempt at a less startling
mode of expression.
\Figure{7}

Let $A'B'$ (\Fig{7}) be a position of~$AB$ infinitely
\index{Infinitely small, the notion of}%
near to it; that is, let $A'PA$~be an infinitely small
angle. With the centre~$P$, and the radii $PA'$ and~$PB$,
describe the infinitely small arcs $A'a$,~$Bb$. An infinitely
small arc of a circle is a straight line perpendicular
to its radius; hence $A'aA$~and~$BbB'$ are right-angled
triangles, the first similar to~$BOA$, the two
having the angle~$A$ in common, and the second similar
to~$B'OA'$. Again, since the angles of~$BOA$, which
are finite, only differ from those of~$B'OA'$ by the infinitely
small angle~$A'PA$, they may be regarded as
\PageSep{50}
equal; whence $A'aA$~and~$B'bB$ are similar to~$BOA$,
and to one another. Also $P$~is the point of which we
are in search, or infinitely near to it; and since $BA = B'A'$,
of which $BP = bP$ and $aP = A'P$, the remainders
$B'b$~and~$Aa$ are equal. Moreover, $Bb$~and~$A'a$
being arcs of circles subtending equal angles, are in
the proportion of the radii $BP$~and~$PA'$.

Hence we have the following proportions:
\begin{gather*}
Aa : A'a :: OA : OB :: a : b \\
Bb : B'b :: OA : OB :: a : b\rlap{.}
\end{gather*}
The composition of which gives, since $Aa = B'b$:
\[
\PadTo[r]{BP + Pa}{Bb} : A'a :: \PadTo[l]{a^{2} + b^{2}}{a^{2}} : b^{2}.
\]
Also
\[
\PadTo[r]{BP + Pa}{Bb} : A'a :: \PadTo[l]{a^{2} + b^{2}}{BP} : Pa,
\]
whence
\[
\PadTo[r]{BP + Pa}{BP} : Pa :: \PadTo[l]{a^{2} + b^{2}}{a^{2}} : b^{2},
\]
and
\[
BP + Pa : Pa :: a^{2} + b^{2} : b^{2}.
\]
But $Pa$~only differs from~$PA$ by the infinitely small
quantity~$Aa$, and $BP + PA = l$, and $a^{2} + b^{2} = l^{2}$;
whence
\[
l : PA :: l^{2} : b^{2},\quad\text{or}\quad PA = \frac{b^{2}}{l},
\]
which is the result already obtained.

In this reasoning we observe four independent
errors, from which others follow: (1)~that $Bb$~and~$A'a$
are straight lines at right-angles to~$Pa$; (2)~that $BOA$\Add{,}~$B'OA'$
are similar triangles; (3)~that $P$~is really the
point of which we are in search; (4)~that $PA$~and~$Pa$
are equal. But at the same time we observe that
every one of these assumptions approaches the truth,
as we diminish the angle~$A'PA$, so that there is no
magnitude, line or angle, so small that the linear or
angular errors, arising from the above-mentioned suppositions,
may not be made smaller.

We now proceed to put the same demonstration
\PageSep{51}
in a stricter form, so as to neglect no quantity during
the process. This should always be done by the beginner,
until he is so far master of the subject as to be
able to annex to the inaccurate terms the ideas necessary
for their rational explanation. To the former figure
add $B\beta$ and~$A\alpha$, the real perpendiculars, with
which the arcs have been confounded. Let $\angle A'PA = d\theta$,
\index{Angle, unit employed in measuring an}%
$PA = p$, $Aa = dp$, $BP = q$, $B'b = dq$; and $OA = a$,
$OB = b$, and $AB = l$. Then\footnote
  {For the unit employed in measuring an angle, see \Title{Study of Mathematics}
  (Chicago, 1898), pages 273--277.}
$A'a = (p - dp)\, d\theta$, $Bb = q\, d\theta$,
and the triangles $A'A\alpha$ and $B'B\beta$ are similar to
\Figure[nolabel]{7}
$BOA$~and~$B'OA'$. The perpendiculars $A'\alpha$ and~$B\beta$
are equal to $PA' \sin d\theta$ and $PB \sin d\theta$, or $(p - dp) \sin d\theta$
and $q \sin d\theta$. Let $a\alpha = \mu$ and $b\beta = \nu$. These
(\PageRef[p.]{9}) will diminish without limit as compared with
$A'\alpha$ and~$B\beta$; and since the ratios of $A'\alpha$ to~$\alpha A$ and $B\beta$
to~$\beta B'$ continue finite (these being sides of triangles
similar to $AOB$ and~$A'OB'$), $a\alpha$~and~$b\beta$ will diminish
indefinitely with respect to $\alpha A$~and~$\beta B'$. Hence the
ratio $A\alpha$ to~$\beta B'$ or $dp + \mu$ to $dq + \nu$ will continually
approximate to that of $dp$ to~$dq$, or a ratio of equality.
\PageSep{52}

The exact proportions, to which those in the last
page are approximations, are as follows:
\begin{alignat*}{3}
dp + \mu &: (p - dp) \sin d\theta &&:: a &&: b, \\
q \sin d\theta &: \PadTo{(p - dp) \sin d\theta}{dq + \nu} &&:: a - da &&: b + db;
\end{alignat*}
by composition of which, recollecting that $dp = dq$
(which is rigorously true) and dividing the two first
terms of the resulting proportion by~$dp$, we have
\[
q\left(1 + \frac{\mu}{dp}\right) : (p - dp)\left(1 + \frac{\nu}{dp}\right)
  :: a(a - da) : b(b + db).
\]

If $d\theta$ be diminished without limit, the quantities
$da$,~$db$, and~$dp$, and also the ratios $\dfrac{\mu}{dp}$ and~$\dfrac{\nu}{dp}$, as
above-mentioned, are diminished without limit, so
that the limit of the proportion just obtained, or the
proportion which gives the limits of the lines into
which $P$~divides~$AB$, is
\begin{alignat*}{3}
q &: p &&:: a^{2} &&: b^{2}, \\
\intertext{hence}
q + p = l &: p &&:: a^{2} + b^{2} = l^{2} &&: b^{2},
\end{alignat*}
the same as before.


\Subsection[An Illustration from Dynamics: Velocity, Acceleration, etc.]
{An Illustration from Dynamics.}

We proceed to apply the preceding principles to
dynamics, or the theory of motion.

Suppose a point moving along a straight line uniformly;
that is, if the whole length described be divided
into any number of equal parts, however great,
each of those parts is described in the same time.
Thus, whatever length is described in the first second
of time, or in any part of the first second, the same
is described in any other second, or in the same part
of any other second. The number of units of length
described in a unit of time is called the \emph{velocity}; thus
\index{Velocity!linear|EtSeq}%
a velocity of $3.01$~feet in a second means that the
\PageSep{53}
point describes three feet and one hundredth in each
second, and a proportional part of the same in any
part of a second. Hence, if $v$~be the velocity, and
$t$~the units of time elapsed from the beginning of the
motion, $vt$~is the length described; and if any length
described be known, the velocity can be determined
by dividing that length by the time of describing it.
Thus, a point which moves uniformly through $3$~feet
in $1\frac{1}{2}$~second, moves with a velocity of $3 ÷ 1\frac{1}{2}$, or $2$~feet
per~second.

Let the point not move uniformly; that is, let different
\index{Continuous quantities}%
\index{Quantities, continuous}%
parts of the line, having the same length, be
described in different times; at the same time let the
motion be \emph{continuous}, that is, not suddenly increased
or decreased, as it would be if the point were composed
of some hard matter, and received a blow while
it was moving. This will be the case if its motion be
represented by some algebraical function of the time,
or if, $t$~being the number of units of time during which
the point has moved, the number of units of length
described can be represented by~$\phi t$. This, for example,
we will suppose to be~$t + t^{2}$, the unit of time
being one second, and the unit of length one inch;
so that $\frac{1}{2} + \frac{1}{4}$, or $\frac{3}{4}$~of an inch, is described in the first
half second; $1 + 1$, or two inches, in the first second;
$2 + 4$, or six inches, in the first two seconds, and so on.

Here we have no longer an evident measure of the
velocity of the point; we can only say that it obviously
increases from the beginning of the motion to
the end, and is different at every two different points.
Let the time~$t$ elapse, during which the point will describe
the distance $t + t^{2}$; let a further time~$dt$ elapse,
during which the point will increase its distance to
$t + dt + (t + dt)^{2}$, which, diminished by~$t + t^{2}$, gives
\PageSep{54}
$dt + 2t\, dt + (dt)^{2}$ for the length described during the
increment of time~$dt$. This varies with the value of~$t$;
thus, in the interval~$dt$ after the first second, the
length described is $3\, dt + dt^{2}$; after the second second,
it is $5\, dt + (dt)^{2}$, and so on. Nor can we, as in the
case of uniform motion, divide the length described,
by the time, and call the result the velocity with which
that length is described; for no length, however small,
is here uniformly described. If we were to divide a
length by the time in which it is described, and also
its first and second halves by the times in which they
are respectively described, the three results would be
all different from one another.

Here a difficulty arises, similar to that already noticed,
when a point moves along a curve; in which,
as we have seen, it is improper to say that it is moving
in any one direction through an arc, however
small. Nevertheless a straight line was found at every
point, which did, more nearly than any other straight
line, represent the direction of the motion. So, in
this case, though it is incorrect to say that there is
any uniform velocity with which the point continues
to move for any portion of time, however small, we
can, at the end of every time, assign a uniform velocity,
which shall represent, more nearly than any
other, the rate at which the point is moving. If we
say that, at the end of the time~$t$, the point is moving
with a velocity~$v$, we must not now say that the length~$v\, dt$
is described in the succeeding interval of time~$dt$;
but we mean that $dt$~may be taken so small, that $v\, dt$~shall
bear to the distance actually described a ratio as
near to equality as we please.

Let the point have moved during the time~$t$, after
which let successive intervals of time elapse, each
\PageSep{55}
\index{Coefficients, differential}%
\index{Differential coefficients}%
equal to~$dt$. At the end of the times, $t$,~$t + dt$, $t + 2\, dt$,
$t + 3\, dt$,~etc., the whole lengths described will be $t + t^{2}$,
$t + dt + (t + dt)^{2}$, $t + 2\, dt + (t + 2\, dt)^{2}$, $t + 3\, dt + (t + 3\, dt)^{2}$,
etc.; the differences of which, or $dt + 2t\, dt + (dt)^{2}$,
$dt + 2t\, dt + 3(dt)^{2}$, $dt + 2t\, dt + 5(dt)^{2}$, etc.,
are the lengths described in the first, second, third,
etc., intervals~$dt$. These are not equal to one another,
as would be the case if the velocity were uniform; but
by making $dt$ sufficiently small, their ratio may be
brought as near to equality as we please, since the
terms $(dt)^{2}$,~$3(dt)^{2}$,~etc., by which they all differ from
the common part $(1 + 2t)\, dt$, may be made as small as
we please, in comparison of this common part. If we
divide the above-mentioned lengths by~$dt$, which does
not alter their ratio, they become $1 + 2t + dt$, $1 + 2t + 3\, dt$,
$1 + 2t + 5\, dt$, etc., which may be brought as
near as we please to equality, by sufficient diminution
of~$dt$. Hence $1 + 2t$ is said to be the velocity of the
point after the time~$t$; and if we take a succession of
equal intervals of time, each equal to~$dt$, and sufficiently
small, the lengths described in those intervals
will bear to $(1 + 2t)\, dt$, the length which would be described
in the same interval with the uniform velocity
$1 + 2t$, a ratio as near to equality as we please. And
observe, that if $\phi t$ is $t + t^{2}$, $\phi' t$~is $1 + 2t$, or the coefficient
of~$h$ in $(t + h) + (t + h)^{2}$.

In the same way it may be shown, that if the point
moves so that $\phi t$~always represents the length described
in the time~$t$, the differential coefficient of~$\phi t$\Add{,}
or~$\phi' t$, is the velocity with which the point is moving
at the end of the time~$t$. For the time~$t$ having elapsed,
the whole lengths described at the end of the times $t$
and $t + dt$ are $\phi t$ and $\phi(t + dt)$; whence the length
described during the time~$dt$ is
\PageSep{56}
\[
\phi(t + dt) - \phi t, \quad\text{or}\quad
\phi't\, dt + \phi'' t\, \frac{(dt)^{2}}{2} + \etc.
\]
Similarly, the length described in the next interval
$dt$ is
\begin{gather*}
\phi(t + 2\, dt) - \phi(t + dt); \quad\text{or}, \displaybreak[0] \\
  \phi t + \phi' t\, 2\, dt + \phi'' t\, \frac{(2\, dt)^{2}}{2} + \etc. \displaybreak[0] \\
-(\phi t + \phi' t\, dt + \phi'' t\, \frac{(dt)^{2}}{2} + \etc.), \displaybreak[0] \\
\intertext{which is}
\phi' t\, dt + 3\phi'' t\, \frac{(dt)^{2}}{2} + \etc.;
\end{gather*}
the length described in the third interval~$dt$ is
$\phi' t\, dt + 5\phi'' t\, \dfrac{(dt)^{2}}{2} + \etc.$,~etc.

Now, it has been shown for each of these, that the
first term can be made to contain the aggregate of all
the rest as often as we please, by making $dt$ sufficiently
small; this first term is $\phi' t\, dt$ in all, or the length
which would be described in the time~$dt$ by the velocity
$\phi' t$ continued uniformly: it is possible, therefore,
to take $dt$ so small, that the lengths actually described
in a succession of intervals equal to~$dt$, shall be as
nearly as we please in a ratio of equality with those
described in the same intervals of time by the velocity~$\phi' t$.
For example, it is observed in bodies which fall
\index{Falling bodies}%
to the earth from a height above it, when the resistance
of the air is removed, that if the time be taken
in seconds, and the distance in feet, the number of
feet fallen through in $t$~seconds is always~$at^{2}$, where
$a = 16\frac{1}{12}$ very nearly; what is the velocity of a body
which has fallen \textit{in~vacuo} for four seconds? Here $\phi t$
being~$at^{2}$, we find, by substituting $t + h$, or $t + dt$, instead
of~$t$, that $\phi' t$~is~$2at$, or $2 × 16\frac{1}{12} × t$; which, at
\PageSep{57}
the end of four seconds, is $32\frac{1}{6} × 4$, or $128\frac{2}{3}$~feet. That
is, at the end of four seconds a falling body moves at
the rate of $128\frac{2}{3}$~feet per~second. By which we do
not mean that it continues to move with this velocity
for any appreciable time, since the rate is always
varying; but that the length described in the interval~$dt$
after the fourth second, may be made as nearly as
we please in a ratio of equality with $128\frac{2}{3} × dt$, by
taking $dt$ sufficiently small. This velocity~$2at$ is said
to be \emph{uniformly} accelerated; since in each second the
\index{Accelerated motion}%
\index{Uniformly accelerated}%
same velocity~$2a$ is gained. And since, when $x$~is the
space described, $\phi' t$~is the limit of~$\dfrac{dx}{dt}$, the velocity is
also this limit; that is, when a point does not move
uniformly, the velocity is not represented by any increment
of length divided by its increment of time,
but by the limit to which that ratio continually tends,
as the increment of time is diminished.


\Subsection{Simple Harmonic Motion.}

We now propose the following problem: A point
\index{Motion!simple harmonic}%
\index{Simple harmonic motion}%
moves uniformly round a circle; with what velocities
do the abscissa and ordinate increase or decrease, at
any given point? (\Fig{8}.)

Let the point~$P$, setting out from~$A$, describe the
arc~$AP$, etc., with the uniform velocity of $a$~inches
per~second. Let $OA = r$, $\angle A0P = \theta$, $\angle POP' = d\theta$,
$0M = x$, $MP = y$, $MM' = dx$, $QP' = dy$.

From the first principles of trigonometry
\begin{alignat*}{4}%[** TN: Re-aligned from the original]
&x      &&= r \cos\theta\Add{,} \\
&x - dx &&= r \cos(\theta + d\theta)
        &&= r \cos\theta \cos d\theta - r \sin\theta \sin d\theta\Add{,} \displaybreak[0] \\
&y      &&= r \sin\theta\Add{,} \\
&y + dy &&= r \sin(\theta + d\theta)
        &&= r \sin\theta \cos d\theta + r \cos\theta \sin d\theta.
\end{alignat*}
\PageSep{58}
Subtracting the second from the first, and the third
from the fourth, we have
\begin{alignat*}{2}
dx &= r \sin\theta \sin d\theta + r \cos\theta(1 - \cos d\theta)\Add{,}
\Tag{(1)} \\
dy &= r \cos\theta \sin d\theta + r \sin\theta(1 - \cos d\theta)\Add{.}
\Tag{(2)} \\
\end{alignat*}
But if $d\theta$ be taken sufficiently small, $\sin d\theta$, and~$d\theta$,
may be made as nearly in a ratio of equality as we
please, and $1 - \cos d\theta$ may be made as small a part
as we please, either of $d\theta$ or $\sin d\theta$. These follow from
\Fig{1}, in which it was shown that $BM$ and the arc~$BA$,
or (if $OA = r$ and $AOB = d\theta$), $r \sin d\theta$ and~$r\, d\theta$,
may be brought as near to a ratio of equality as we
\Figure{8}
please, which is therefore true of $\sin d\theta$ and~$d\theta$. Again,
it was shown that~$AM$, or $r - r \cos d\theta$, can be made
as small a part as we please, either of~$BM$ or the arc~$BA$,
that is, either of $r \sin d\theta$, or~$r\, d\theta$; the same is
therefore true of $1 - \cos d\theta$, and either $\sin d\theta$ or~$d\theta$.
Hence, if we write equations \Eq{(1)}~and~\Eq{(2)} thus,
\[
dx = r \sin\theta\, d\theta\quad (1)\qquad\qquad
dy = r \cos\theta\, d\theta\quad (2),
\]
we have equations, which, though never exactly true,
are such that by making $d\theta$ sufficiently small, the
errors may be made as small parts of~$d\theta$ as we please.
Again, since the arc~$AP$ is uniformly described, so
also is the angle~$POA$; and since an arc~$a$ is described
\PageSep{59}
in one second, the angle~$\dfrac{a}{r}$ is described in the same
\index{Velocity!angular}%
time; this is, therefore, the \emph{angular velocity}.\footnote
  {The same considerations of velocity which have been applied to the
  motion of a point along a line may also be applied to the motion of a line
  round a point. If the angle so described be always increased by equal angles
  in equal portions of time, the angular velocity is said to be uniform, and is
  measured by the number of angular units described in a unit of time. By
  similar reasoning to that already described, if the velocity with which the
  angle increases be not uniform, so that at the end of the time~$t$ the angle described
  is $\theta = \phi t$, the angular velocity is~$\phi' t$, or the limit of the ratio~$\dfrac{d\theta}{dt}$.}
If we
divide equations \Eq{(1)}~and~\Eq{(2)} by~$dt$, we have
\[
%[** TN: Signs OK; De Morgan absorbs the - in dx/dt at the bottom of p. 57]
\frac{dx}{dt} = r \sin\theta\, \frac{d\theta}{dt}\qquad
\frac{dy}{dt} = r \cos\theta\, \frac{d\theta}{dt};
\]
these become more nearly true as $dt$~and~$d\theta$ are diminished,
so that if for $\dfrac{dx}{dt}$,~etc., the limits of these ratios
be substituted, the equations will become rigorously
true. But these limits are the velocities of $x$,~$y$, and~$\theta$,
the last of which is also~$\dfrac{a}{r}$; hence
\begin{alignat*}{2}
\text{velocity of~$x$} &= r \sin\theta × \frac{a}{r} &&= a \sin\theta, \\
\text{velocity of~$y$} &= r \cos\theta × \frac{a}{r} &&= a \cos\theta;
\end{alignat*}
that is, the point~$M$ moves towards~$O$ with a variable
velocity, which is always such a part of the velocity
of~$P$, as $\sin\theta$~is of unity, or as $PM$~is of~$OB$; and the
distance~$PM$ increases, or the point~$N$ moves from~$O$,
with a velocity which is such a part of the velocity of~$P$
as $\cos\theta$~is of unity, or as $OM$~is of~$OA$. [The motion
of the point~$M$ or the point~$N$ is called in physics
a \emph{simple harmonic motion}.]

In the language of Leibnitz, the results of the two
\index{Leibnitz}%
foregoing sections would be expressed thus: If a
point move, but not uniformly, it may still be considered
as moving uniformly for any infinitely small
\index{Infinitely small, the notion of}%
\PageSep{60}
time; and the velocity with which it moves is the infinitely
small space thus described, divided by the infinitely
small time.


\Subsection{The Method of Fluxions.}

The foregoing process contains the method employed
\index{Fluxions}%
by Newton, known by the name of the \emph{Method
\index{Newton}%
of Fluxions}. If we suppose $y$ to be any function of~$x$,
and that $x$~increases with a given velocity, $y$~will also
increase or decrease with a velocity depending: (1)~upon
the velocity of~$x$; (2)~upon the function which
$y$ is of~$x$. These velocities Newton called the fluxions
of $y$~and~$x$, and denoted them by $\dot{y}$~and~$\dot{x}$. Thus, if
$y = x^{2}$, and if in the interval of time~$dt$, $x$~becomes
$x + dx$, and $y$~becomes $y + dy$, we have $y + dy = (x + dx)^{2}$,
and $dy = 2x\, dx + (dx)^{2}$, or $\dfrac{dy}{dt} = 2x\, \dfrac{dx}{dt} + \dfrac{dx}{dt}\, dx$.
If we diminish~$dt$, the term $\dfrac{dx}{dt}\, dx$ will diminish
without limit, since one factor continually approaches
to a given quantity, viz., the velocity of~$x$,
and the other diminishes without limit. Hence we
obtain the velocity of $y = 2x × \text{the velocity of~$x$}$, or
$\dot{y} = 2x\, \dot{x}$, which is used in the method of fluxions instead
of $dy = 2x\, dx$ considered in the manner already
described. The processes are the same in both methods,
since the ratio of the velocities is the limiting
ratio of the corresponding increments, or, according
to Leibnitz, the ratio of the infinitely small increments.
\index{Leibnitz}%
We shall hereafter notice the common objection
to the Method of Fluxions.


\Subsection{Accelerated Motion.}

When the velocity of a material point is suddenly
\index{Accelerated motion}%
\index{Motion!accelerated}%
\index{Uniformly accelerated}%
increased, an \emph{impulse} is said to be given to it, and the
\index{Impulse}%
\PageSep{61}
magnitude of the impulse or impulsive force is in proportion
\index{Force|(}%
to the velocity created by it. Thus, an impulse
which changes the velocity from $50$ to $70$~feet
per~second, is twice as great as one which changes it
from $50$ to $60$~feet. When the velocity of the point is
altered, not suddenly but continuously, so that before
the velocity can change from $50$ to $70$~feet, it goes
through all possible intermediate velocities, the point
is said to be acted on by an \emph{accelerating force}. \emph{Force}
is a name given to that which causes a change in the
velocity of a body. It is said to act uniformly, when
the velocity acquired by the point in any one interval
of time is the same as that acquired in any other interval
of equal duration. It is plain that we cannot,
by supposing any succession of impulses, however
small, and however quickly repeated, arrive at a uniformly
accelerated motion; because the length described
between any two impulses will be uniformly
described, which is inconsistent with the idea of continually
accelerated velocity. Nevertheless, by diminishing
the magnitude of the impulses, and increasing
their number, we may come as near as we please
to such a continued motion, in the same way as, by
diminishing the magnitudes of the sides of a polygon,
and increasing their number, we may approximate as
near as we please to a \Typo{continous}{continuous} curve.

Let a point, setting out from a state of rest, increase
its velocity uniformly, so that in the time~$t$, it
may acquire the velocity~$v$---what length will have
been described during that time~$t$? Let the time~$t$
and the velocity~$v$ be both divided into $n$~equal parts,
each of which is $t'$ and~$v'$, so that $nt' = t$, and $nv' = v$.
Let the velocity~$v'$ be communicated to the point at
rest; after an interval of~$t'$ let another velocity~$v'$ be
\PageSep{62}
communicated, so that during the second interval~$t'$
the point has a velocity~$2v'$; during the third interval
let the point have the velocity~$3v'$, and so on; so that
in the last or $n$\th~interval the point has the velocity~$nv'$.
The space described in the first interval is, therefore,~$v't'$;
in the second,~$2v't'$; in the third~$3v't'$; and
so on, till in the $n$\th~interval it is~$nv't'$. The whole
space described is, therefore,
\[
v't' + 2v't' + 3v't' + \dots + (n - 1)v't' + nv't'\Add{,}
\]
or
\[
[1 + 2 + 3 \Add{+} \dots + (n - 1) + n]v't'
  = n · \frac{(n + 1)}{2}\, v't'
  = \frac{n^{2} v't' + nv't'}{2}.
\]
In this substitute $v$ for~$nv'$, and $t$ for~$nt'$, which gives
for the space described $\frac{1}{2}v(t + t')$. The smaller we
suppose~$t'$, the more nearly will this approach to~$\frac{1}{2}vt$.
But the smaller we suppose~$t'$, the greater must be~$n$,
the number of parts into which $t$~is divided; and the
more nearly do we render the motion of the point uniformly
accelerated. Hence the limit to which we approximate
by diminishing~$t'$ without limit, is the length
described in the time~$t$ by a uniformly accelerated
velocity, which shall increase from~$0$ to~$v$ in that time.
This is~$\frac{1}{2}vt$, or half the length which would have been
described by the velocity~$v$ continued uniformly from
the beginning of the motion.

It is usual to measure the accelerating force by the
\index{Accelerating force}%
velocity acquired in one second. Let this be~$g$; then
since the same velocity is acquired in every other second,
the velocity acquired in $t$~seconds will be~$gt$, or
$v = gt$. Hence the space described is $\frac{1}{2}gt × t$, or~$\frac{1}{2}gt^{2}$.
If the point, instead of being at rest at the beginning
of the acceleration, had had the velocity~$a$, the lengths
\PageSep{63}
described in the successive intervals would have been
$at' + v't'$, $at' + 2v't'$, etc.; so that to the space described
by the accelerated motion would have been added~$nat'$,
or~$at$, and the whole length would have been
$at + \frac{1}{2}gt^{2}$. By similar reasoning, had the force been
a uniformly \emph{retarding} force, that is, one which diminished
\index{Force|)}%
the initial velocity~$a$ equally in equal times, the
length described in the time~$t$ would have been $at - \frac{1}{2}gt^{2}$.

Now let the point move in such a way, that the
velocity is accelerated or retarded, but not uniformly;
that is, in different times of equal duration, let different
velocities be lost or gained. For example, let the
point, setting out from a state of rest, move in such a
\Figure{9}
way that the number of inches passed over in $t$~seconds
is always~$t^{3}$. Here $\phi t = t^{3}$, and the velocity acquired
by the body at the end of the time~$t$, is the coefficient
of~$dt$ in $(t + dt)^{3}$, or $3t^{2}$~inches per~second.
Let the point (\Fig{9}) be at~$A$ at the end of the time~$t$;
and let $AB$,~$BC$, $CD$,~etc., be lengths described in
successive equal intervals of time, each of which is~$dt$.
Then the velocities at $A$,~$B$,~$C$,~etc., are $3t^{2}$, $3(t + dt)^{2}$,
$3(t + 2\, dt)^{3}$, etc., and the lengths $AB$,~$BC$, $CD$,~etc.,
are $(t + dt)^{3} - t^{3}$, $(t + 2\, dt)^{2} - (t + dt)^{3}$, $(t + 3\, dt)^{3} - (t + 2\, dt)^{3}$,
etc.
\[
\ArrayCompress
\begin{array}{cl}
\ColHead{VELOCITY AT} & \\
A & 3t^{2}\Add{,} \\
B & 3t^{2} + \Z6t\, dt + \Z3(dt)^{2}\Add{,} \\
C & 3t^{2} +  12t\, dt +  12(dt)^{2}\Add{,} \\
\PageSep{64}
\ColHead{LENGTH OF} & \\
AB & 3t^{2}\, dt + \Z3t(dt)^{2} + \Z\Z(dt)^{3}\Add{,} \\
BC & 3t^{2}\, dt + \Z9t(dt)^{2} + \Z 7(dt)^{3}\Add{,} \\
CD & 3t^{2}\, dt +  15t(dt)^{2} +   19(dt)^{3}\Add{.}
\end{array}
\]

If we could, without error, reject the terms containing~$(dt)^{2}$
in the velocities, and those containing~$(dt)^{3}$
in the lengths, we should then reduce the motion
of the point to the case already considered, the
initial velocity being~$3t^{2}$, and the accelerating force~$6t$.
For we have already shown that $a$~being the initial
velocity, and $g$~the accelerating force, the space described
in the time~$t$ is $at + \frac{1}{2}gt^{2}$. Hence, $3t^{2}$~being
the initial velocity, and $6t$~the accelerating force, the
space in the time~$dt$ is $3t^{2}\, dt + 3t(dt)^{2}$, which is the
same as~$AB$ after $(dt)^{3}$~is rejected. The velocity acquired
is~$gt$, and the whole velocity is, therefore,
$a + gt$, or making the same substitutions $3t^{2} + 6t\, dt$.
This is the velocity at~$B$, after the term~$3(dt)^{2}$ is
rejected. Again, the velocity being $3t^{2} + 6t\, dt$, and
the force~$6t$, the space described in the time~$dt$ is
$(3t^{2} + 6t\, dt)\, dt + 3t(dt)^{2}$, or $3t^{2}\, dt + 9t(dt)^{2}$. This is
what the space~$BC$ becomes after $7(dt)^{3}$~is rejected.
The velocity acquired is~$6t\, dt$; and the whole velocity
is $3t^{2} + 6t\, dt + 6t\, dt$, or $3t^{2} + 12t\, dt$; which is the velocity
at~$C$ after $12(dt)^{2}$~is rejected.

But as the terms involving $(dt)^{2}$ in the velocities,
etc., cannot be rejected without error, the above supposition
of a uniform force cannot be made. Nevertheless,
as we may take $dt$ so small that these terms
shall be as small parts as we please of those which
precede, the results of the erroneous and correct suppositions
may be brought as near to equality as we
please; hence we conclude, that though there is no
force, which, continued uniformly, would preserve
\PageSep{65}
the motion of the point~$A$, so that $OA$~should always
be~$t^{2}$ in inches, yet an interval of time may be taken
so small, that the length actually described by~$A$ in
that time, and the one which would be described if
the force~$6t$ were continued uniformly, shall have a
ratio as near to equality as we please. Hence, on a
principle similar to that by which we called~$3t^{3}$ the
velocity at~$A$, though, in truth, no space, however
small, is described with that velocity, we call~$6t$ the
accelerating force at~$A$. And it must be observed
that $6t$~is the differential coefficient of~$3t^{2}$, or the coefficient
of~$dt$, in the development of~$3(t + dt)^{2}$.

Generally, let the point move so that the length
described in any time~$t$ is~$\phi t$. Hence the length described
at the end of the time $t + dt$ is $\phi(t + dt)$, and
that described in the interval~$dt$ is $\phi(t + dt) - \phi t$, or
\[
\phi' t\, dt + \phi'' t\, \frac{(dt)^{2}}{2}
             + \phi''' t\, \frac{(dt)^{3}}{2·3} + \etc.\Add{,}
\]
in which $dt$ may be taken so small, that either of the
first two terms shall contain the aggregate of all the
rest, as often as we please. These two first terms are
$\phi' t\, dt + \frac{1}{2}\phi'' t (dt)^{2}$, and represent the length described
during~$dt$, with a uniform velocity~$\phi' t$, and an accelerating
force~$\phi'' t$. The interval~$dt$ may then generally
be taken so small, that this supposition shall represent
the motion during that interval as nearly as we please.


\Subsection{Limiting Ratios of Magnitudes that Increase
Without Limit.}

We have hitherto considered the limiting ratio of
\index{Increase without limit|EtSeq}%
\index{Limiting ratios|EtSeq}%
\index{Ratios, limiting|EtSeq}%
quantities only as to their state of \emph{decrease}: we now
proceed to some cases in which the limiting ratio of
different magnitudes which \emph{increase} without limit is
investigated.
\PageSep{66}

It is easy to show that the increase of two magnitudes
may cause a decrease of their ratio; so that, as
the two increase without limit, their ratio may diminish
without limit. The limit of any ratio may be found
by rejecting any terms or aggregate of terms~($Q$) which
are connected with another term~($P$) by the sign of
addition or subtraction, provided that by increasing~$x$,
$Q$~may be made as small a part of~$P$ as we please.
For example, to find the limit of $\dfrac{x^{2} + 2x + 3}{2x^{2} + 5x}$, when
$x$~is increased without limit. By increasing~$x$ we can,
as will be shown immediately, cause $2x + 3$ and~$5x$ to
be contained in $x^{2}$ and~$2x^{2}$, as often as we please; rejecting
these terms, we have $\dfrac{x^{2}}{2x^{2}}$, or~$\frac{1}{2}$, for the limit.

The demonstration is as follows: Divide both
numerator and denominator by~$x^{2}$, which gives $1 + \dfrac{2}{x} + \dfrac{3}{x^{2}}$,
and $2 + \dfrac{5}{x}$, for the numerator and denominator
of a fraction equal in value to the one proposed.
These can be brought as near as we please to $1$~and~$2$
by making $x$ sufficiently great, or $\dfrac{1}{x}$~sufficiently small;
and, consequently, their ratio can be brought as near
as we please to~$\dfrac{1}{2}$.

We will now prove the following: That in any
series of decreasing powers of~$x$, any one term will, if
$x$~be taken sufficiently great, contain the aggregate of
all which follow, as many times as we please. Take,
for example,
\[
% [** TN: On two lines in the original]
ax^{m} + bx^{m-1} + cx^{m-2} + \dots + px + q
  + \frac{r}{x} + \frac{s}{x^{2}} + \etc.
\]
\PageSep{67}
The ratio of the several terms will not be altered if we
divide the whole by~$x^{m}$, which gives
\[
a + \frac{b}{x} + \frac{c}{x^{2}} + \dots
  + \frac{p}{x^{m-1}} + \frac{q}{x^{m}}  + \frac{r}{x^{m+1}}
  + \frac{s}{x^{m+2}} + \etc.
\]
It has been shown that by taking $\dfrac{1}{x}$ sufficiently small,
that is, by taking $x$ sufficiently great, any term of this
series may be made to contain the aggregate of the
succeeding terms, as often as we please; which relation
is not altered if we multiply every term by~$x^{m}$,
and so restore the original series.

It follows from this, that $\dfrac{(x + 1)^{m}}{x^{m}}$ has unity for its
limit when $x$~is increased without limit. For $(x + 1)^{m}$
is $x^{m} + mx^{m-1} + \etc.$, in which $x^{m}$~can be made as
great as we please with respect to the rest of the
series. Hence $\dfrac{(x + 1)^{m}}{x^{m}} = 1 + \dfrac{mx^{m-1} + \etc.}{x^{m}}$, the numerator
of which last fraction decreases indefinitely
as compared with its denominator.

In a similar way it may be shown that the limit of
$\dfrac{x^{m}}{(x + 1)^{m+1} - x^{m+1}}$, when $x$~is increased, is~$\dfrac{1}{m + 1}$. For
since $(x + 1)^{m+1} = x^{m+1} + (m + 1)x^{m} + \frac{1}{2}(m + 1)m x^{m-1} + \etc.$,
this fraction is
\[
\frac{x^{m}}{(m + 1)x^{m} + \frac{1}{2}(m + 1)m x^{m-1} + \etc.}
\]
in which the first term of the denominator may be
made to contain all the rest as often as we please;
that is, if the fraction be written thus, $\dfrac{x^{m}}{(m + 1)x^{m} + A}$,
$A$~can be made as small a part of~$(m + 1)x^{m}$ as we
\PageSep{68}
please. Hence this fraction can, by a sufficient increase
of~$x$, be brought as near as we please to
$\dfrac{x^{m}}{(m + 1)x^{m}}$, or~$\dfrac{1}{m + 1}$.

A similar proposition may be shown of the fraction
$\dfrac{(x + b)^{m}}{(x + a)^{m+1} - x^{m+1}}$, which may be immediately reduced
to the form $\dfrac{x^{m} + B}{(m + 1)ax^{m} + A}$, where $x$~may be taken
so great that $x^{m}$~shall contain $A$~and~$B$ any number of
times.

We will now consider the sums of $x$~terms of the
following series, each of which may evidently be made
as great as we please, by taking a sufficient number
of its terms,
\begin{alignat*}{7}
&1    &&+ 2      &&+ 3     &&+ 4      &&+ \dots &&+ \; x - 1    &&+ x\Add{,}
\tag*{(1)} \\
&1^{2} &&+ 2^{2}  &&+ 3^{2}  &&+ 4^{2}  &&+ \dots &&+ (x - 1)^{2} &&+ x^{2}\Add{,}
\tag*{(2)} \displaybreak[0]\\
&1^{3} &&+ 2^{3}  &&+ 3^{3}  &&+ 4^{3}  &&+ \dots &&+ (x - 1)^{3} &&+ x^{3}\Add{,}
\tag*{(3)} \\
\DotRow{14} \displaybreak[0]\\
&1^{m} &&+ 2^{m}  &&+ 3^{m}  &&+ 4^{m}  &&+ \dots &&+ (x - 1)^{m} &&+ x^{m}\Add{.}
\tag*{($m$)}
\end{alignat*}
We propose to inquire what is the limiting ratio of
any one of these series to the last term of the succeeding
one; that is, to what do the ratios of $(1 + 2 + \dots + x)$
to~$x^{2}$, of $(1^{2} + 2^{2}\Add{+} \dots + x^{2})$ to~$x^{3}$, etc.,
approach, when $x$~is increased without limit.

To give an idea of the method of increase of these
series, we shall first show that $x$~may be taken so
great, that the last term of each series shall be as
small a part as we please of the sum of all those which
precede. To simplify the symbols, let us take the
third series $1^{3} + 2^{3} + \dots + x^{3}$, in which we are to
show that $x^{3}$~may be made less than any given part\Typo{.}{,}
\PageSep{69}
say one thousandth, of the sum of those which precede,
or of $1^{3} + 2^{3} \Add{+} \dots + (x - 1)^{3}$.

First, $x$~may be taken so great that $x^{3}$ and $(x - 1000)^{3}$
shall have a ratio as near to equality as we
please. For the ratio of these quantities being the
same as that of $1$ to $\left(1 - \dfrac{1000}{x}\right)^{3}$, and $\dfrac{1000}{x}$ being as
small as we please if $x$ may be as great as we please, it
follows that $1 - \dfrac{1000}{x}$, and, consequently, $\left(1 - \dfrac{1000}{x}\right)^{3}$
may be made as near to unity as we please, or the
ratio of $1$ to $\left(1 - \dfrac{1000}{x}\right)^{3}$, may be brought as near as
we please to that of $1$ to~$1$, or a ratio of equality. But
this ratio is that of $x^{3}$ to~$(x - 1000)^{3}$. Similarly the
ratios of $x^{3}$ to~$(x - 999)^{3}$, of $x^{3}$ to~$(x - 998)^{3}$, etc., up
to the ratio of $x^{3}$ to~$(x - 1)^{3}$ may be made as near as
we please to ratios of equality; there being one thousand
in all. If, then, $(x - 1)^{3} = \alpha x^{3}$, $(x - 2)^{3} = \beta x^{3}$,
etc., up to $(x - 1000)^{3} = \omega x^{3}$, $x$~can be taken so great
that each of the fractions $\alpha$,~$\beta$,~etc., shall be as near
to unity, or $\alpha + \beta + \dots + \omega$ as near\footnote
  {Observe that this conclusion depends upon the \emph{number} of quantities $\alpha$,~$\beta$,~etc.,
  being \emph{determinate}. If there be \emph{ten} quantities, each of which can be
  brought as near to unity as we please, their sum can be brought as near to~$10$
  as we please; for, take any fraction~$A$, and make each of those quantities
  differ from unity by less than the tenth part of~$A$, then will the sum differ
  from~$10$ by less than~$A$. This argument fails, if the number of quantities be
  unlimited.}
to~$1000$ as we
please. Hence
%[** TN: In-line in the original]
\[
\frac{1}{\alpha + \beta + \dots + \omega}\Add{,}
\]
which is
\[
\frac{x^{3}}{\alpha x^{3} + \beta x^{3} + \dots + \omega x^{3}},
\]
or
\[
\frac{x^{3}}{(x - 1)^{3} + (x - 2)^{2} + \dots + (x - 1000)^{3}},
\]
\PageSep{70}
can be brought as near to~$\dfrac{1}{1000}$ as we please; and by
the same reasoning, the fraction
\[
\frac{x^{3}}{(x - 1)^{3} + \dots + (x - 1001)^{3}}
\]
may be brought as near to~$\dfrac{1}{1001}$ as we please; that is,
may be made less than~$\dfrac{1}{1000}$. Still more then may
\[
\frac{x^{3}}{(x - 1)^{3} + \dots + (x - 1001)^{3} + \dots + 2^{3} + 1^{3}}
\]
be made less than~$\dfrac{1}{1000}$, or $x^{3}$~may be less than the
thousandth part of the sum of all the preceding terms.

In the same way it may be shown that a term may
be taken in any one of the series, which shall be less
than any given part of the sum of all the preceding
terms. It is also true that the difference of any two
succeeding terms may be made as small a part of
either as we please. For $(x + 1)^{m} - x^{m}$, when developed,
will only contain exponents less than~$m$, being
$mx^{m-1} + m\dfrac{m - 1}{2}\, x^{m-2} + \etc.$; and we have shown
(\PageRef{66}) that the sum of such a series may be made
less than any given part of~$x^{m}$. It is also evident
that, whatever number of terms we may sum, if a
sufficient number of succeeding terms be taken, the
sum of the latter shall exceed that of the former in
any ratio we please.

Let there be a series of fractions
\[
\frac{a}{pa + b},\quad
\frac{a'}{pa' + b'},\quad
\frac{a''}{pa'' + b''},\quad \etc.,
\]
in which $a$,~$a'$,~etc., $b$,~$b'$,~etc., increase without limit;
but in which the ratio of $b$~to~$a$, $b'$~to~$a'$, etc., diminishes
without limit. If it be allowable to begin by
\PageSep{71}
supposing $b$~as small as we please with respect to~$a$,
or $\dfrac{b}{a}$~as small as we please, the first, and all the succeeding
fractions, will be as near as we please to~$\dfrac{1}{p}$,
which is evident from the equations
\[
\frac{a}{pa + b} = \frac{1}{p + \dfrac{b}{a}},\quad
\frac{a'}{pa' + b'} = \frac{1}{p + \dfrac{b'}{a'}},\quad \etc.
\]
Form a new fraction by summing the numerators and
denominators of the preceding, such as
\[
\frac{a + a' + a'' + \etc.}
     {p(a + a' + a'' + \etc.) + b + b' + b'' + \etc.\Typo{,}{}},
\]
the \emph{etc.}\ extending to any given number of terms.

This may also be brought as near to~$\dfrac{1}{p}$ as we please.
For this fraction is the same as
\[
\text{$1$~divided by } p + \frac{b + b' + \etc.}{a + a' + \etc.};
\]
and it can be shown\footnote
  {See \Title{Study of Mathematics} (Reprint Edition, Chicago: The Open Court
  Publishing Co.), page~270.}
that
\[
\frac{b + b' + \etc.}{a + a' + \etc.}
\]
must lie between the least and greatest of the fractions
$\dfrac{b}{a}$,~$\dfrac{b'}{a'}$,~etc.
If, then, each of these latter fractions
can be made as small as we please, so also can
\[
\frac{b + b' + \etc.}{a + a' + \etc.}.
\]
No difference will be made in this result, if we use
the following fraction,
\[
\frac{A + (a + a' + a'' + \etc.)}
     {B + p(a + a' + a'' + \etc.) + b + b' + b'' + \etc.}\Add{,}
\Tag{(1)}
\]
\PageSep{72}
$A$~and~$B$ being given quantities; provided that we
can take a number of the original fractions sufficient
to make $a + a' + a'' + \etc.$, as great as we please,
compared with $A$~and~$B$. This will appear on dividing
the numerator and denominator of~\Eq{(1)} by $a + a' + a'' + \etc$.

Let the fractions be
\begin{gather*}
\frac{(x + 1)^{3}}{(x + 1)^{4} - x^{4}},\quad
\frac{(x + 2)^{3}}{(x + 2)^{4} - (x + 1)^{4}}, \\
\frac{(x + 3)^{3}}{(x + 3)^{4} - (x + 2)^{4}},\quad \etc.
\end{gather*}
The first of which, or $\dfrac{(x + 1)^{3}}{4x^{3} + \etc.}$ may, as we have
shown, be within any given difference of~$\dfrac{1}{4}$, and the
others still nearer, by taking a value of~$x$ sufficiently
great. Let us suppose each of these fractions to be
within $\dfrac{1}{100000}$ of~$\dfrac{1}{4}$. The fraction formed by summing
the numerators and denominators of these fractions
($n$~in number) will be within the same degree of
nearness to~$\frac{1}{4}$. But this is
\[
\frac{(x + 1)^{3} + (x + 2)^{3} + \dots + (x + n)^{3}}{(x + 1)^{4} - x^{4}}\Add{,}
\Tag{(2)}
\]
all the terms of the denominator disappearing, except
two from the first and last. If, then, we add~$x^{4}$ to
the denominator, and $1^{3} + 2^{3} + 3^{3} \Add{+} \dots + x^{3}$ to the numerator,
we can still take $n$ so great that $(x + 1)^{3} + \dots + (x + n)^{3}$
shall contain $1^{3} + \dots + x^{3}$ as often
as we please, and that $(x + n)^{4} - x^{4}$ shall contain~$x^{4}$
in the same manner. To prove the latter, observe
that the ratio of $(x + n)^{4} - x^{4}$ to~$x^{4}$ being $\left(1 + \dfrac{n}{x}\right)^{4}$,
can be made as great as we please, if it be permitted
\PageSep{73}
to take for~$n$ a number containing~$x$ as often as we
please. Hence, by the preceding reasoning, the fraction,
with its numerator and denominator thus increased,
or
\[
\frac{1^{3} + 2^{3} + 3^{3} + \dots + x^{3} + (x + 1)^{3} + \dots + (x + n)^{3}}
     {(x + n)^{4}}
\Tag{(3)}
\]
may be brought to lie within the same degree of nearness
to~$\frac{1}{4}$ as~\Eq{(2)}; and since this degree of nearness
could be named at pleasure, it follows that \Eq{(3)}~can
be brought as near to~$\frac{1}{4}$ as we please. Hence the
limit of the ratio of $(1^{3} + 2^{3} + \dots + x^{3})$ to~$x^{4}$, as $x$~is
increased without limit, is~$\frac{1}{4}$; and, in a similar manner,
it may be proved that the limit of the ratio of
$(1^{m} + 2^{m} + \dots + x^{m})$ to~$x^{m+1}$ is the same as that of
$\dfrac{(x + 1)^{m}}{(x + 1)^{m+1} - x^{m+1}}$ or $\dfrac{1}{m + 1}$.

This result will be of use when we come to the
first principles of the integral calculus. It may also
\index{Integral Calculus}%
be noticed that the limits of the ratios which $x\, \dfrac{x - 1}{2}$,
$x\, \dfrac{x - 1}{2}\, \dfrac{x - 2}{3}$, etc., bear to $x^{2}$,~$x^{3}$, etc., are severally $\dfrac{1}{2}$,
$\dfrac{1}{2·3}$, etc.; the limit being that to which the ratios approximate
as $x$~increases without limit. For $x\, \dfrac{x - 1}{2} ÷ x^{2} = \dfrac{x - 1}{2x}$,
$x\, \dfrac{x - 1}{2}\, \dfrac{x - 2}{3} ÷ x^{3} = \dfrac{x - 1}{2x}\, \dfrac{x - 2}{3x}$, etc.,
and the limits of $\dfrac{x - 1}{2}$, $\dfrac{x - 2}{3}$, are severally equal to
unity.

We now resume the elementary principles of the
Differential Calculus.
\PageSep{74}


\Subsection[Recapitulation of Results Reached in the Theory of Functions.]
{Recapitulation of Results.}

The following is a recapitulation of the principal
results which have hitherto been noticed in the general
theory of functions:
\index{Functions!recapitulation of results in the theory of}%

(1) That if in the equation $y = \phi(x)$, the variable~$x$
receives an increment~$dx$, $y$~is increased by the series
\[
\phi' x\, dx + \phi'' x\, \frac{(dx)^{2}}{2}
             + \phi''' x\, \frac{(dx)^{3}}{2·3} + \etc.
\]

(2) That $\phi'' x$ is derived in the same manner from~$\phi' x$,
that $\phi' x$~is from~$\phi x$; viz., that in like manner as
$\phi' x$~is the coefficient of~$dx$ in the development of
$\phi(x + dx)$, so $\phi'' x$~is the coefficient of~$dx$ in the development
of $\phi'(x + dx)$; similarly $\phi''' x$~is the coefficient
of~$dx$ in the development of~$\phi''(x + dx)$, and
so on.

(3) That $\phi' x$ is the limit of~$\dfrac{dy}{dx}$, or the quantity to
which the latter will approach, and to which it may
be brought as near as we please, when $dx$~is diminished.
It is called the differential coefficient of~$y$.

(4) That in every case which occurs in practice,
$dx$~may be taken so small, that any term of the series
above written may be made to contain the aggregate
of those which follow, as often as we please; whence,
though $\phi' x\, dx$~is not the actual increment produced
by changing~$x$ into~$x + dx$ in the function~$\phi x$, yet, by
taking $dx$ sufficiently small, it may be brought as near
as we please to a ratio of equality with the actual increment.


\Subsection[Approximations by the Differential Calculus.]
{Approximations.}

The last of the above-mentioned principles is of
the greatest utility, since, by means of it, $\phi' x\, dx$~may
\PageSep{75}
\index{Errors, in the valuation of quantities}%
be made as nearly as we please the actual increment;
and it will generally happen in practice, that $\phi' x\, dx$
may be used for the increment of~$\phi x$ without sensible
error; that is, if in~$\phi x$, $x$~be changed into $x + dx$, $dx$~being
very small, $\phi x$~is changed into $\phi x + \phi' x\, dx$,
very nearly. Suppose that $x$ being the correct value
of the variable, $x + h$ and $x + k$ have been successively
substituted for it, or the errors $h$~and~$k$ have
been committed in the valuation of~$x$, $h$~and~$k$ being
very small. Hence $\phi(x + h)$ and $\phi(x + k)$ will be
erroneously used for~$\phi x$. But these are nearly $\phi x + \phi' x\, h$
and $\phi x + \phi' x\, k$, and the errors committed in
taking~$\phi x$ are $\phi' x\, h$ and $\phi' x\, k$, very nearly. These
last are in the proportion of $h$ to~$k$, and hence results
a proposition of the utmost importance in every practical
application of mathematics, viz., that if two different,
but small, errors be committed in the valuation
of any quantity, the errors arising therefrom at
the end of any process, in which both the supposed
values of~$x$ are successively adopted, are very nearly
in the proportion of the errors committed at the beginning.
For example, let there be a right-angled
triangle, whose base is~$3$, and whose other side should
be~$4$, so that the hypothenuse should be $\sqrt{3^{2} + 4^{2}}$
or~$5$. But suppose that the other side has been twice
erroneously measured, the first measurement giving
$4.001$, and the second $4.002$, the errors being $.001$
and~$.002$. The two values of the hypothenuse thus
obtained are
\[
\sqrt{3^{2} + 4.001^{2}}, \quad\text{or}\quad \sqrt{25.008001},
\]
and
\[
\sqrt{3^{2} + 4.002^{2}}, \quad\text{or}\quad \sqrt{25.016004},
\]
which are very nearly $5.0008$ and $5.0016$. The errors
of the hypothenuse are then $.0008$ and $.0016$ nearly;
and these last are in the proportion of $.001$ and~$.002$.
\PageSep{76}

It also follows, that if $x$~increase by successive equal
steps, any function of~$x$ will, for a few steps, increase
so nearly in the same manner, that the supposition of
such an increase will not be materially wrong. For,
if $h$,~$2h$,~$3h$, etc., be successive small increments given
to~$x$, the successive increments of~$\phi x$ will be $\phi' x\, h$,
$\phi' x\, 2h$, $\phi' x\, 3h$,~etc.\ nearly; which being proportional
to $h$,~$2h$,~$3h$, etc., the increase of the function is nearly
doubled, trebled, etc., if the increase of~$x$ be doubled,
trebled,~etc.

This result may be rendered conspicuous by reference
to any astronomical ephemeris, in which the
\index{Astronomical ephemeris}%
positions of a heavenly body are given from day to
day. The intervals of time at which the positions are
given differ by $24$~hours, or nearly $\frac{1}{365}$\th~part of the
whole year. And even for this interval, though it can
hardly be called \emph{small} in an astronomical point of view,
the increments or decrements will be found so nearly
the same for four or five days together, as to enable
the student to form an idea how much more near they
would be to equality, if the interval had been less, say
one hour instead of twenty-four. For example, the
sun's longitude on the following days at noon is written
\index{Sun's longitude}%
underneath, with the increments from day to day.
\[
\ArrayCompress
\begin{array}{c*{2}{>{\ }c}c}
\ColHead[September]{1834 \\ September} &
\ColHead[Sun's longitude]{Sun's longitude \\ at noon.} &
\ColHead{Increments.} &
\ColHead[Proportion which the differences]{Proportion which the differences \\
  of the increments bear to the \\
  whole increments.} \\
%
1\text{st} & 158\rlap{$°$}\ \ 30\rlap{$'$}\ \ 35\rlap{$''$}
                          & \Low{58\rlap{$'$}\ \Z9\rlap{$''$}} & \\
2\text{nd} & 159\ \ 28\ \  44 & \Low{58\ 12} & \frac{3}{3489} \\
3\text{rd} & 160\ \ 26\ \  56 & \Low{58\ 13} & \frac{1}{3492} \\
4\text{th} & 161\ \ 25\ \ \Z9 & \Low{58\ 14} & \frac{1}{3493} \\
5\text{th} & 162\ \ 23\ \  23
\end{array}
\]
The sun's longitude is a function of the time; that is,
the number of years and days from a given epoch
being given, and called~$x$, the sun's longitude can be
\PageSep{77}
found by an algebraical expression which may be
called~$\phi x$. If we date from the first of January,~1834,
$x$~is~$.666$, which is the decimal part of a year between
the first days of January and September. The increment
is one day, or nearly $.0027$~of a year. Here $x$~is
successively made equal to~$.666$, $.666 + 0027$, $.666 + 2 × .0027$,
etc.; and the intervals of the corresponding
values of~$\phi x$, if we consider only minutes, are the
same; but if we take in the seconds, they differ from
one another, though only by very small parts of themselves,
as the last column shows.


\Subsection[Solution of Equations by the Differential Calculus.]
{Solution of Equations.}

This property is also used\footnote
  {See \Title{Study of Mathematics} (Reprint Edition, Chicago: The Open Court
  Publishing Co., 1898), page~169 et~seq.}
in finding logarithms
\index{Equations, solution of}%
intermediate to those given in the tables; and may
be applied to find a nearer solution to an equation,
than one already found. For example, suppose it required
to find the value of~$x$ in the equation $\phi x = 0$,
$a$~being a near approximation to the required value.
Let $a + h$ be the real value, in which $h$~will be a small
quantity. It follows that $\phi(a + h) = 0$, or, which is
nearly true, $\phi a + \phi' a\, h = 0$. Hence the real value of~$h$
is nearly~$-\dfrac{\phi a}{\phi' a}$, or the value $a - \dfrac{\phi a}{\phi' a}$ is a nearer
approximation to the value of~$x$. For example, let
$x^{2} + x - 4 = 0$ be the equation. Here $\phi x = x^{2} + x - 4$,
and $\phi(x + h) = (x + h)^{2} + x + h - 4 = x^{2} + x - 4 + (2x + 1)h + h^{2}$;
so that $\phi' x = 2x + 1$. A near value
of~$x$ is~$1.57$; let this be~$a$. Then $\phi a = .0349$, and
$\phi' a = 4.14$. Hence $-\dfrac{\phi a}{\phi' a} = -.00843$. Hence
$1.57 - .00843$, or~$1.56157$, is a nearer value of~$x$. If
\PageSep{78}
we proceed in the same way with~$1.5616$, we shall
find a still nearer value of~$x$, viz., $1.561553$. We
have here chosen an equation of the second degree,
in order that the student may be able to verify the
result in the common way; it is, however, obvious
that the same method may be applied to equations
of higher degrees, and even to those which are not
to be treated by common algebraical method, such as
$\tan x = ax$.


\Subsection{Partial and Total Differentials.}

We have already observed, that in a function of
\index{Differentials!partial|EtSeq}%
\index{Differentials!total|EtSeq}%
\index{Partial!differentials|EtSeq}%
\index{Total!differentials|EtSeq}%
more quantities than one, those only are mentioned
which are considered as variable; so that all which
we have said upon functions of one variable, applies
\index{Functions!of several variables|EtSeq}%
equally to functions of several variables, so far as a
\index{Variables!functions of several|EtSeq}%
change in one only is concerned. Take for example
$x^{2} y + 2xy^{3}$. If $x$~be changed into $x + dx$, $y$~remaining
the same, this function is increased by $2xy\, dx + 2y^{3}\, dx + \etc.$,
in which, as in \PageRef{29}, no terms are contained
in the~\emph{etc.}\ except those which, by diminishing~$dx$,
can be made to bear as small a proportion as we
please to the first terms. Again, if $y$~be changed into
$y + dy$, $x$~remaining the same, the function receives
the increment $x^{2}\, dy + 6xy^{2}\, dy + \etc.$; and if $x$~be changed
into $x + dx$, $y$~being at the same time changed into
$y + dy$, the increment of the function is $(2xy + 2y^{3})\, dx + (x^{2} + 6xy^{2})\, dy + \etc$.
If, then, $u = x^{2} y + 2xy^{3}$, and
$du$~denote the increment of~$u$, we have the three following
equations, answering to the various suppositions
above mentioned, \\
(1) when $x$~only varies,
\[
du = (2xy + 2y^{3})\, dx + \etc.
\]
\PageSep{79}
(2) when $y$~only varies,
\[
du = (x^{2} + 6xy^{2})\, dy + \etc.
\]
(3) when both $x$~and~$y$ vary,
\[
du = (2xy + 2y^{3})\, dx + (x^{2} + 6xy^{2})\, dy + \etc.
\]
in which, however, it must be remembered, that $du$~does
not stand for the same thing in any two of the
three equations: it is true that it always represents
an increment of~$u$, but as far as we have yet gone, we
have used it indifferently, whether the increment of~$u$
was the result of a change in $x$~only, or $y$~only, or both
together.

To distinguish the different increments of~$u$, we
must therefore seek an additional notation, which,
\index{Calculus, notation of|EtSeq}%
\index{Notation!of the Differential Calculus|EtSeq}%
without sacrificing the~$du$ that serves to remind us
that it was $u$ which received an increment, may also
point out from what supposition the increment arose.
For this purpose we might use $d_{x}u$~and~$d_{y}u$, and $d_{x,y}u$,
to distinguish the three; and this will appear to the
learner more simple than the one in common use,
which we shall proceed to explain. We must, however,
remind the student, that though in matters of
reasoning, he has a right to expect a solution of every
difficulty, in all that relates to notation, he must trust
entirely to his instructor; since he cannot judge between
the convenience or inconvenience of two symbols
without a degree of experience which he evidently
cannot have had. Instead of the notation above
described, the increments arising from a change in $x$
and~$y$ are severally denoted by $\dfrac{du}{dx}\, dx$ and $\dfrac{du}{dy}\, dy$, on
the following principle: If there be a number of results
obtained by the same species of process, but on
different suppositions with regard to the quantities
\PageSep{80}
used; if, for example, $p$~be derived from some supposition
with regard to~$a$, in the same manner as are $q$
and~$r$ with regard to $b$~and~$c$, and if it be inconvenient
and unsymmetrical to use separate letters $p$,~$q$, and~$r$,
for the three results, they may be distinguished by
using the same letter~$p$ for all, and writing the three
results thus, $\dfrac{p}{a}\, a$, $\dfrac{p}{b}\, b$, $\dfrac{p}{c}\, c$. Each of these, in common
algebra, is equal to~$p$, but the letter~$p$ does not
stand for the same thing in the three expressions.
The first is the~$p$, so to speak, which belongs to~$a$, the
second that which belongs to~$b$, the third that which
belongs to~$c$. Therefore the numerator of each of the
fractions $\dfrac{p}{a}$,~$\dfrac{p}{b}$, and~$\dfrac{p}{c}$, must never be separated
from its denominator, because the value of the former
depends, in part, upon the latter; and one~$p$ cannot
be distinguished from another without its denominator.
The numerator by itself only indicates what operation
is to be performed, and on what quantity; the
denominator shows what quantity is to be made use
of in performing it. Neither are we allowed to say
that $\dfrac{p}{a}$ divided by~$\dfrac{p}{b}$ is~$\dfrac{b}{a}$; for this supposes that $p$~means
the same thing in both quantities.

In the expressions $\dfrac{du}{dx}\, dx$, and $\dfrac{du}{dy}\, dy$, each denotes
that $u$~has received an increment; but the first points
out that~$x$, and the second that~$y$, was supposed to increase,
in order to produce that increment; while $du$~by
itself, or sometimes $d.u$, is employed to express
the increment derived from both suppositions at once.
And since, as we have already remarked, it is not the
ratios of the increments themselves, but the limits of
those ratios, which are the objects of investigation in
\PageSep{81}
the Differential Calculus, here, as in \PageRef{28}, $\dfrac{du}{dx}\, dx$,
and $\dfrac{du}{dy}\, dy$, are generally considered as representing
those terms which are of use in obtaining the limiting
ratios, and do not include those terms, which, from
\index{Limiting ratios}%
\index{Ratios, limiting}%
their containing higher powers of $dx$~or~$dy$ than the
first, may be made as small as we please with respect
to $dx$~or~$dy$. Hence in the example just given, where
$u = x^{2} y + 2xy^{3}$, we have
\begin{align*}
&\dfrac{du}{dx}\, dx = (2xy + 2y^{3})\, dx,
  &&\text{or}\quad \frac{du}{dx} = 2xy + 2y^{3}\Add{,} \\
&\dfrac{du}{dy}\, dy = (x^{2} + 6xy^{2})\, dy,
  &&\text{or}\quad \frac{du}{dy} = x^{2} + 6xy^{2}\Add{,} \\
&du \quad\text{or}\quad d.u = \frac{du}{dx}\, dx + \frac{du}{dy}\,dy.
\end{align*}

The last equation gives a striking illustration of
the method of notation. Treated according to the
common rules of algebra, it is $du = du + du$, which is
absurd, but which appears rational when we recollect
that the second~$du$ arises from a change in $x$~only, the
third from a change in $y$~only, and the first from a
change in both. The same equation may be proved
to be generally true for all functions of $x$~and~$y$, if we
bear in mind that no term is retained, or need be retained,
as far as the limit is concerned, which, when
$dx$~or~$dy$ is diminished, diminishes without limit as
compared with them. In using $\dfrac{du}{dx}$ and $\dfrac{du}{dy}$ as differential
coefficients of~$u$ with respect to $x$~and~$y$, the objection
(\PageRef{27}) against considering these as the
limits of the ratios, and not the ratios themselves,
does not hold, since the numerator is not to be separated
from its denominator.
\PageSep{82}

Let $u$ be a function of $x$~and~$y$, represented\footnote
  {The symbol $\phi(x, y)$ must not be confounded with~$\phi(xy)$. The former represents
  any function of $x$~and~$y$; the latter a function in which $x$~and~$y$ only
  enter so far as they are contained in their product. The second is therefore
  a particular case of the first; but the first is not necessarily represented by
  the second. For example, take the function $xy + \sin xy$, which, though it
  contains both $x$~and~$y$, yet can only be altered by such a change in $x$~and~$y$ as
  will alter their product, and if the product be called~$p$, will be $p + \sin p$. This
  may properly be represented by~$\phi(xy)$; whereas $x + xy^{2}$ cannot be represented
  in the same way, since other functions besides the product are contained
  in it.}
by~$\phi(x, y)$.
It is indifferent whether $x$~and~$y$ be changed
\index{Coefficients, differential}%
\index{Differential coefficients}%
at once into $x + dx$ and $y + dy$, or whether $x$~be first
changed into $x + dx$, and $y$~be changed into $y + dy$ in
the result. Thus, $x^{2} y + y^{3}$ will become $(x + dx)^{2}(y + dy) + (y + dy)^{3}$
in either case. If $x$~be changed
into $x + dx$, $u$~becomes $u + \ux\, dx + \etc.$, (where $\ux$~is
what we have called the differential coefficient of~$u$
with respect to~$x$, and is itself a function of $x$~and~$y$;
and the corresponding increment of~$u$ is $\ux\, dx + \etc.$)\Add{.}
If in this result $y$~be changed into $y + dy$, $u$~will assume
the form $u + \uy\, dy + \etc.$, where $\uy$~is the differential
coefficient of~$u$ with respect to~$y$; and the increment
which $u$~receives will be $\uy + \etc$. Again,
when $y$~is changed into $y + dy$, $\ux$,~which is a function
of $x$~and~$y$, will assume the form $\ux + p\, dy + \etc.$; and
$u + \ux\, dx + \etc.$\ becomes $u + \uy\, dy + \etc. + (\ux + p\, dy + \etc.)\, dx + \etc.$,
or $u + \uy\, dy + \ux\, dx + p\, dx\, dy + \etc.$,
in which the term $p\, dx\, dy$ is useless in finding the limit.
For since $dy$~can be made as small as we please,
$p\, dx\, dy$ can be made as small a part of~$p\, dx$ as we please,
and therefore can be made as small a part of~$dx$ as
we please. Hence on the three suppositions already
made, we have the following results:
%[** TN: Re-formatted]
\begin{itemize}
\item[(1)] when $x$~only is changed into~$x + dx$,
$u$~receives the increment
$\ux\, dx + \etc$.\Add{,}

\item[(2)] when $y$~only is changed into~$y + dy$,
$u$~receives the increment
$\uy\, dy + \etc$.\Add{,}

\item[(3)] when $x$~becomes $x + dx$ and $y$~becomes $y + dy$ at once,
$u$~receives the increment
$\ux\, dx + \uy\, dy + \etc$.\Add{,}
\end{itemize}
\PageSep{83}
the \emph{etc.}\ in each case containing those terms only which
can be made as small as we please, with respect to
the preceding terms. In the language of Leibnitz,
\index{Leibnitz}%
we should say that if $x$~and~$y$ receive infinitely small
\index{Infinitely small, the notion of}%
increments, the sum of the infinitely small increments
of~$u$ obtained by making these changes separately, is
equal to the infinitely small increment obtained by
making them both at once. As before, we may correct
this inaccurate method of speaking. The several
increments in (1),~(2), and~(3), maybe expressed by
$\ux\, dx + P$, $\uy\, dy + Q$, and $\ux\, dx + \uy\, dy + R$; where $P$,~$Q$,
and~$R$ can be made such parts of $dx$~or~$dy$ as we
please, by taking $dx$~or~$dy$ sufficiently small. The sum
of the two first is $\ux\, dx + \uy\, dy + P + Q$, which differs
from the third by $P + Q - R$; which, since each of
its terms can be made as small a part of $dx$~or~$dy$ as
we please, can itself be made less than any given part
of $dx$~or~$dy$.

This theorem is not confined to functions of two
variables only, but may be extended to those of any
number whatever. Thus, if $z$~be a function of $p$,~$q$,~$r$,
and~$s$, we have
\[
d.z \quad\text{or}\quad
dz = \frac{dz}{dp}\, dp
   + \frac{dz}{dq}\, dq
   + \frac{dz}{dr}\, dr
   + \frac{dz}{ds}\, ds + \etc.
\]
in which $\dfrac{dz}{dp}\, dp + \etc.$\ is the increment which a change
in $p$~\emph{only} gives to~$z$, and so on. The \emph{etc.}\ is the representative
of an infinite series of terms, the aggregate
of which diminishes continually with respect to $dp$,~$dq$,~etc.,
as the latter are diminished, and which, therefore,
\PageSep{84}
has no effect on the \emph{limit} of the ratio of~$d.z$ to
any other quantity.


\Subsection[Application of the Theorem for Total Differentials to the Determination of Total Resultant Errors.]
{Practical Application of the Preceding Theorem.}

We proceed to an important practical use of this
\index{Errors, in the valuation of quantities}%
theorem. If the increments $dp$,~$dq$,~etc., be small,
this last-mentioned equation, (the terms included in
the \emph{etc.}\ being omitted,) though not actually true, is
sufficiently near the truth for all practical purposes;
which renders the proposition, from its simplicity, of
the highest use in the applications of mathematics.
For if any result be obtained from a set of \textit{data}, no
one of which is exactly correct, the error in the result
would be a very complicated function of the errors in
the \textit{data}, if the latter were considerable. When they
are small, the error in the results is very nearly the
sum of the errors which would arise from the error in
each \textit{datum}, if all the others were correct. For if $p$,~$q$,~$r$,
and~$s$, are the \emph{presumed} values of the \textit{data}, which
give a certain value~$z$ to the function required to be
found; and if $p + dp$, $q + dq$, etc., be the \emph{correct} values
of the \textit{data}, the correction of the function~$z$ will be
very nearly made, if $z$~be increased by $\dfrac{dz}{dp}\, dp + \dfrac{dz}{dq}\, dq + \dfrac{dz}{dr}\, dr + \dfrac{dz}{ds}\, ds$,
being the sum of terms which would
arise from each separate error, if each were made in
turn by itself.

For example: A transit instrument is a telescope
\index{Transit instrument}%
mounted on an axis, so as to move in the plane of the
meridian only, that is, the line joining the centres of
the two glasses ought, if the telescope be moved, to
pass successively through the zenith and the pole.
Hence can be determined the exact time, as shown by
a clock, at which any star passes a vertical thread,
\PageSep{85}
fixed inside the telescope so as apparently to cut the
field of view exactly in half, which thread will always
cover a part of the meridian, if the telescope be correctly
adjusted. In trying to do this, three errors
may, and generally will be committed, in some small
degree. (1)~The axis of the telescope may not be exactly
level; (2)~the ends of the same axis may not be
exactly east and west; (3)~the line which joins the
centres of the two glasses, instead of being perpendicular
to the axis of the telescope, may be inclined
to it. If each of these errors were considerable, and
the time at which a star passed the thread were observed,
the calculation of the time at which the same
star passes the real meridian would require complicated
formulæ, and be a work of much labor. But if
the errors exist in small quantities only, the calculation
is very much simplified by the preceding principle.
For, suppose only the first error to exist, and
calculate the corresponding error in the time of passing
the thread. Next suppose only the second error,
and then only the third to exist, and calculate the
effect of each separately, all which may be done by
simple formulæ. The effect of all the errors will then
be the sum of the effects of each separate error, at
least with sufficient accuracy for practical purposes.
The formulæ employed, like the equations in \PageRef{28},
are not actually true in any case, but approach more
near to the truth as the errors are diminished.


\Subsection{Rules for Differentiation.}

In order to give the student an opportunity of exercising
\index{Differentiation!of the common functions}%
himself in the principles laid down, we will
so far anticipate the treatises on the Differential Calculus
as to give the results of all the common rules
\PageSep{86}
for differentiation; that is, assuming $y$~to stand for
various functions of~$x$, we find the increment of~$y$ arising
from an increment in the value of~$x$, or rather,
that term of the increment which contains the first
power of~$dx$. This term, in theory, is the only one
on which the \emph{limit} of the ratio of the increments depends;
in practice, it is sufficiently near to the real
increment of~$y$, if the increment of~$x$ be small.

{\Loosen (1) $y = x^{m}$, where $m$~is either whole or fractional,
\index{Differentiation!of the common functions}%
positive or negative; then $dy = mx^{m-1}\, dx$. Thus the
increment of~$x^{\efrac{2}{3}}$ or the first term of $(x + dx)^{\efrac{2}{3}} - x^{\efrac{2}{3}}$
is $\frac{2}{3}x^{\efrac{2}{3}-1}\, dx$, or~$\dfrac{2\, dx}{3x^{\efrac{1}{3}}}$. Again, if $y = x^{8}$, $dy = 8x^{7}\, dx$.
When the exponent is negative, or when $y = \dfrac{1}{x^{m}}$,
$dy = -\dfrac{m\, dx}{x^{m+1}}$, or when $y = x^{-m}$, $dy = -mx^{-m-1}\, dx$,
which is according to the rule. The negative sign
indicates that an increase in~$x$ decreases the value
of~$y$; which, in this case, is evident.}

(2) $y = a^{x}$. Here $dy = a^{x}\log a\, dx$ where the logarithm
(as is always the case in analysis, except
where the contrary is specially mentioned) is the Naperian
or hyperbolic logarithm. When $a$~is the base
of these logarithms, that is when $a = 2.7182818 = e$,
\index{Logarithms}%
or when $y = e^{x}$, $dy = e^{x}\, dx$.

(3) $y = \log x$ (the Naperian logarithm). Here
$dy = \dfrac{dx}{x}$. If $y = \text{common log}~x$, $dy = -.4342944\, \dfrac{dx}{x}$.

(4) $y = \sin x$, $dy = \cos x\, dx$; $y = \cos x$, $dy = -\sin x\, dx$;
$y = \tan x$, $dy = \dfrac{dx}{\cos^{2} x}$.


\Subsection[Illustration of the Rules for Differentiation.]
{Illustration of the Preceding Formulæ.}

At the risk of being tedious to some readers, we
will proceed to illustrate these formulæ by examples
\PageSep{87}
from the tables of logarithms and sines, let $y = \text{common log}~x$.
\index{Logarithms}%
\index{Sines}%
If $x$~be changed into $x + dx$, the real increment
of~$y$ is
\[
.4342944 \left(\frac{dx}{x} - \tfrac{1}{2}\, \frac{(dx)^{2}}{x^{2}}
  + \tfrac{1}{3}\, \frac{(dx)^{3}}{x^{3}} - \etc.\right),
\]
in which the law of continuation is evident. The corresponding
series for Naperian logarithms is to be
found in \PageRef{20}. From the first term of this the
limit of the ratio of $dy$~to~$dx$ can be found; and if $dx$~be
\index{Ratio!of two increments}%
small, this will represent the increment with sufficient
accuracy. Let $x = 1000$, whence $y = \text{common log}~ 1000 =3$;
and let $dx = 1$, or let it be required to
find the common logarithm of $1000 + 1$, or~$1001$. The
first term of the series is therefore $.4342944 × \frac{1}{1000}$, or
$.0004343$, taking seven decimal places only. Hence
$\log 1001 = \log 1000 + .0004343$ or $3.0004343$ nearly.
The tables give $3.0004341$, differing from the former
only in the $7$\Chg{th}{\th}~place of decimals.

{\Loosen Again, let $y = \sin x$; from which, by \PageRef{20}, as
before, if $x$~be increased by~$dx$, $\sin x$~is increased by
$\cos x\, dx - \frac{1}{2}\sin x\, (dx)^{2} - \etc.$, of which we take only
the first term. Let $x = 16°$, in which case $\sin x = .2756374$,
and $\cos x = .9612617$.} Let $dx = 1'$, or, as
it is represented in analysis, where the angular unit is
that angle whose arc is equal to the radius,\footnote
  {See \Title{Study of Mathematics} (Chicago; The Open Court Pub. Co.), page~273
  et~seq.}
$\frac{60}{206265}$.
Hence $\sin 16°\, 1' = \sin 16° + .9612617 × \frac{60}{206265} =
.2756374 + .0002797 = .2759171$, nearly. The tables
give~$.2759170$. These examples may serve to show
how nearly the real ratio of two increments approaches
to their limit, when the increments themselves are
small.
\PageSep{88}


\Subsection{Differential Coefficients of Differential
Coefficients.}

When the differential coefficient of a function of~$x$
\index{Coefficients, differential}%
\index{Differential coefficients!of higher orders}%
\index{Finite differences|EtSeq}%
\index{Orders, differential coefficients of higher}%
\index{Successive differentiation|EtSeq}%
has been found, the result, being a function of~$x$, may
be also differentiated, which gives the differential coefficient
of the differential coefficient, or, as it is called,
the \emph{second} differential coefficient. Similarly the differential
coefficient of the second differential coefficient
is called the third differential coefficient, and so on.
We have already had occasion to notice these successive
differential coefficients in \PageRef{22}, where it appears
that $\phi' x$~being the first differential coefficient of~$\phi x$,
$\phi'' x$~is the coefficient of~$h$ in the development
$\phi'(x + h)$, and is therefore the differential coefficient
of~$\phi' x$, or what we have called the second differential
coefficient of~$\phi x$. Similarly $\phi''' x$~is the third differential
coefficient of~$\phi x$. If we were strictly to adhere
to our system of notation, we should denote the
several differential coefficients of~$\phi x$ or~$y$ by
\[
\frac{dy}{dx}\Add{,}\quad
\frac{d.\dfrac{dy}{dx}}{dx}\Add{,}\quad
\frac{d.\dfrac{d.\frac{dy}{dx}}{dx}}{dx}\Add{,}\quad \etc.
\]
In order to avoid so cumbrous a system of notation,
the following symbols are usually preferred,
\[
\frac{dy}{dx}\Add{,}\quad
\frac{d^{2} y}{dx^{2}}\Add{,}\quad
\frac{d^{3} y}{dx^{3}}\Add{,}\quad \etc.
\]


\Subsection{Calculus of Finite Differences. Successive
Differentiation.}

We proceed to explain the manner in which this
\index{Differentiation!successive|EtSeq}%
notation is connected with our previous ideas on the
subject.
\PageSep{89}

When in any function of~$x$, an increase is given to~$x$,
\index{Differences!calculus of}%
which is not supposed to be as small as we please,
it is usual to denote it by~$\Delta x$ instead of~$dx$, and the
corresponding increment of~$y$ or~$\phi x$, by~$\Delta y$ or~$\Delta\phi x$,
instead of~$dy$ or~$d\phi x$. The symbol~$\Delta x$ is called the
\emph{difference} of~$x$, being the difference between the value
of the variable~$x$, before and after its increase.

Let $x$ increase at successive steps by the same difference;
that is, let a variable, whose first value is~$x$,
successively become $x + \Delta x$, $x + 2\Delta x$, $x + 3\Delta x$, etc.,
and let the successive values of~$\phi x$ corresponding to
these values of~$x$ be $y$,~$y_{1}$, $y_{2}$,~$y_{3}$,~etc.; that is, $\phi x$~is
called~$y$, $\phi(x + \Delta x)$ is~$y_{1}$, $\phi(x + 2\Delta x)$ is~$y_{2}$, etc., and,
generally, $\phi(x + m\Delta x)$ is~$y_{m}$. Then, by our previous
definition $y_{1} - y$ is~$\Delta y$, $y_{2} - y_{1}$ is~$\Delta y_{1}$, $y_{3} - y_{2}$ is~$\Delta y_{2}$,
etc., the letter~$\Delta$ before a quantity always denoting
the increment it would receive if $x + \Delta x$ were substituted
for~$x$. Thus $y_{3}$ or $\phi(x + 3\Delta x)$ becomes $\phi(x + \Delta x + 3\Delta x)$,
or $\phi(x + 4\Delta x)$, when $x$~is changed into
$x + \Delta x$, and receives the increment $\phi(x + 4\Delta x) - \phi(x + 3\Delta x)$, or $y_{4} - y_{3}$. If $y$~be a function which decreases
when $x$~is increased, $y_{1} - y$, or $\Delta y$ is negative.

It must be observed, as in \PageRef{26}, that $\Delta x$~does
not depend upon~$x$, because $x$~occurs in it; the symbol
merely signifies an increment given to~$x$, which
increment is not necessarily dependent upon the value
of~$x$. For instance, in the present case we suppose
it a given quantity; that is, when $x + \Delta x$ is changed
into $x + \Delta x + \Delta x$, or $x + 2\Delta x$, $x$~is changed, and $\Delta x$~is
not.

In this way we get the two first of the columns underneath,
in which each term of the \emph{second} column is
formed by subtracting the term which immediately
precedes it in the first column from the one which immediately
\PageSep{90}
follows. Thus $\Delta y$ is $y_{1} - y$, $\Delta y_{1}$ is $y_{2} - y_{1}$,
etc.
\begin{gather*}%**** Tall, bad page break
\left.
\begin{alignedat}{2}
& \PadTo[l]{\phi(x + 4\Delta x)}{\phi(x)} && y \\
& \phi(x + \Z\Delta x)\qquad              && y_{1} \\
& \phi(x +  2\Delta x)                    && y_{2} \\
& \phi(x +  3\Delta x)                    && y_{3} \\
& \phi(x +  4\Delta x)                    && y_{4}
\end{alignedat}\
\right|
%
\left.
\begin{aligned}
& \Delta y \\
& \Delta y_{1} \\
& \Delta y_{2} \\
& \Delta y_{3}
\end{aligned}\
\right|
%
\left.
\begin{aligned}
& \Delta^{2} y \\
& \Delta^{2} y_{1} \\
& \Delta^{2} y_{2}
\end{aligned}\
\right|
%
\left.
\begin{aligned}
& \Delta^{3} y \\
& \Delta^{3} y_{1}
\end{aligned}\
\right|
%
\begin{aligned}
& \Delta^{4} y
\end{aligned} \\
\PadTo{\phi(x +  4\Delta x)}{\etc.}\phantom{\qquad\qquad\qquad\qquad\qquad\qquad}
\end{gather*}

In the first column is to be found a series of successive
values of the same function~$\phi x$, that is, it contains
terms produced by substituting successively in~$\phi x$
the quantities $x$, $x + \Delta x$, $x + 2\Delta x$, etc., instead of~$x$.
The second column contains the successive values
of another function $\phi(x + \Delta x) - \phi x$, or~$\Delta \phi x$, made by
the same substitutions; if, for example, we substitute
$x + 2\Delta x$ for~$x$, we obtain $\phi(x + 3\Delta x) - \phi(x + 2\Delta x)$,
or $y_{3} - y_{2}$, or~$\Delta y_{2}$. If, then, we form the successive
differences of the terms in the second column, we obtain
a new series, which we might call the differences
of the differences of the first column, but which are
called the \emph{second differences} of the first column. And
as we have denoted the operation which deduces the
second column from the first by~$\Delta$, so that which deduces
the third from the second may be denoted by~$\Delta\Delta$,
which is abbreviated into~$\Delta^{2}$. Hence as $y_{1} - y$
was written~$\Delta y$, $\Delta y_{1} - \Delta y$ is written~$\Delta\Delta y$, or~$\Delta^{2} y$. And
the student must recollect, that in like manner as $\Delta$~is
not the symbol of a number, but of an operation,
so $\Delta^{2}$~does not denote a number multiplied by itself,
but an operation repeated upon its own result; just
as the logarithm of the logarithm of~$x$ might be written
$\log^{2} x$; $(\log x)^{2}$~being reserved to signify the square
of the logarithm of~$x$. We do not enlarge on this notation,
as the subject is discussed in most treatises on
\PageSep{91}
algebra.\footnote
  {The reference of the original text is to ``the treatise on \Title{Algebraical Expressions},''
  Number~105 of the Library of Useful Knowledge,---the same series
  in which the present work appeared. The first six pages of this treatise are
  particularly recommended by De~Morgan in relation to the present point.---\Ed.}
Similarly the terms of the fourth column,
or the differences of the second differences, have the
prefix~$\Delta\Delta\Delta$ abbreviated into~$\Delta^{3}$, so that $\Delta^{2} y_{1} - \Delta^{2} y = \Delta^{3} y$, etc.

When we have occasion to examine the results
which arise from supposing $\Delta x$~to diminish without
limit, we use~$dx$ instead of~$\Delta x$, $dy$~instead of~$\Delta y$, $d^{2} y$~instead
of~$\Delta^{2} y$, and so on. If we suppose this case, we
can show that the ratio which the term in any column
bears to its corresponding term in any preceding column,
diminishes without limit. Take for example,
$d^{2} y$~and~$dy$. The latter is $\phi(x + dx) - \phi x$, which, as
we have often noticed already, is of the form $p\, dx + q\, (dx)^{2} + \etc.$,
in which $p$,~$q$,~etc., are also functions
of~$x$. To obtain~$d^{2} y$, we must, in this series, change~$x$
into $x + dx$, and subtract $p\, dx + q\, (dx)^{2} + \etc.$\ from
the result. But since $p$,~$q$,~etc., are functions of~$x$,
this change gives them the form
\[
p + p'\, dx + \etc.,\quad
q + q'\, dx + \etc.;
\]
so that $d^{2} y$~is
\begin{gather*}
(p + p'\, dx + \etc.)\, dx + (q + q'\, dx + \etc.)(dx)^{2} + \etc. \\
{} - (p\, dx + q\, (dx)^{2} + \etc.)
\end{gather*}
in which the first power of~$dx$ is destroyed. Hence
(\PageRefs{42}{44}), the ratio of $d^{2} y$ to~$dx$ diminishes without
limit, while that of $d^{2} y$ to~$(dx)^{2}$ has a finite limit,
except in those particular cases in which the second
power of~$dx$ is destroyed, in the previous subtraction,
as well as the first. In the same way it may be shown
that the ratio of $d^{3} y$ to $dx$~and $(dx)^{2}$ decreases without
limit, while that of $d^{3} y$ to~$(dx)^{3}$ remains finite; and so
\PageSep{92}
on. Hence we have a succession of ratios $\dfrac{dy}{dx}$, $\dfrac{d^{2} y}{dx^{2}}$, $\dfrac{d^{3} y}{dx^{3}}$,
etc., which tend towards finite limits when $dx$~is diminished.

{\Loosen We now proceed to show that in the development
of $\phi(x + h)$, which has been shown to be of the form}
\[
\phi x + \phi' x\, h
       + \phi'' x\, \frac{h^{2}}{2}
       + \phi''' x\, \frac{h^{3}}{2·3} + \etc.,
\]
in the same manner as $\phi' x$~is the limit of~$\dfrac{dy}{dx}$ (\PageRef{23}),
so $\phi'' x$~is the limit of~$\dfrac{d^{2} y}{dx^{2}}$, $\phi''' x$~is that of~$\dfrac{d^{3} y}{dx^{3}}$, and so
forth.

From the manner in which the preceding table
was formed, the following relations are seen immediately:
\begin{gather*}
\begin{alignedat}{4}
y_{1} &= y     &&+ \Delta y    &
  \Delta y_{1} &= \Delta y    &&+ \Delta^{2} y\Add{,} \\
y_{2} &= y_{1} &&+ \Delta y_{1}\qquad &
  \Delta y_{2} &= \Delta y_{1} &&+ \Delta^{2} y_{1}\Add{,}
\end{alignedat} \\
\begin{alignedat}{3}
\Delta^{2} y_{1} &= \Delta^{2} y     &&+ \Delta^{3} y     && \etc.\Add{,} \\
\Delta^{2} y_{2} &= \Delta^{2} y_{1} &&+ \Delta^{3} y_{1}\ && \etc.
\end{alignedat}
\end{gather*}
Hence $y_{1}$,~$y_{2}$,~etc., can be expressed in terms of $y$,~$\Delta y$,
$\Delta^{2} y$,~etc. For $y_{1} = y + \Delta y$;
%[** TN: Next two displays in-line in the original]
\begin{align*}
y_{2} = y_{1} + \Delta y_{1}
  &= (y + \Delta y) + (\Delta y + \Delta^{2} y)
   = y + 2\Delta y + \Delta^{2} y. \displaybreak[0]
\intertext{In the same way $\Delta y_{2} = \Delta y + 2\Delta^{2} y + \Delta^{3} y$;
hence}
y_{3} = y_{2} + \Delta y_{2}
  &= (y + 2\Delta y + \Delta^{2} y) + (\Delta y + 2\Delta^{2} y + \Delta^{3} y) \\
  &= y + 3\Delta y + 3\Delta^{2} y + \Delta^{3} y.
\end{align*}

Proceeding in this way we have
\begin{alignat*}{5}
y_{1} = y &{}+{}& \Delta y\Add{\rlap{,}} \\
y_{2} = y &{}+{}&2\Delta y &{}+{}&  \Delta^{2} y\Add{\rlap{,}} \\
y_{3} = y &{}+{}&3\Delta y &{}+{}& 3\Delta^{2} y &{}+{}&  \Delta^{3} y\Add{\rlap{,}} \\
y_{4} = y &{}+{}&4\Delta y &{}+{}& 6\Delta^{2} y &{}+{}& 4\Delta^{3} y &{}+{}& \Delta^{5} y\Add{\rlap{,}} \\
y_{5} = y &{}+{}&5\Delta y &{}+{}&10\Delta^{2} y &{}+{}&10\Delta^{3} y &{}+{}&5\Delta^{5} y
  &+ \Delta^{6} y,\ \etc.\Add{,}
\end{alignat*}
from the whole of which it appears that $y_{n}$ or $\phi(x + n\Delta x)$
is a series consisting of $y$,~$\Delta y$,~etc., up to~$\Delta^{n} y$,
severally multiplied by the coefficients which occur in
the expansion $(1 + a)^{n}$, or
\PageSep{93}
\begin{align*}%[** TN: Re-formatted]
y_{n} &= \phi(x + n\Delta x) \\
  &= y + n\Delta y + n\frac{n - 1}{2}\, \Delta^{2} y
       + n\frac{n - 1}{2}\, \frac{n - 2}{3}\, \Delta^{3} y + \etc.
\end{align*}

Let us now suppose that $x$~becomes $x + h$ by $n$~equal
steps; that is, $x$,~$x + \dfrac{h}{n}$, $x + \dfrac{2h}{n}$, etc.~\dots\ $x + \dfrac{nh}{h}$
or~$x + h$, are the successive values of~$x$, so that
$n\Delta x = h$. Since the product of a number of factors is
not altered by multiplying one of them, provided we
divide another of them by the same quantity, multiply
every factor which contains~$n$ by~$\Delta x$, and divide the
accompanying difference of~$y$ by $\Delta x$ as often as there
are factors which contain~$n$, substituting~$h$ for~$n\Delta x$,
which gives
\begin{align*}
\phi(x + n\Delta x)
  &= y + n\Delta x\, \frac{\Delta y}{\Delta x}
       + n\Delta x\, \frac{n\Delta x - \Delta x}{2}\, \frac{\Delta^{2} y}{(\Delta x)^{2}} \\
  &\quad+ n\Delta x\, \frac{n\Delta x -  \Delta x}{2}\,
                 \frac{n\Delta x - 2\Delta x}{3}\, \frac{\Delta^{3} y}{(\Delta x)^{3}} + \etc.\Add{,}
\end{align*}
or
\begin{align*}
\iffalse %[** TN: Commented code matches the original]
\phi(x + h)
  &= \PadTo{y + n\Delta x}{y + h}\, \frac{\Delta y}{\Delta x}
       + \PadTo{n\Delta x}{h}\,
         \PadTo{\dfrac{n\Delta x - \Delta x}{2}}{\dfrac{h - \Delta x}{2}}\,
         \frac{\Delta^{2} y}{(\Delta x)^{2}} \\
  &+ \PadTo{n\Delta x}{h}\,
     \PadTo{\dfrac{n\Delta x -  \Delta x}{2}}{\dfrac{h -  \Delta x}{2}}\,
     \PadTo{\dfrac{n\Delta x - 2\Delta x}{3}}{\dfrac{h - 2\Delta x}{3}}\,
     \frac{\Delta^{3} y}{(\Delta x)^{3}} + \etc.
\fi
\phi(x + h)
  = y &+ h\, \frac{\Delta y}{\Delta x}
       + h\, \frac{h - \Delta x}{2}\, \frac{\Delta^{2} y}{(\Delta x)^{2}} \\
  &+ h\, \frac{h -  \Delta x}{2}\, \frac{h - 2\Delta x}{3}\,
         \frac{\Delta^{3} y}{(\Delta x)^{3}} + \etc.
\end{align*}
If $h$ remain the same, the more steps we make between
$x$~and~$x + h$, the smaller will each of those
steps be, and the number of steps may be increased,
until each of them is as small as we please. We can
therefore suppose $\Delta x$ to decrease without limit, without
affecting the truth of the series just deduced.
Write $dx$ for~$\Delta x$, etc., and recollect that $h - dx$,
$h - 2\, dx$, etc., continually approximate to~$h$. The series
then becomes
\[
\phi(x + h) = y + \frac{dy}{dx}\, h
  + \frac{d^{2} y}{dx^{2}}\, \frac{h^{2}}{2}
  + \frac{d^{3} y}{dx^{3}}\, \frac{h^{3}}{2·3} + \etc.\Add{,}
\]
\PageSep{94}
in which, according to the view taken of the symbols
$\dfrac{dy}{dx}$~etc.\ in \PageRefs{26}{27}, $\dfrac{dy}{dx}$~stands for the \emph{limit} of the
ratio of the increments, $\dfrac{dy}{dx}$~is $\phi' x$, $\dfrac{d^{2} y}{dx^{2}}$~is $\phi'' x$,~etc.
According to the method proposed in \PageRefs{28}{29},
the series written above is the first term of the development
of~$\phi(x + h)$, the remaining terms (which we
might include under an additional~$+$ etc.)\ being such
as to diminish without limit in comparison with the
first, when $dx$~is diminished without limit. And we
may show that the limit of~$\dfrac{d^{2} y}{dx^{2}}$ is the differential coefficient
of the limit of~$\dfrac{dy}{dx}$; or if by these fractions
themselves are understood their limits, that $\dfrac{d^{2} y}{dx^{2}}$ is the
differential coefficient of~$\dfrac{dy}{dx}$: for since $dy$, or $\phi(x + dx) - \phi x$,
becomes $dy + d^{2} y$, when $x$~is changed into
$x + dx$; and since $dx$~does not change in this process,
$\dfrac{dy}{dx}$ will become $\dfrac{dy}{dx} + \dfrac{d^{2} y}{dx}$, or its increment is~$\dfrac{d^{2} y}{dx}$. The
ratio of this to~$dx$ is~$\dfrac{d^{2} y}{(dx)^{2}}$, the limit of which, in the
definition of \PageRef{22}, is the differential coefficient of~$\dfrac{dy}{dx}$.
Similarly the limit of~$\dfrac{d^{3} y}{dx^{3}}$ is the differential coefficient
of the limit of~$\dfrac{d^{2} y}{dx^{2}}$; and so on.


\Subsection{Total and Partial Differential Coefficients.
Implicit Differentiation.}

We now proceed to apply the principles laid down,
\index{Differentiation!implicit|EtSeq}%
\index{Implicit!differentiation|EtSeq}%
to some cases in which the variable enters into its
function in a less direct and more complicated manner.
\PageSep{95}

For example, let $z$ be a given function of $x$~and~$y$,
and let $y$~be another given function of~$x$; so that $z$
contains $x$ both directly and indirectly; the latter as
it contains~$y$, which is a function of~$x$. This will be
the case if $z = x\log y$, where $y = \sin x$. If we were to
substitute for~$y$ its value in terms of~$x$, the value of~$z$
would then be a function of $x$~only; in the instance
just given it would be $x\log\sin x$. But if it be not convenient
to combine the two equations at the beginning
of the process, let us first consider $z$ as a function of
$x$~and~$y$, in which the two variables are independent.
In this case, if $x$~and~$y$ respectively receive the increments
$dx$~and~$dy$, the whole increment of~$z$, or~$d.z$, (or
at least that part which gives the limit of the ratios)
is represented by
\[
\frac{dz}{dx}\, dx + \frac{dz}{dy}\, dy.
\]
If $y$ be now considered as a function of~$x$, the consequence
is that $dy$, instead of being independent of~$dx$,
is a series of the form $p\, dx + q\, (dx)^{2} + \etc.$, in which $p$~is
the differential coefficient of~$y$ with respect to~$x$.
Hence
\[
d.z = \frac{dz}{dx}\, dx + \frac{dz}{dy}\, p\, dx \quad\text{or}\quad
\frac{d.z}{dx} = \frac{dz}{dx} + \frac{dz}{dx}\, p,
\]
in which the difference between $\dfrac{d.z}{dx}$ and $\dfrac{dz}{dx}$ is this,
that in the second, $x$~is only considered as varying
where it is directly contained in~$z$, or $z$~is considered
in the form in which it first appeared, as a function of
$x$~and~$y$, where $y$~is independent of~$x$; in the first, or
$\dfrac{d.z}{dx}$, the \emph{total variation} of~$z$ is denoted, that is, $y$~is
\index{Total!variations}%
\index{Variations, total}%
now considered as a function of~$x$, by which means if
$x$ become $x + dx$, $z$~will receive a different increment
\PageSep{96}
from that which it would have received, had $y$ been
independent of~$x$. {\Loosen In the instance above cited, where
$z = x\log y$ and $y = \sin x$, if the first equation be taken,
and $x$ becomes $x + dx$, $y$~remaining the same, $z$~becomes
$x\log y + \log y\, dx$ or $\dfrac{dz}{dx}$ is~$\log y$.} If $y$~only varies,
since (\PageRef{20}) $z$~will then become
\[
x\log y + x\, \frac{dy}{y} - \etc.,
\]
$\dfrac{dz}{dy}$ is~$\dfrac{x}{y}$ And $\dfrac{dy}{dx}$~is $\cos x$ when $y = \sin x$ (\PageRef{20}).
Hence $\dfrac{dz}{dx} + \dfrac{dz}{dy}\, p$, or $\dfrac{dz}{dx} + \dfrac{dz}{dy}\, \dfrac{dy}{dx}$ is $\log y + \dfrac{x}{y} \cos x$,
or $\log\sin x + \dfrac{x}{\sin x} \cos x$. This is~$\dfrac{d.z}{dx}$, which might
have been obtained by a more complicated process, if
$\sin x$ had been substituted for~$y$, before the operation
commenced. It is called the \emph{complete} or \emph{total} differential
\index{Coefficients, differential}%
\index{Complete Differential Coefficients}%
coefficient with respect to~$x$, the word \emph{total} indicating
that every way in which $z$ contains~$x$ has been
used; in opposition to~$\dfrac{dz}{dx}$, which is called the \emph{partial}
\index{Partial!differential coefficients}%
differential coefficient, $x$~having been considered as
varying only where it is directly contained in~$z$.

Generally, the complete differential coefficient of~$z$
with respect to~$x$, will contain as many terms as there
are different ways in which $z$ contains~$x$. From looking
at a complete differential coefficient, we may see
in what manner the function contained its variable.
Take, for example, the following,
\[
\frac{d.z}{dx}
  = \frac{dz}{dx} + \frac{dz}{dy}\, \frac{dy}{dx}
    + \frac{dz}{da}\, \frac{da}{dy}\, \frac{dy}{dx}
    + \frac{dz}{da}\, \frac{da}{dx}.
\]

Before proceeding to demonstrate this formula, we
will collect from itself the hypothesis from which it
\PageSep{97}
\index{Functions!direct and indirect}%
must have arisen. When $x$~is contained in~$z$, we shall
say that $z$~is a \emph{direct}\footnote
  {It may be right to warn the student that this phraseology is new, to the
  best of our knowledge. The nomenclature of the Differential Calculus has
  by no means kept pace with its wants; indeed the same may be said of algebra
  generally. [Written in~1832.---\Ed.]}
function of~$x$. When $x$~is contained
in~$y$, and $y$~is contained in~$z$, we shall say that
$z$~is an indirect function of~$x$ \emph{through}~$y$. It is evident
\index{Indirect function}%
that an indirect function may be reduced to one which
is direct, by substituting for the quantities which contain~$x$,
their values in terms of~$x$.

The first side of the equation~$\dfrac{d.z}{dx}$ is shown by the
point to be a complete differential coefficient, and indicates
that $z$~is a function of~$x$ in several ways; either
directly, and indirectly through one quantity at least,
or indirectly through several. If $z$~be a direct function
\index{Direct function}%
only, or indirectly through one quantity only, the
symbol~$\dfrac{dz}{dx}$, without the point, would represent its
total differential coefficient with respect to~$x$.

On the second side of the equation we see:

(1) $\dfrac{dz}{dx}$: which shows that $z$~is a direct function of~$x$,
and is that part of the differential coefficient which
we should get by changing $x$ into $x + dx$ throughout~$z$,
not supposing any other quantity which enters into~$z$
to contain~$x$.

(2) $\dfrac{dz}{dy}\, \dfrac{dy}{dx}$: which shows that $z$~is an indirect function
of~$x$ through~$y$. If $x$~and~$y$ had been supposed to
vary independently of each other, the increment of~$z$,
(or those terms which give the limiting ratio of this
increment to any other,) would have been $\dfrac{dz}{dx}\, dx + \dfrac{dz}{dy}\, dy$,
in which, if $dy$~had arisen from~$y$ being a function
\PageSep{98}
of~$x$, $dy$~would have been a series of the form
$p\, dx + q\, (dx)^{2} + \etc.$, of which only the differential coefficient~$p$
would have appeared in the limit. Hence
$\dfrac{dz}{dy}\, dy$ would have given~$\dfrac{dz}{dy}\, p$, or~$\dfrac{dz}{dy}\, \dfrac{dy}{dx}$.

(3) $\dfrac{dz}{da}\, \dfrac{da}{dy}\, \dfrac{dy}{dx}$: this arises from $z$ containing~$a$, which
contains~$y$, which contains~$x$. If $z$~had been differentiated
with respect to $a$~only, the increment would
have been represented by~$\dfrac{dz}{da}\, da$; if $da$~had arisen from
an increment of~$y$, this would have been expressed by
$\dfrac{dz}{da}\, \dfrac{da}{dy}\, dy$; if~$y$ had arisen from an increment given to~$x$,
this would have been expressed by $\dfrac{dz}{da}\, \dfrac{da}{dy}\, \dfrac{dy}{dx}\, dx$,
which, after $dx$~has been struck out, is the part of the
differential coefficient answering to that increment.

(4) $\dfrac{dz}{da}\, \dfrac{da}{dx}$: arising from $a$~containing $x$~directly,
and $z$~therefore containing $x$ indirectly through~$a$.

Hence $z$~is directly a function of $x$,~$y$, and~$a$, of
which $y$~is a function of~$x$, and $a$~of $y$~and~$x$.

If we suppose $x$,~$y$ and~$a$ to vary independently,
we have
\[
d.z = \frac{dz}{dx}\, dx + \frac{dz}{dy}\, dy + \frac{dz}{da}\, da + \etc.
\quad\text{(\PageRefs{28}{29})}.
\]
But as $a$~varies as a function of $y$~and~$x$,
\[
da = \frac{da}{dx}\, dx + \frac{da}{dy}\, dy.
\]
If we substitute this instead of~$da$, and divide by~$dx$,
taking the limit of the ratios, we have the result first
given.

For example, let (1) $z = x^{2} ya^{3}$, (2) $y = x^{2}$, and (3) $a = x^{3} y$.
Taking the first equation only, and substituting
\PageSep{99}
$x + dx$ for~$x$ etc., we find $\dfrac{dz}{dx} = 2xya^{3}$, $\dfrac{dz}{dy} = x^{2} a^{3}$,
and $\dfrac{dz}{da} = 3x^{2} ya^{2}$. From the second $\dfrac{dy}{dx} = 2x$, and from
the third $\dfrac{da}{dx} = 3x^{2} y$, and $\dfrac{da}{dy} = x^{3}$. Substituting these
in the value of~$\dfrac{d.z}{dx}$, we find
\begin{alignat*}{3}
%[** TN: Reformatted first line]
\frac{d.z}{dx} \text{ \ or \ }
   \frac{dz}{dx} &+ \frac{dz}{dy}\, \frac{dy}{dx}
    &&+ \frac{dz}{da}\, \frac{da}{dy}\, \frac{dy}{dx}
    &&+ \frac{dz}{da}\, \frac{da}{dx} \\
  = 2xya^{3} &+ x^{2} a^{3} × 2x &&+ 3x^{2} ya^{2} × x^{3} × 2x &&+ 3x^{2} ya^{2} × 3x^{2} y \\
  = 2xya^{3} &+ 2x^{3} a^{3} &&+ 6x^{6} ya^{2} &&+ 9x^{4} y^{2} a^{2}\Add{.}
\end{alignat*}
If for $y$~and~$a$ in the first equation we substitute their
values $x^{2}$ and~$x^{3} y$, or~$x^{5}$, we have $z = x^{19}$, the differential
coefficient of which\Add{ is}~$19x^{18}$. This is the same as
arises from the formula just obtained, after $x^{2}$~and~$x^{5}$
have been substituted for $y$~and~$a$; for this formula
then becomes
\[
2x^{18} + 2x^{18} + 6x^{18} + 9x^{18} \quad\text{or}\quad 19x^{18}.
\]

In saying that $z$~is a function of $x$~and~$y$, and that
$y$~is a function of~$x$, we have first supposed~$x$ to vary,
$y$~remaining the same. The student must not imagine
that $y$~is \emph{then} a function of~$x$; for if so, it would vary
when $x$~varied. There are two parts of the total differential
coefficient, arising from the direct and indirect
manner in which $z$ contains~$x$. That these two
parts may be obtained separately, and that their sum
constitutes the complete differential coefficient, is the
theorem we have proved. The first part~$\dfrac{dz}{dx}$ is what
\emph{would} have been obtained if $y$~had \emph{not} been a function
of~$x$; and on this supposition we therefore proceed to
find it. The other part $\dfrac{dz}{dy}\, \dfrac{dy}{dx}$ is the product (1)~of~$\dfrac{dz}{dy}$,
which would have resulted from a variation of $y$~only,
not considered as a function of~$x$; and (2)~of~$\dfrac{dy}{dx}$,
\PageSep{100}
the coefficient which arises from considering~$y$ as a
function of~$x$. These partial suppositions, however
useful in obtaining the total differential coefficient,
\index{Coefficients, differential}%
\index{Total!differential coefficient}%
cannot be separately admitted or used, except for this
purpose; since if $y$~be a function of~$x$, $x$~and~$y$ must
vary together.

If $z$~be a function of~$x$ in various ways, the theorem
obtained may be stated as follows:

Find the differential coefficient belonging to each
of the ways in which $z$ will contain~$x$, as if it were the
only way; the sum of these results (with their proper
signs) will be the total differential coefficient.

Thus, if $z$~only contains $x$ indirectly through~$y$,
$\dfrac{dz}{dx}$~is $\dfrac{dz}{dy}\, \dfrac{dy}{dx}$. If $z$ contains~$a$, which contains~$b$, which
contains~$x$, $\dfrac{dz}{dx} = \dfrac{dz}{da}\, \dfrac{da}{db}\, \dfrac{db}{dx}$.

This theorem is useful in the differentiation of complicated
\index{Differentiation!of complicated functions|EtSeq}%
functions; for example, let $z = \log(x^{2} + a^{2})$.
If we make $y = x^{2} + a^{2}$, we have $z = \log y$, and $\dfrac{dz}{dy} = \dfrac{1}{y}$;
while from the first equation $\dfrac{dy}{dx} = 2x$. Hence $\dfrac{dy}{dx}$ or
$\dfrac{dz}{dy}\, \dfrac{dy}{dx}$ is $\dfrac{2x}{y}$ or $\dfrac{2x}{x^{2} + a^{2}}$.

%[** TN: \log\log\sin x is never real valued when x is real]
If $z = \log\log\sin x$, or the logarithm of the logarithm
of~$\sin x$, let $\sin x = y$ and $\log y = a$; whence
$z= \log a$, and contains~$x$, because $a$ contains~$y$, which
contains~$x$. Hence
\[
\frac{dz}{dx} = \frac{dz}{da}\, \frac{da}{dy}\, \frac{dy}{dx};
\]
but since $z = \log a$,
\[
\frac{dz}{da} = \frac{1}{a};
\]
\PageSep{101}
since $a = \log y$,
\[
\frac{da}{dy} = \frac{1}{y};
\]
and since $y = \sin x$,
\[
\frac{dy}{dx} = \cos x.
\]
Hence
\[
\frac{dz}{dx} = \frac{dz}{da}\, \frac{da}{dy}\, \frac{dy}{dx}
  = \frac{1}{a}\, \frac{1}{y} \cos x
  = \frac{\cos x}{\log\sin x \sin x}.
\]

We now put some rules in the form of applications
of this theorem, though they may be deduced more
simply.


\Subsection[Applications of the Theorem for Implicit Differentiation.]
{Applications of the Preceding Theorem.}

(1) Let $z = ab$, where $a$~and~$b$ are functions of~$x$.
The general formula, since $z$ contains~$x$ indirectly
through $a$~and~$b$, is (in this case as well as in those
which follow,)
\[
\frac{dz}{dx}
  = \frac{dz}{da}\, \frac{da}{dx} + \frac{dz}{db}\, \frac{db}{dx}.
\]

We must leave $\dfrac{da}{dx}$ and $\dfrac{db}{dx}$ as we find them, until we
know \emph{what} functions $a$~and~$b$ are of~$x$; but as we
know what function $z$~is of $a$~and~$b$, we substitute for
$\dfrac{dz}{da}$ and~$\dfrac{dz}{db}$. Since $z = ab$, if $a$~becomes $a + da$, $z$~becomes
$ab + b\, da$, whence $\dfrac{dz}{db} = b$. In this case, and part
of the following, the limiting ratio of the increments
is the same as that of the increments themselves.
Similarly $\dfrac{dz}{db} = a$, whence from $z = ab$ follows
\[
\frac{dz}{dx} = b\, \frac{da}{dx} + a\, \frac{db}{dx}.
\]
\PageSep{102}

%[** TN: [sic] "become", twice]
(2) Let $z = \dfrac{a}{b}$. If $a$~become $a + da$, $z$~becomes
$\dfrac{a + da}{b}$ or $\dfrac{a}{b} + \dfrac{da}{b}$, and $\dfrac{dz}{da}$ is~$\dfrac{1}{b}$. If $b$~become $b + db$, $z$~becomes
$\dfrac{a}{b + db}$, or $\dfrac{a}{b} - \dfrac{a\, db}{b^{2}} + \etc.$, whence $\dfrac{dz}{db}$ is~$-\dfrac{a}{b^{2}}$\Add{.}
Hence from $z = \dfrac{a}{b}$ follows
\[
\frac{dz}{dx} = \frac{1}{b}\, \frac{da}{dx} - \frac{a}{b^{2}}\, \frac{db}{dx}
  = \frac{b\, \dfrac{da}{dx} - a\, \dfrac{db}{dx}}{b^{2}}.
\]

(3) Let $z = a^{b}$. Here $(a + da)^{b} = a^{b} + ba^{b-1}\, da + \etc.$\
(\PageRef{21}), whence $\dfrac{dz}{da} = ba^{b-1}$. Again,
$a^{b+db} = a^{b}\, a^{db} = a^{b}(1 + \log a\, db + \etc.)$ whence $\dfrac{dz}{db} = a^{b} \log a$.
Therefore from $z = a^{b}$ follows
\[
\frac{dz}{dx} = ba^{b-1}\, \frac{da}{dx} + a^{b} \log a\, \frac{db}{dx}.
\]


\Subsection{Inverse Functions.}

If $y$~be a function of~$x$, such as $y = \phi x$, we may,
\index{Functions!inverse|EtSeq}%
\index{Inverse functions|EtSeq}%
by solution of the equation, determine $x$ in terms of~$y$,
or produce another equation of the form $x = \psi y$.
For example, when $y = x^{2}$, $x = y^{\efrac{1}{2}}$. It is not necessary
that we should be able to solve the equation
$y = \phi x$ in finite terms, that is, so as to give a value
of~$x$ without infinite series; it is sufficient that $x$~can
be so expressed that the value of~$x$ corresponding to
any value of~$y$ may be found as near as we please
from $x = \psi y$, in the same manner as the value of~$y$
corresponding to any value of~$x$ is found from $y = \phi x$.

The equations $y = \phi x$, and $x = \psi y$, are connected,
being, in fact, the same relation in different forms;
and if the value of~$y$ from the first be substituted in
\PageSep{103}
the second, the second becomes $x = \psi(\phi x)$, or as it is
more commonly written, $\psi\phi x$. That is, the effect of
the operation or set of operations denoted by~$\psi$ is destroyed
by the effect of those denoted by~$\phi$; as in the
instances $(x^{2})^{\efrac{1}{2}}$, $(x^{3})^{\efrac{1}{3}}$, $e^{\log x}$, angle whose sine is~$(\sin x)$,
etc., each of which is equal to~$x$.

By differentiating the first equation $y = \phi x$, we obtain
$\dfrac{dy}{dx} = \phi' x$, and from the second $\dfrac{dx}{dy} = \psi' y$. But
whatever values of $x$~and~$y$ together satisfy the first
equation, satisfy the second also; hence, if when $x$~becomes
$x + dx$ in the first, $y$~becomes $y + dy$; the same
$y + dy$ substituted for~$y$ in the second, will give the
same $x + dx$. Hence $\dfrac{dx}{dy}$ as deduced from the second,
and $\dfrac{dy}{dx}$ as deduced from the first, are reciprocals for
every value of~$dx$. The limit of one is therefore the
reciprocal of the limit of the other; the student may
easily prove that if $a$~is always equal to~$\dfrac{1}{b}$, and if $a$~continually
approaches to the limit~$\alpha$, while $b$~at the
same time approaches the limit~$\beta$, $\alpha$~is equal to~$\dfrac{1}{\beta}$.
But $\dfrac{dx}{dy}$ or $\psi' y$, deduced from $x = \psi y$, is expressed in
terms of~$y$, while $\dfrac{dy}{dx}$ or $\phi' x$, deduced from $y = \phi x$ is
expressed in terms of~$x$. Therefore $\psi' y$ and $\phi' x$ are
reciprocals for all such values of $x$~and~$y$ as satisfy
either of the two first equations.

For example let $y = e^{x}$, from which $x = \log y$. From
the first (\PageRef{20}) $\dfrac{dy}{dx} = e^{x}$; from the second $\dfrac{dx}{dy} = \dfrac{1}{y}$;
and it is evident that $e^{x}$~and~$\dfrac{1}{y}$ are reciprocals, whenever
$y = e^{x}$.

If we differentiate the above equations twice, we get
\PageSep{104}
$\dfrac{d^{2} y}{dx^{2}} = \phi'' x$, and $\dfrac{d^{2} x}{dy^{2}} = \psi'' x$. There is no very obvious
analogy between $\dfrac{d^{2} y}{dx^{2}}$ and $\dfrac{d^{2} x}{dy^{2}}$; indeed no such appears
from the method in which these coefficients were first
formed. Turn to the table in \PageRef{90}, and substitute
$d$ for~$\Delta$ throughout, to indicate that the increments
may be taken as small as we please. We there substitute
in~$\phi x$ what we will call a set of \emph{equidistant} values
\index{Equidistant values}%
\index{Values!equidistant}%
of~$x$, or values in arithmetical progression, viz.,
$x$,~$x + dx$, $x + 2\, dx$,~etc. The resulting values of~$y$,
or $y$,~$y_{1}$, etc., are not equidistant, except in one function
only, when $y = ax + b$, where $a$~and~$b$ are constant.
Therefore $dy$,~$dy_{1}$, etc., are not equal; whence
arises the next column of second differences, or $d^{2} y$,
$d^{2} y_{1}$, etc. The limiting ratio of $d^{2} y$ to~$(dx)^{2}$, expressed
by~$\dfrac{d^{2} y}{dx^{2}}$, is the second differential coefficient of~$y$ with
respect to~$x$. If from $y = \phi x$ we deduce $x = \psi y$, and
take a set of equidistant values of~$y$, viz., $y$,~$y + dy$,
$y + 2\, dy$, etc., to which the corresponding values of~$x$
are $x$,~$x_{1}$, $x_{2}$,~etc., a similar table may be formed,
which will give $dx$,~$dx_{1}$, etc., $d^{2} x$,~$d^{2} x_{1}$, etc., and the
limit of the ratio of~$d^{2} x$ to~$(dy)^{2}$ or $\dfrac{d^{2} x}{dy^{2}}$ is the second
differential coefficient of~$x$ with respect to~$y$. These
are entirely different suppositions, $dx$~being given in
the first table, and $dy$~varying; while in the second $dy$~is
given and $dx$~varies. We may show how to deduce
one from the other as follows:

When, as before, $y = \phi x$ and $x = \psi y$, we have
\[
\frac{dy}{dx} = \phi' x = \frac{1}{\psi' y} = \frac{1}{p},
\]
if $\psi' y$ be called~$p$. Calling this~$u$, and considering it
\PageSep{105}
as a function of~$x$ from containing~$p$, which contains~$y$,
which contains~$x$, we have
\[
\frac{du}{dp}\, \frac{dp}{dy}\, \frac{dy}{dx}
\]
for its differential coefficient with respect to~$x$. But
since
\[
u = \frac{1}{p},
\]
therefore
\[
\frac{du}{dp} = -\frac{1}{p^{2}};
\]
since $p = \psi' y$, therefore
\[
\frac{dp}{dy} = \psi'' y;
\]
and $\psi'' y$ is the differential coefficient of~$\psi' y$, and is
$\dfrac{d^{2} x}{dy^{2}}$. Also $\dfrac{1}{p^{2}}$~is
\[
\frac{1}{(\psi' y)^{2}} \quad\text{or}\quad
(\phi' x)^{2} \quad\text{or}\quad
\left(\frac{dy}{dx}\right)^{2}.
\]
Hence the differential coefficient of $u$ or~$\dfrac{dy}{dx}$, with respect
to~$x$, which is~$\dfrac{d^{2} y}{dx^{2}}$, is also
\[
-\left(\frac{dy}{dx}\right)^{2} \frac{d^{2} x}{dy^{2}}\, \frac{dy}{dx}
\quad\text{or}\quad
-\left(\frac{dy}{dx}\right)^{3} \frac{d^{2} x}{dy^{2}}.
\]

{\Loosen If $y = e^{x}$, whence $x = \log y$, we have $\dfrac{dy}{dx} = e^{x}$ and
$\dfrac{d^{2} y}{dx^{2}} = e^{x}$. But $\dfrac{dx}{dy} = \dfrac{1}{y}$ and $\dfrac{d^{2} x}{dy^{2}} = -\dfrac{1}{y^{2}}$. Therefore}
\[
-\left(\frac{dy}{dx}\right)^{3} \frac{d^{2} x}{dy^{2}} \quad\text{is}\quad
-e^{3x} \left(-\frac{1}{y^{2}}\right) \quad\text{or}\quad
\frac{e^{3x}}{y^{2}} \quad\text{or}\quad
\frac{e^{3x}}{e^{2x}},
\]
which is~$e^{x}$, the value just found for~$\dfrac{d^{2} y}{dx^{2}}$.
\PageSep{106}

In the same way $\dfrac{d^{3} y}{dx^{3}}$ might be expressed in terms
of $\dfrac{dy}{dx}$, $\dfrac{d^{2} y}{dx^{2}}$, and~$\dfrac{d^{3} x}{dy^{3}}$; and so on.


\Subsection{Implicit Functions.}

The variable which appears in the denominator of
the differential coefficients is called the \emph{independent}
variable. In any function, one quantity at least is
changed at pleasure; and the changes of the rest,
with the limiting ratio of the changes, follow from the
form of the function. The number of independent
variables depends upon the number of quantities
\index{Variables!independent and dependent}%
which enter into the equations, and upon the number
of equations which connect them. If there be only
one equation, all the variables except one are independent,
or may be changed at pleasure, without ceasing
to satisfy the equation; for in such a case the
common rules of algebra tell us, that as long as one
quantity is left to be determined from the rest, it can
be determined by one equation; that is, the values of
all but one are at our pleasure, it being still in our
power to satisfy one equation, by giving a proper
value to the remaining one. Similarly, if there be
two equations, all variables except two are independent,
and so on. If there be two equations with two
unknown quantities only, there are no variables; for
by algebra, a finite number of values, and a finite
number only, can satisfy these equations; whereas it
is the nature of a variable to receive any value, or at
least any value which will not give impossible values
for other variables. If then there be $m$~equations containing
$n$~variables, ($n$~must be greater than~$m$), we
have $n - m$~independent variables, to each of which
\index{Independent variables}%
\PageSep{107}
we may give what values we please, and by the equations,
deduce the values of the rest. We have thus
various sets of differential coefficients, arising out of
the various choices which we may make of independent
variables.

If, for example, $a$,~$b$, $x$,~$y$, and~$z$, being variables,
we have
\begin{align*}
\phi(a, b, x, y, z) &= 0, \\
\psi(a, b, x, y, z) &= 0, \\
\chi(a, b, x, y, z) &= 0,
\end{align*}
we have two independent variables, which may be
either $x$~and~$y$, $x$~and~$z$, $a$~and~$b$, or any other combination.
If we choose $x$~and~$y$, we should determine
$a$,~$b$, and~$z$ in terms of $x$~and~$y$ from the three equations;
in which case we can obtain
\[
\frac{da}{dx},\quad \frac{da}{dy},\quad \frac{db}{dx},\quad \etc.
\]

When $y$~is a function of~$x$, as in $y = \phi x$, it is called
\index{Explicit functions}%
\index{Functions!implicit and explicit}%
an \emph{explicit} function of~$x$. This equation tells us not
only that $y$~is a function of~$x$, but also what function
it is. The value of~$x$ being given, nothing more is
necessary to determine the corresponding value of~$y$,
than the substitution of the value of~$x$ in the several
terms of~$\phi x$.

But it may happen that though $y$~is a function of~$x$,
\index{Implicit!function}%
the relation between them is contained in a form
from which $y$~must be deduced by the solution of an
equation. For example, in $x^{2} - xy + y^{2} = a$, when $x$~is
known, $y$~must be determined by the solution of an
equation of the second degree. Here, though we know
that $y$~must be a function of~$x$, we do not know, without
further investigation, what function it is. In this
case $y$~is said to be \emph{implicitly} a function of~$x$, or an implicit
\PageSep{108}
function. By bringing all the terms on one side
of the equation, we may always reduce it to the form
$\phi(x, y) = 0$. Thus, in the case just cited, we have
$x^{2} - xy + y^{2} - a = 0$.

{\Loosen We now want to deduce the differential coefficient
$\dfrac{dy}{dx}$ from an equation of the form $\phi(x, y) = 0$. If we
take the equation $u = \phi(x, y)$, in which when $x$~and~$y$
become $x + dx$ and $y + dy$, $u$~becomes $u + du$, we have,
by our former principles,}
\[
du = \ux\, dx + \uy\, dy + \etc., \text{(\PageRef{82})},
\]
in which $\ux$~and~$\uy$ can be directly obtained from the
equation, as in \PageRef{82}. Here $x$~and~$y$ are independent,
as also $dx$~and~$dy$; whatever values are given to
them, it is sufficient that $u$~and~$du$ satisfy the two last
equations. But if $x$~and~$y$ must be always so taken
that $u$ may~$= 0$, (which is implied in the equation
$\phi(x, y) = 0$,) we have $u = 0$, and $du = 0$; and this,
whatever may be the values of $dx$ and~$dy$. Hence $dx$
and~$dy$ are connected by the equation
\[
0 = \ux\, dx + \uy\, dy + \etc.,
\]
and their limiting ratio must be obtained by the equation
\[
\ux\, dx + \uy\, dy = 0, \quad\text{or}\quad \frac{dy}{dx} = -\frac{\ux}{\uy};
\]
{\Loosen $y$~and~$x$ are no longer independent; for, one of them
being given, the other must be so taken that the equation
$\phi(x, y) = 0$ may be satisfied. The quantities $\ux$
\index{Functions!implicit and explicit}%
\index{Implicit!function}%
and~$\uy$ we have denoted by $\dfrac{du}{dx}$ and~$\dfrac{du}{dy}$, so that}
\[
\frac{dy}{dx} = -\frac{\;\dfrac{du}{dx}\;}{\dfrac{du}{dy}}\Add{.}
\Tag{(1)}
\]
\PageSep{109}

We must again call attention to the different meanings
of the same symbol~$du$ in the numerator and denominator
of the last fraction. Had $du$, $dx$, and~$dy$
been common algebraical quantities, the first meaning
the same thing throughout, the last equation would
not have been true until the negative sign had been
removed. We will give an instance in which $du$~shall
mean the same thing in both.

Let $u = \Chg{\phi(x)}{\phi x}$, and let $u = \psi y$, in which two equations
is implied a third $\phi x = \psi y$; and $y$~is a function
of~$x$. Here, $x$~being given, $u$~is known from the first
equation; and $u$~being known, $y$~is known from the
second. Again, $x$~and~$dx$ being given, $du$, which is
$\phi(x + dx) - \phi x$ is known, and being substituted in
the result of the second equation, we have $du = \psi(y + dy) - \psi y$,
which $dy$~must be so taken as to
satisfy. From the first equation we deduce $du = \phi'x\, dx + \etc.$\
and from the second $du = \psi' y\, dy + \etc.$,
whence
\[
\phi' x\, dx + \etc. = \psi' y\, dy + \etc.;
\]
the \emph{etc.}\ only containing terms which disappear in finding
the limiting ratios. Hence,
\[
\frac{dy}{dx} = \frac{\phi' x}{\psi' y}
  = \frac{\;\dfrac{du}{dx}\;}{\dfrac{du}{dy}}\Add{,}
\Tag{(2)}
\]
a result in accordance with common algebra.

But the equation~\Eq{(1)} was obtained from $u = \phi(x, y)$,
on the supposition that $x$~and~$y$ were always so taken
that $u$ should~$= 0$, while \Eq{(2)}~was obtained from $u = \Chg{\phi(x)}{\phi x}$
and $u = Sy$, in which no new supposition can be
made; since one more equation between $u$,~$x$, and~$y$
would give three equations connecting these three
quantities, in which case they would cease to be variable
(\PageRef{106}).
\PageSep{110}

As an example of~\Eq{(1)} let $xy - x = 1$, or $xy - x - 1 = 0$.
From $u = xy - x - 1$ we deduce (\PageRef{81})
$\dfrac{du}{dx} = y - 1$, $\dfrac{du}{dy} = x$; whence, by equation~\Eq{(1)},
\[
\frac{dy}{dx} = -\frac{y - 1}{x}.
\Tag{(3)}
\]
By solution of $xy - x = 1$, we find $y = 1 + \dfrac{1}{x}$, and
\[
dy = \left(1 + \frac{1}{x + dx}\right) - \left(1 + \frac{1}{x}\right)
   = -\frac{dx}{x^{2}} + \etc.\footnote{See \PageRef{26}.}
\]
Hence $\dfrac{dy}{dx}$ (meaning the limit) is~$-\dfrac{1}{x^{2}}$, which will also
be the result of~\Eq{(3)} if $1 + \dfrac{1}{x}$ be substituted for~$y$.


\Subsection{Fluxions, and the Idea of Time.}

To follow this subject farther would lead us beyond
\index{Time, idea of|EtSeq}%
our limits; we will therefore proceed to some
observations on the differential coefficient, which, at
this stage of his progress, may be of use to the student,
who should never take it for granted that because
he has made some progress in a science, he understands
the first principles, which are often, if not
always, the last to be learned well. If the mind were
so constituted as to receive with facility any perfectly
new idea, as soon as the same was legitimately applied
in mathematical demonstration, it would doubtless
be an advantage not to have any notion upon a
mathematical subject, previous to the time when it is
to become a subject of consideration after a strictly
mathematical method.

This not being the case, it is a cause of embarrassment
to the student, that he is introduced at once to a
definition so refined as that of the limiting ratio which
\PageSep{111}
the increment of a function bears to the increment of
its variable. Of this he has not had that previous experience,
which is the case in regard to the words
\emph{force}, \emph{velocity}, or \emph{length}. Nevertheless, he can easily
\index{Velocity!linear}%
conceive a mathematical quantity in a state of continuous
increase or decrease, such as the distance between
two points, one of which is in motion. The
number which represents this line (reference being
made to a given linear unit) is in a corresponding
state of increase or decrease, and so is every function
of this number, or every algebraical expression in the
formation of which it is required. And the nature of
the change which takes place in the function, that is,
whether the function will increase or decrease when
the variable increases; whether that increase or decrease
corresponding to a given change in the variable
will be smaller or greater, etc., depends on the
manner in which the variable enters as a component
part of its function.

Here we want a new word, which has not been invented
for the world at large, since none but mathematicians
consider the subject; which word, if the
change considered were change of place, depending
upon change of time, would be \emph{velocity}. Newton
adopted this word, and the corresponding idea, expressing
many numbers in succession, instead of at
once, by supposing a point to generate a straight line
by its motion, which line would at different instants
contain any different numbers of linear units.

To this it was objected that the idea of \emph{time} is introduced,
which is foreign to the subject. We may
answer that the notion of time is only necessary, inasmuch
as we are not able to consider more than one
thing at a time. Imagine the diameter of a circle divided
\PageSep{112}
into a million of equal parts, from each of which
a perpendicular is drawn meeting the circle. A mind
which could at a view take in every one of these lines,
and compare the differences between every two contiguous
perpendiculars with one another, could, by
subdividing the diameter still further, prove those
propositions which arise from supposing a point to
move uniformly along the diameter, carrying with it
a perpendicular which lengthens or shortens itself so
as always to have one extremity on the circle. But
we, who cannot consider all these perpendiculars at
once, are obliged to take one after another. If one
perpendicular only were considered, and the differential
coefficient of that perpendicular deduced, we might
certainly appear to avoid the idea of time; but if all
the states of a function are to be considered, corresponding
to the different states of its variable, we
have no alternative, with our bounded faculties, but
to consider them in succession; and succession, disguise
it as we may, is the identical idea of time introduced
in Newton's Method of Fluxions.
\index{Fluxions}%


\Subsection{The Differential Coefficient Considered with Respect
to its Magnitude.}

The differential coefficient corresponding to a particular
\index{Coefficients, differential}%
\index{Contiguous values}%
\index{Differential coefficients!as the index of the change of a function}%
\index{Logarithms|EtSeq}%
\index{Values!contiguous}%
value of the variable, is, if we may use the
phrase, the \emph{index} of the change which the function
would receive if the value of the variable were increased.
Every value of the variable, gives not only
a different value to the function, but a different quantity
of increase or decrease in passing to what we may
call \emph{contiguous} values, obtained by a given increase of
the variable.

If, for example, we take the common logarithm of~$x$,
\PageSep{113}
and let $x$ be~$100$, we have common $\log 100 = 2$. If
$x$~be increased by~$2$, this gives common $\log 102 = 2.0086002$,
the ratio of the increment of the function
\index{Increment}%
to that of the variable being that of $.0086002$ to~$2$, or
$.0043001$. In passing from $1000$ to~$1003$, we have the
logarithms $3$ and~$3.0013009$, the above-mentioned ratio
being~$.0004336$, little more than a tenth of the
former. We do not take the increments themselves,
but the proportion they bear to the changes in the
variable which gave rise to them; so in estimating
the rate of motion of two points, we either consider
lengths described in the same time, or if that cannot
be done, we judge, not by the lengths described in
different times, but by the proportion of those lengths
to the times, or the proportions of the units which
express them.

The above rough process, though from it some
might draw the conclusion that the logarithm of~$x$ is
increasing faster when $x = 100$ than when $x = 1000$,
is defective; for, in passing from $100$ to~$102$, the
change of the logarithm is not a sufficient index of the
change which is taking place when $x$ is~$100$; since,
for any thing we can be supposed to know to the contrary,
the logarithm might be decreasing when $x = 100$,
and might afterwards begin to increase between
$x = 100$ and $x = 102$, so as, on the whole, to cause
the increase above mentioned. The same objection
would remain good, however small the increment
might be, which we suppose $x$ to have. If, for example,
we suppose $x$ to change from $x = 100$ to $x = 100.00001$,
which increases the logarithm from~$2$ to~$2.00000004343$,
we cannot yet say but that the logarithm
may be decreasing when $x = 100$, and may begin
to increase between $x = 100$ and $x = 100.00001$.
\PageSep{114}

In the same way, if a point is moving, so that at
the end of $1$~second it is at $3$~feet from a fixed point,
and at the end of $2$~seconds it is at $5$~feet from the
fixed point, we cannot say which way it is moving at
the end of one second. \emph{On the whole}, it increases its
distance from the fixed point in the second second;
but it is possible that at the end of the first second it
may be moving back towards the fixed point, and may
turn the contrary way during the second second. And
the same argument holds, if we attempt to ascertain
the way in which the point is moving by supposing
any finite portion to elapse after the first second. But
if on adding any interval, \emph{however small}, to the first
second, the moving point does, during that interval,
increase its distance from the fixed point, we can then
certainly say that at the end of the first second the
point is moving from the fixed point.

On the same principle, we cannot say whether the
logarithm of~$x$ is increasing or decreasing when $x$~increases
and becomes~$100$, unless we can be sure that
any increment, however small, added to~$x$, will increase
the logarithm. Neither does the ratio of the
increment of the function to the increment of its variable
furnish any distinct idea of the change which is
taking place when the variable has attained or is passing
through a given value. For example, when $x$~passes
from $100$ to~$102$, the difference between $\log 102$
and $\log 100$ is the united effect of all the changes
which have taken place between $x = 100$ and $x = 100\frac{1}{10}$;
$x = 100\frac{1}{10}$ and $x = 100\frac{2}{10}$, and so on. Again,
the change which takes place between $x = 100$ and
$x = 100\frac{1}{10}$ may be further compounded of those which
take place between $x = 100$ and $x = 100\frac{1}{100}$; $x = 100\frac{1}{100}$
and $x = 100\frac{2}{100}$, and so on. The objection
\PageSep{115}
becomes of less force as the increment diminishes,
but always exists unless we take the limit of the ratio
of the increments, instead of that ratio.

How well this answers to our previously formed
ideas on such subjects as direction, velocity, and
force, has already appeared.


\Subsection{The Integral Calculus.}

We now proceed to the Integral Calculus, which
\index{Integral Calculus|EtSeq}%
is the inverse of the Differential Calculus, as will afterwards
appear.

We have already shown, that when two functions
\emph{increase} or \emph{decrease} without limit, their \emph{ratio} may either
increase or decrease without limit, or may tend to
some finite limit. Which of these will be the case depends
upon the manner in which the functions are related
to their variable and to one another.

This same proposition may be put in another form,
as follows: If there be two functions, the first of which
\emph{decreases} without limit, on the same supposition which
makes the second \emph{increase} without limit, the \emph{product}
of the two may either remain finite, and never exceed
a certain finite limit; or it may increase without limit,
or diminish without limit.

For example, take $\cos\theta$ and~$\tan\theta$. As the angle~$\theta$
\emph{approaches} a right angle, $\cos\theta$~diminishes without
limit; it is nothing when $\theta$~\emph{is} a right angle; and any
fraction being named, $\theta$~can be taken so near to a
right angle that $\cos\theta$~shall be smaller. Again, as $\theta$~approaches
to a right angle, $\tan\theta$~increases without
limit; it is called \emph{infinite} when $\theta$~is a right angle, by
which we mean that, let any number be named, however
great, $\theta$~can be taken so near a right angle that
$\tan\theta$~shall be greater. Nevertheless the product $\cos\theta × \tan\theta$,
\PageSep{116}
of which the first factor diminishes without limit,
while the second increases without limit, is always
finite, and tends towards the limit~$1$; for $\cos\theta × \tan\theta$
is always~$\sin\theta$, which last approaches to~$1$ as $\theta$~approaches
to a right angle, and is~$1$ when $\theta$~\emph{is} a right
angle.

Generally, if $A$~diminishes without limit at the
same time as $B$~increases without limit, the product~$AB$
may, and often will, tend towards a finite limit.
This product~$AB$ is the representative of~$A$ divided by~$\dfrac{1}{B}$
or the ratio of $A$ to~$\dfrac{1}{B}$. If $B$~increases without
limit, $\dfrac{1}{B}$~decreases without limit; and as $A$~also decreases
without limit, the ratio of $A$ to~$\dfrac{1}{B}$ may have a
finite limit. But it may also diminish without limit;
as in the instance of $\cos^{2}\theta × \tan\theta$, when $\theta$~approaches
to a right angle. Here $\cos^{2}\theta$~diminishes without limit,
and $\tan\theta$~increases without limit; but $\cos^{2}\theta × \tan\theta$
being $\cos\theta × \sin\theta$, or a diminishing magnitude multiplied
by one which remains finite, diminishes without
limit. Or it may increase without limit, as in the case
of $\cos\theta × \tan^{2}\theta$, which is also $\sin\theta × \tan\theta$; which last
has one factor finite, and the other increasing without
limit. We shall soon see an instance of this.

If we take any numbers, such as $1$~and~$2$, it is evident
that between the two we may interpose any number
of fractions, however great, either in arithmetical
progression, or according to any other law. Suppose,
for example, we wish to interpose $9$~fractions in arithmetical
progression between $1$~and~$2$. These are $1\frac{1}{10}$,
$1\frac{2}{10}$,~etc., up to~$1\frac{9}{10}$; and, generally, if $m$~fractions in
arithmetical progression be interposed between $a$~and~$a + h$,
the complete series is
\PageSep{117}
\begin{multline*}
a,\quad a + \frac{h}{m + 1},\quad
        a + \frac{2h}{m + 1},\quad \etc.\Add{,} \dots\\
\text{up to } a + \frac{mh}{m + 1},\quad a + h\Add{.}
\Tag{(1)}
\end{multline*}
The sum of these can evidently be made as great as
we please, since no one is less than the given quantity~$a$,
and the number is as great as we please. Again,
if we take~$\phi x$, any function of~$x$, and let the values
just written be successively substituted for~$x$, we shall
have the series
\begin{multline*}
\phi a,\quad \phi\left(a + \frac{h}{m + 1}\right),\quad
  \phi\left(a + \frac{2h}{m + 1}\right),\quad \etc.\Add{,} \dots\\
\text{up to } \phi(a + h);
\Tag{(2)}
\end{multline*}
the sum of which may, in many cases, also be made
as great as we please by sufficiently increasing the
number of fractions interposed, that is, by sufficiently
increasing~$m$. But though the two sums increase without
limit when $m$~increases without limit, it does not
therefore follow that their ratio increases without
limit; indeed we can show that this cannot be the
case when all the separate terms of~\Eq{(2)} remain finite.

For let $A$~be greater than any term in~\Eq{(2)}, whence,
as there are $(m + 2)$~terms, $(m + 2)A$~is greater than
their sum. Again, every term of~\Eq{(1)}, except the first,
being greater than~$a$, and the terms being $m + 2$~in
number, $(m + 2)a$~is less than the sum of the terms in~\Eq{(1)}.
Consequently,
\[
\frac{(m + 2)A}{(m + 2)a}
  \text{ is greater than the ratio }
  \frac{\text{sum of terms in~\Eq{(2)}}}{\text{sum of terms in~\Eq{(1)}}},
\]
since its numerator is greater than the last numerator,
and its denominator less than the last denominator.
But
\PageSep{118}
\[
\frac{(m + 2)A}{(m + 2)a} = \frac{A}{a},
\]
which is independent of~$m$, and is a finite quantity.
Hence the ratio of the sums of the terms is always
finite, whatever may be the number of terms, at least
unless the terms in~\Eq{(2)} increase without limit.

As the number of interposed values increases, the
interval or difference between them diminishes; if,
therefore, we multiply this difference by the sum of
the values, or form
\begin{multline*}
\frac{h}{m + 1} \Biggl[
  \phi a + \phi\left(a + \frac{h}{m + 1}\right) + \\
  \phi\left(a + \frac{2h}{m + 1}\right) \Add{+} \dots + \phi(a + h)
\Biggr]\Add{,}
\end{multline*}
we have a product, one term of which diminishes, and
the other increases, when $m$~is increased. The product
\emph{may} therefore remain finite, or never pass a certain
limit, when $m$~is increased without limit, and we
shall show that this \emph{is} the case.

As an example, let the given function of~$x$ be~$x^{2}$,
and let the intermediate values of~$x$ be interposed between
$x = a$ and $x = a + h$. Let $v = \dfrac{h}{m + 1}$, whence
the above-mentioned product is
\begin{multline*}%[** TN: Re-formatted from the original]
v\bigl\{a^{2} + (a + v)^{2} + (a + 2v)^{2} + \dots
  + \bigl(a + (m + 1)v\bigr)^{2}\bigr\} \\
  = (m + 2)va^{2} + 2av^{2} \{1 + 2 + 3 + \dots + (m + 1)\} \\
  + v^{3} \{1^{2} + 2^{2} + 3^{2} + \dots + (m + 1)^{2}\};
\end{multline*}
{\Loosen of which, $1 + 2 + \dots + (m + 1) = \frac{1}{2}(m + 1)(m + 2)$
and (\PageRef{73}), $1^{2} + 2^{2} + \dots + (m + 1)^{2}$ approaches
without limit to a ratio of equality with $\frac{1}{3}(m + 1)^{3}$,
when $m$~is increased without limit. Hence this last
sum may be put under the form $\frac{1}{3}(m + 1)^{3} (1 + \alpha)$,
\PageSep{119}
where $\alpha$~diminishes without limit when $m$~is increased
without limit. Making these substitutions, and putting
for~$v$ its value $\dfrac{h}{m + 1}$, the above expression becomes}
\[
\frac{m + 2}{m + 1}\, ha^{2} + \frac{m + 2}{m + 1}\, ha^{2}
  + (1 + \alpha)\, \frac{h^{3}}{3},
\]
in which $\dfrac{m + 2}{m + 1}$ has the limit~$1$ when $m$~increases without
limit, and $1 + \alpha$~has also the limit~$1$, since, in that
case, $\alpha$~diminishes without limit. Therefore the limit
of the last expression is
\[
ha^{2} + ha^{2} + \frac{h^{3}}{3} \quad\text{or}\quad \frac{(a + h)^{3} - a^{3}}{3}.
\]

{\Loosen This result may be stated as follows: If the variable~$x$,
setting out from a value~$a$, becomes successively
$a + dx$, $a + 2\,dx$, etc., until the total increment
is~$h$, the smaller $dx$ is taken, the more nearly will the
sum of all the values of~$x^{2}\, dx$, or $a^{2}\, dx + (a + dx)^{2}\, dx + (a + 2\, dx)^{2}\, dx + \etc.$,
be equal to}
\[
\frac{(a + h)^{3} - a^{3}}{3},
\]
and to this the aforesaid sum may be brought within
any given degree of nearness, by taking $dx$ sufficiently
small.

This result is called the \emph{integral} of~$x^{2}\, dx$, between
\index{Integral Calculus!notation of}%
\index{Integrals!definition of|EtSeq}%
\index{Notation!of the Integral Calculus}%
the limits $a$~and~$a + h$, and is written $\int x^{2}\, dx$, when it
is not necessary to specify the limits, $\int_{a}^{a+h} x^{2}\, dx$,
or\footnote
  {This notation $\int x^{2}\, dx\Ibar_{a}^{a+h}$ appears to me to avoid the objections which
  may be raised against $\int_{a+h}^{a} x^{2}\, dx$ as contrary to analogy, which would require
  that $\int^{2} x^{2}\, dx^{2}$ should stand for the second integral of~$x^{2}\, dx$. It will be found
  convenient in such integrals as $\int z\, dx\Ibar_{b}^{a}\, dy\Ibar_{0}^{\phi x}$. There is as yet no general agreement
  on this point of notation.---\textit{De~Morgan}, 1832.}
$\int x^{2}\, dx\Ibar_{a}^{a+h}$, or $\int x^{2}\, dx\Ibar_{x=a}^{x=a+h}$ in the contrary case. We
\PageSep{120}
now proceed to show the connexion of this process
with the principles of the Differential Calculus.


\Subsection{Connexion of the Integral with the Differential
Calculus.}

Let $x$ have the successive values $a$, $a + dx$, $a + 2\, dx$,
etc.,~\dots\ up to $a + m\, dx$, or $a + h$, $h$~being a given
quantity, and $dx$ the $m$\th~part of~$h$, so that as $m$~is increased
without limit, $dx$~is diminished without limit.
Develop the successive values $\phi x$, or $\phi a$, $\phi(a + dx)$\Add{,}~\dots\
(\PageRef{21}),\par
{\footnotesize\begin{alignat*}{6}
& \phi a &&= \phi a\Add{,} \\
&\phi(a + dx) &&= \phi a &&+ \phi' a\, dx
  &&+ \phi'' a\, \frac{(dx)^{2}}{2}
  &&+ \phi''' a\, \frac{(dx)^{3}}{2·3} &&+ \etc.\Add{,} \displaybreak[0] \\
&\phi(a + 2\, dx) &&= \phi a &&+ \phi' a\, 2\, dx
  &&+ \phi'' a\, \frac{(2\, dx)^{2}}{2}
  &&+ \phi''' a\, \frac{(2\, dx)^{3}}{2·3} &&+ \etc.\Add{,} \displaybreak[0] \\
&\phi(a + 3\, dx) &&= \phi a &&+ \phi' a\, 3\, dx
  &&+ \phi'' a\, \frac{(3\, dx)^{2}}{2}
  &&+ \phi''' a\, \frac{(3\, dx)^{3}}{2·3} &&+ \etc.\Add{,} \\
\DotRow{12} \\
&\phi(a + m\, dx) &&= \phi a &&+ \phi' a\, m\, dx
  &&+ \phi'' a\, \frac{(m\, dx)^{2}}{2}
  &&+ \phi''' a\, \frac{(m\, dx)^{3}}{2·3} &&+ \etc.
\end{alignat*}}%
If we multiply each development by~$dx$ and add the
results, we have a series made up of the following
terms, arising from the different columns,
\begin{alignat*}{7}
&\phi a &&×{}      &&    &&    &&        &&    && \phantom{()}m\, dx\Add{,} \\
&\phi' a &&× (1  &&+ 2 &&+ 3 &&+ \dots &&+ m &&)\, (dx)^{2}\Add{,} \\
&\phi'' a &&× (1^{2} &&+ 2^{2} &&+ 3^{2} &&+ \dots &&+ m^{2} &&)\, \frac{(dx)^{3}}{2}\Add{,} \\
&\phi''' a &&× (1^{3} &&+ 2^{3} &&+ 3^{3} &&+ \dots &&+ m^{3} &&)\, \frac{(dx)^{4}}{2·3}
\quad \etc.\Add{,}
\end{alignat*}
and, as in the last example, we may represent (\PageRef{73}),
\begin{alignat*}{6}
&1  &&+ 2  &&+ 3 &&+ \dots &&+ m  &&\text{by}\quad
  \tfrac{1}{2}m^{2}(1 + \alpha)\Add{,} \displaybreak[0] \\
&1^{2} &&+ 2^{2} &&+ 3^{2} &&+ \dots &&+ m^{2} &&\text{by}\quad
  \tfrac{1}{3}m^{3}(1 + \beta)\Add{,} \displaybreak[0] \\
&1^{3} &&+ 2^{3} &&+ 3^{3} &&+ \dots &&+ m^{3}\quad &&\text{by}\quad
  \tfrac{1}{4}m^{4}(1 + \gamma)
\quad \etc.\Add{,}
\end{alignat*}
\PageSep{121}
where $\alpha$,~$\beta$,~$\gamma$, etc., diminish without limit, when $m$~is
increased without limit. If we substitute these values,
and also put $\dfrac{h}{m}$ instead of~$dx$, we have, for the
sum of the terms,
\begin{align*}
\phi a\, h + \phi' a\, \frac{h^{2}}{2} (1 + \alpha)
  &+ \phi'' a\, \frac{h^{3}}{2·3} (1 + \beta) \\
  &+ \phi''' a\, \frac{h^{4}}{2·3·4} (1 + \gamma) + \etc.
\end{align*}
which, when $m$~is increased without limit, in consequence
of which $\alpha$,~$\beta$,~etc., diminish without limit,
continually approaches to
\[
\phi a\, h + \phi' a\, \frac{h^{2}}{2}
  + \phi'' a\, \frac{h^{3}}{2·3}
  + \phi''' a\, \frac{h^{4}}{2·3·4} + \etc.\Add{,}
\]
which is the limit arising from supposing $x$ to increase
from~$a$ through $a + dx$, $a + 2\, dx$, etc., up to~$a + h$,
multiplying every value of~$\phi x$ so obtained by~$dx$, summing
the results, and decreasing~$dx$ without limit.

This is the integral of $\phi x\, dx$ from $x = a$ to $x = a + h$.
\index{Integrals!relations between differential coefficients and}%
It is evident that this series bears a great resemblance
to the development in \PageRef{21}, deprived
of its first term. Let us suppose that $\psi a$~is the function
of which $\phi a$~is the differential coefficient, that is,
that $\psi' a = \phi a$. These two functions being the same,
their differential coefficients will be the same, that is,
$\psi'' a = \phi' a$. Similarly $\psi''' a = \phi'' a$, and so on. Substituting
these, the above series becomes
\[
\psi' a\, h + \psi'' a\, \frac{h^{2}}{2}
  + \psi''' a\, \frac{h^{3}}{2·3}
  + \psi^{\text{iv}} a\, \frac{h^{4}}{2·3·4}
  + \etc.\Add{,}
\]
{\Loosen which is (\PageRef{21}) the same as $\psi(a + h) - \psi a$. That
is, the integral of $\phi x\, dx$ between the limits $a$~and~$a + h$,
is $\psi(a + h) - \psi a$, where $\psi x$~is the function, which,
\PageSep{122}
when differentiated, gives~$\phi x$. For $a + h$ we may
write~$b$, so that $\psi b - \psi a$ is the integral of~$\phi x\, dx$ from
$x = a$ to $x = b$. Or we may make the second limit indefinite
by writing~$x$ instead of~$b$, which gives $\psi x - \psi a$,
which is said to be the integral of~$\phi x\, dx$, beginning
when $x = a$, the summation being supposed to be continued
from $x = a$ until $x$~has the value which it may
be convenient to give it.}


\Subsection{Nature of Integration.}

Hence results a new branch of the inquiry, the reverse
of the Differential Calculus, the object of which
is, not to find the differential coefficient, having given
the function, but to find the function, having given
the differential coefficient. This is called the Integral
Calculus.

From the definition given, it is obvious that the
value of an integral is not to be determined, unless
we know the values of~$x$ corresponding to the beginning
and end of the summation, whose limit furnishes
the integral. We might, instead of defining the integral
in the manner above stated, have made the
word mean merely the converse of the differential coefficient;
thus, if $\phi x$~be the differential coefficient of~$\psi x$,
$\psi x$~might have been called the integral of~$\phi x\, dx$.
We should then have had to show that the integral,
thus defined, is equivalent to the limit of the summation
already explained. We have preferred bringing
the former method before the student first, as it is
most analogous to the manner in which he will deduce
integrals in questions of geometry or mechanics.
\index{Integrals!indefinite}%

With the last-mentioned definition, it is also obvious
that every function has an unlimited number of
integrals. For whatever differential coefficient~$\psi x$
\PageSep{123}
gives, $C + \psi x$ will give the same, if $C$~be a constant,
that is, not varying when $x$~varies. In this case, if $x$
become $x + h$, $C + \psi x$ becomes $C + \psi x + \psi' x\, h + \etc.$,
from which the subtraction of the original form $C + \psi x$
gives $\psi' x\, h + \etc.$; whence, by the process in \PageRef{23},
$\psi' x$~is the differential coefficient of $C + \psi' x$ as well as
of~$\psi x$. As many values, therefore, positive or negative,
as can be given to~$C$, so many different integrals
\index{Integrals!indefinite}%
can be found for~$\psi' x$; and these answer to the various
limits between which the summation in our original
definition may be made. To make this problem definite,
not only $\psi' x$ the function to be integrated, must
be given, but also that value of~$x$ from which the summation
is to begin. If this be~$a$, the integral of~$\psi' x$ is,
as before determined, $\psi x - \psi a$, and $C = -\psi a$. We
may afterwards end at any value of~$x$ which we please.
If $x = a$, $\psi x - \psi a = 0$, as is evident also from the
formation of the integral. We may thus, having given
an integral in terms of~$x$, find the value at which it
began, by equating the integral to zero, and finding
the value of~$x$. Thus, since $x^{2}$, when differentiated,
gives~$2x$, $x^{2}$~is the integral of~$2x$, beginning at $x = 0$;
and $x^{2} - 4$~is the integral beginning at~$x = 2$.

In the language of Leibnitz, an integral would be
\index{Leibnitz}%
the sum of an infinite number of infinitely small quantities,
which are the differentials or infinitely small increments
of a function. Thus, a circle being, according
to him, a rectilinear polygon of an infinite number
of infinitely small sides, the sum of these would be
the circumference of the figure. As before (\PageRefs{13}{14},
\PageNo{38}~et~seq., \PageNo{48}~et~seq.) we proceed to interpret
this inaccuracy of language. If, in a circle, we successively
describe regular polygons of $3$,~$4$, $5$,~$6$,~etc.,
sides, we may, by this means, at last attain to a polygon
\PageSep{124}
whose side shall differ from the arc of which it is
the chord, by as small a fraction, either of the chord
or arc, as we please (\PageRefs{7}{11}). That is, $A$~being
the arc, $C$~the chord, and $D$~their difference, there is
no fraction so small that $D$~cannot be made a smaller
part of~$C$. Hence, if $m$~be the number of sides of the
polygon, $mC + mD$ or $mA$ is the real circumference;
and since $mD$~is the same part of~$mC$, which $D$~is of~$C$,
$mD$~may be made as small a part of~$mC$ as we please;
so that $mC$, or the sum of all the sides of the polygon,
can be made as nearly equal to the circumference as
we please.

As in other cases, the expressions of Leibnitz are
\index{Leibnitz}%
the most convenient and the shortest, for all who can
immediately put a rational construction upon them;
this, and the fact that, good or bad, they have been,
and are, used in the works of Lagrange, Laplace,
\index{Lagrange}%
\index{Laplace}%
Euler, and many others, which the student who really
\index{Euler}%
desires to know the present state of physical science,
cannot dispense with, must be our excuse for continually
bringing before him modes of speech, which,
taken quite literally, are absurd.


\Subsection{Determination of Curvilinear Areas. The Parabola.}

We will now suppose such a part of a curve, each
\index{Curvilinear areas, determination of|EtSeq}%
\index{Parabola, the|EtSeq}%
ordinate of which is a given function of the corresponding
abscissa, as lies between two given ordinates;
for example,~$MPP'M'$. Divide the line~$MM'$
into a number of equal parts, which we may suppose
as great as we please, and construct \Fig[Figure]{10}. Let
$O$~be the origin of co-ordinates, and let $OM$, the value
of~$x$, at which we begin, be~$a$; and $OM'$, the value
at which we end, be~$b$. Though we have only divided~$MM'$
\PageSep{125}
into four equal parts in the figure, the reasoning
to which we proceed would apply equally, had we divided
it into four million of parts. The sum of the
parallelograms $Mr$,~$mr$,~$m'r''$, and~$m''R$, is less than
the area~$MPP'M'$, the value of which it is our object
to investigate, by the sum of the curvilinear triangles
$Prp$,~$pr'p'$,~$p'r''p''$, and~$p''RP'$. The sum of these triangles
is less than the sum of the parallelograms $Qr$,~$qr'$,~$q'r''$,
and~$q''R$; but these parallelograms are together
\Figure{10}
equal to the parallelogram~$q''w$, as appears by
inspection of the figure, since the base of each of the
above-mentioned parallelograms is equal to~$m''M$, or~$q''P'$,
and the altitude~$P'w$ is equal to the sum of the
altitudes of the same parallelograms. Hence the sum
of the parallelograms $Mr$,~$mr'$,~$m'r''$, and~$m''R$, differs
from the curvilinear area~$MPP'M'$ by less than the
parallelogram~$q''w$. But this last parallelogram may
be made as small as we please by sufficiently increasing
the number of parts into which $MM'$~is divided;
\PageSep{126}
for since one side of it,~$P'w$, is always less than~$P'M'$,
and the other side~$P'q''$, or~$m''M'$, is as small a part as
we please of~$MM'$ the number of square units in~$q''w$,
is the product of the number of linear units in $P'w$
and~$P'q''$, the first of which numbers being finite, and
the second as small as we please, the product is
as small as we please. Hence the curvilinear area~$MPP'M'$
is the limit towards which we continually
approach, but which we never reach, by dividing $MM'$
into a greater and greater number of equal parts, and
adding the parallelograms $Mr$,~$mr'$,~etc., so obtained.
If each of the equal parts into which $MM'$ is divided
be called~$dx$, we have $OM = a$, $Om = a + dx$, $Om' = a + 2\,dx$,
etc. And $MP$,~$mp$, $m'p'$,~etc., are the values
of the function which expresses the ordinates, corresponding
to $a$,~$a + dx$, $a + 2\, dx$,~etc., and may therefore
be represented by $\phi a$, $\phi(a + dx)$, $\phi(a + 2\, dx)$,
etc. These are the altitudes of a set of parallelograms,
the base of each of which is~$dx$; hence the
sum of their area is
\[
\phi a\, dx + \phi(a + dx)\, dx + \phi(a + 2\, dx)\, dx + \etc.,
\]
and the limit of this, to which we approach by diminishing~$dx$,
is the area required.

This limit is what we have defined to be the integral
of~$\phi x\, dx$ from $x = a$ to $x = b$; or if $\psi x$~be the
function, which, when differentiated, gives $\phi x$, it is
$\psi b - \psi a$. Hence, $y$~being the ordinate, the area included
between the axis of~$x$, any two values of~$y$, and
the portion of the curve they cut off, is $\int y\, dx$, beginning
at the one ordinate and ending at the other.

Suppose that the curve is a part of a parabola
of which $O$~is the vertex, and whose equation\footnote
  {If the student has not any acquaintance with the conic sections, he must
  nevertheless be aware that there is some curve whose abscissa and ordinate
  are connected by the equation $y^{2} = px$. This, to him, must be the definition
  of \emph{parabola}: by which word he must understand, a curve whose equation is
  $y^{2} = px$.}
is
\PageSep{127}
therefore $y^{2} = px$ where $p$~is the double ordinate which
passes through the focus. Here $y = p^{\efrac{1}{2}} x^{\efrac{1}{2}}$, and we
must find the integral of~$p^{\efrac{1}{2}} x^{\efrac{1}{2}}\, dx$, or the function
whose differential coefficient is~$p^{\efrac{1}{2}} x^{\efrac{1}{2}}$, $p^{\efrac{1}{2}}$~being a constant.
If we take the function~$cx^{n}$, $c$~being independent
of~$x$, and substitute $x + h$ for~$x$, we have for the
development $cx^{n} + cnx^{n-1}\, h + \etc$. Hence the differential
coefficient of~$cx^{n}$ is~$cnx^{n-1}$; and as $c$~and~$n$ may
be any numbers or fractions we please, we may take
them such that $cn$~shall $= p^{\efrac{1}{2}}$ and $n - 1 = \frac{1}{2}$, in which
case $n = \frac{3}{2}$ and $c = \frac{2}{3}p^{\efrac{1}{2}}$. Therefore the differential coefficient
of~$\frac{2}{3} p^{\efrac{1}{2}} x^{\efrac{3}{2}}$ is~$p^{\efrac{1}{2}} x^{\efrac{1}{2}}$, and conversely, the integral
of~$p^{\efrac{1}{2}} x^{\efrac{1}{2}}\, dx$ is~$\frac{2}{3} p^{\efrac{1}{2}} x^{\efrac{3}{2}}$.

{\Loosen The area~$MPP'M'$ of the parabola is therefore
\index{Parabola, the}%
~$\frac{2}{3} p^{\efrac{1}{2}} b^{\efrac{3}{2}} - \frac{2}{3} p^{\efrac{1}{2}} a^{\efrac{3}{2}}$. If we begin the integral at the vertex~$O$,
in which case $a = 0$, we have for the area~$OM'P'$,
$\frac{2}{3} p^{\efrac{1}{2}} b^{\efrac{3}{2}}$, where $b = OM'$. This is~$\frac{2}{3} p^{\efrac{1}{2}} b^{\efrac{1}{2}} × b$, which, since
$p^{\efrac{1}{2}} b^{\efrac{1}{2}} = M'P'$ is $\frac{2}{3}P'M' × OM'$, or two-thirds of the rectangle\footnote
  {This proposition is famous as having been discovered by Archimedes
\index{Archimedes}%
  at a time when such a step was one of no small magnitude.}
contained by $OM'$~and~$M'P'$.}


\Subsection{Method of Indivisibles.}

We may mention, in illustration of the preceding
\index{Indivisibles!method of|EtSeq}%
problem, a method of establishing the principles of
the Integral Calculus, which generally goes by the
name of the \emph{Method of Indivisibles}. A line is considered
as the sum of an infinite number of points, a
surface of an infinite number of lines, and a solid of
an infinite number of surfaces. One line twice as long
as another would be said to contain twice as many
\PageSep{128}
points, though the number of points in each is unlimited.
To this there are two objections. First, the
word infinite, in this absolute sense, really has no
\index{Infinite@\emph{Infinite}, the word}%
meaning, since it will be admitted that the mind has
no conception of a number greater than any number.
The word infinite\footnote
  {See \Title{Study of Mathematics} (Chicago: The Open Court Publishing~Co\Add{.}),
  page~123 et~seq.}
can only be justifiably used as an
abbreviation of a distinct and intelligible proposition;
for example, when we say that $a + \dfrac{1}{x}$ is equal to~$a$
when $x$~is infinite, we only mean that as $x$ is increased,
$a + \dfrac{1}{x}$~becomes nearer to~$a$, and may be made as near
to it as we please, if $x$~may be as great as we please.
The second objection is, that the notion of a line
being the sum of a number of points is not true, nor
does it approach nearer the truth as we increase the
number of points. If twenty points be taken on a
straight line, the sum of the twenty-one lines which
lie between point and point is equal to the whole line;
which cannot be if the points by themselves constitute
any part of the line, however small. Nor will the sum
of the points be a part of the line, if twenty thousand
be taken instead of twenty. There is then, in this
method, neither the rigor of geometry, nor that approach
to truth, which, in the method of Leibnitz,
\index{Leibnitz}%
may be carried to any extent we please, short of absolute
correctness. We would therefore recommend to
the student not to regard any proposition derived
from this method as true on that account; for falsehoods,
as well as truths, may be deduced from it. Indeed,
the primary notion, that the number of points
in a line is proportional to its length, is manifestly incorrect.
Suppose (\Fig{6}, \PageRef{48}) that the point~$Q$
\PageSep{129}
moves from $A$ to~$P$. It is evident that in whatever
number of points $OQ$ cuts~$AP$, it cuts~$MP$ in the same
number. But $PM$~and~$PA$ are not equal. A defender
of the system of indivisibles, if there were such a person,
\index{Indivisibles!notion of, in mechanics|EtSeq}%
would say something equivalent to supposing
that the points on the two lines are of \emph{different sizes},
which would, in fact, be an abandonment of the
method, and an adoption of the idea of Leibnitz, using
\index{Leibnitz}%
the word \emph{point} to stand for the infinitely small
\index{Point@\emph{Point}, the word}%
line.

This notion of indivisibles, or at least a way of
speaking which looks like it, prevails in many works
on mechanics. Though a point is not treated as a
length, or as any part of space whatever, it is considered
as having weight; and two points are spoken of
as having different weights. The same is said of a
line and a surface, neither of which can correctly be
supposed to possess weight. If a solid be of the same
density throughout, that is, if the weight of a cubic
inch of it be the same from whatever part it is cut, it
is plain that the weight may be found by finding the
number of cubic inches in the whole, and multiplying
this number by the weight of one cubic inch. But if
the weight of every two cubic inches is different, we
can only find the weight of the whole by the integral
calculus.

Let $AB$ (\Fig{11}) be a line possessing weight, or
\index{Points, the number of, in a straight line}%
a very thin parallelepiped of matter, which is such,
that if we were to divide it into any number of equal
parts, as in the figure, the weight of the several parts
would be different. We suppose the weight to vary
continuously, that is, if two contiguous parts of equal
length be taken, as $pq$~and~$qr$, the ratio of the weights
\PageSep{130}
of these two parts may, by taking them sufficiently
small, be as near to equality as we please.

The \emph{density} of a body is a mathematical term, which
\index{Density, continuously varying|EtSeq}%
\index{Specific gravity, continuously varying|EtSeq}%
may be explained as follows: A cubic inch of gold
weighs more than a cubic inch of water; hence gold
is \emph{denser} than water. If the first weighs $19$~times as
much as the second, gold is said to be $19$~times more
dense than water, or the density of gold is $19$~times
that of water. Hence we might define the density by
the weight of a cubic inch of the substance, but it is
usual to take, not this weight, but the proportion
which it bears to the same weight of water. Thus,
when we say the \emph{density}, or \emph{specific gravity} (these terms
are used indifferently), of cast iron is~$7.207$, we mean
\index{Iron bar continually varying in density, weight of|EtSeq}%
\index{Weight of an iron bar of which the density varies from point to point|EtSeq}%
that if any vessel of pure water were emptied and
filled with cast iron, the iron would weigh $7.207$~times
as much as the water.

If the density of a body were uniform throughout,
we might easily determine it by dividing the weight
of any bulk of the body, by the weight of an equal
bulk of water. In the same manner (\PageRef[pages]{52} et~seq.)\
we could, from our definition of velocity, determine
any uniform velocity by dividing the length described
by the time. But if the density vary continuously,
no such measure can be adopted. For if by the side
of~$AB$ (which we will suppose to be of iron) we placed
a similar body of water similarly divided, and if we
divided the weight of the part~$pq$ of iron by the weight
of the same part of water, we should get different
densities, according as the part~$pq$ is longer or shorter.
The water is supposed to be homogeneous, that is,
any part of it~$pr$, being twice the length of~$pq$, is twice
the weight of~$pq$, and so on. The iron, on the contrary,
being supposed to vary in density, the doubling
\PageSep{131}
the length gives either more or less than twice the
weight. But if we suppose $q$ to move towards~$p$, both
on the iron and the water, the limit of the ratio~$pq$ of
iron to $pq$~of water, may be chosen as a measure of
the density of~$p$, on the same principle as in \PageRefs{54}{55},
the limit of the ratio of the length described to
the time of describing it, was called the velocity. If
we call $k$ this limit, and if the weight varies continuously,
though no part~$pq$, however small, of iron,
would be exactly $k$~times the same part of water in
weight, we may nevertheless take $pq$ so small that
these weights shall be as nearly as we please in the
ratio of $k$~to~$1$.

Let us now suppose that this density, expressed
by the limiting ratio aforesaid, is always $x^{2}$ at any
\Figure{11}
point whose distance from~$A$ is $x$~feet; that is, the
density at~$q$, $2$~feet distance from~$A$, is~$4$, and so on.
Let the whole distance $AB = a$. If we divide~$a$ into
$n$~equal parts, each of which is~$dx$, so that $n\, dx = a$,
and if we call~$b$ the area of the section of the parallelepiped,
($b$~being a fraction of a square foot,) the
solid content of each of the parts will be $b\, dx$ in
cubic feet; and if $w$~be the weight of a cubic foot of
water, the weight of the same bulk of water will be~$wb\, dx$.
If the solid~$AB$ were homogeneous in the immediate
neighborhood of the point~$p$, the density being
then~$x^{2}$, would give $x^{2} × bw\, dx$ for the weight of the
same part of the substance. This is not true, but can
be brought as near to the truth as we please, by taking
$dx$ sufficiently small, or dividing~$AB$ into a sufficient
\PageSep{132}
number of parts. Hence the real weight of~$pq$
may be represented by $bwx^{2}\, dx + \alpha$, where $\alpha$~may be
made as small a part as we please of the term which
precedes it.

In the sum of any number of these terms, the sum
arising from the term~$\alpha$ diminishes without limit as
compared with the sum arising from the term~$bwx^{2}\, dx$;
for if $\alpha$~be less than the thousandth part of~$p$, $\alpha'$~less
than the thousandth part of~$p'$, etc., then $\alpha + \alpha' + \etc$.\
will be less than the thousandth part of~$p + p' + \etc.$:
which is also true of any number of quantities, and of
any fraction, however small, which each term of one
set is of its corresponding term in the other. Hence
the taking of the integral of~$bwx^{2}\, dx$ dispenses with
the necessity of considering the term~$\alpha$; for in taking
the integral, we find a limit which supposes $dx$ to
have decreased without limit, and the \emph{integral} which
would arise from~$\alpha$ has therefore diminished without
limit.

The integral of~$bwx^{2}\, dx$ is~$\frac{1}{3}bwx^{3}$, which taken from
$x = 0$ to $x = a$ is~$\frac{1}{3}bwa^{3}$. This is therefore the weight
in pounds of the bar whose length is $a$~feet, and whose
section is $b$~square feet, when the density at any point
distant by $x$~feet from the beginning is~$x^{2}$; $w$~being
the weight in pounds of a cubic foot of water.


\Subsection{Concluding Remarks on the Study of the Calculus.}

We would recommend it to the student, in pursuing
\index{Advice for studying the Calculus}%
\index{Approximate solutions in the Integral Calculus}%
\index{Rough methods of solution in the Integral Calculus}%
any problem of the Integral Calculus, never for
one moment to lose sight of the manner in which he
would do it, if a rough solution for practical purposes
only were required. Thus, if he has the area of a
curve to find, instead of merely saying that~$y$, the
ordinate, being a certain function of the abscissa~$x$,
\PageSep{133}
$\int y\, dx$ within the given limits would be the area required;
and then proceeding to the mechanical solution
of the question: let him remark that if an approximate
solution only were required, it might be
obtained by dividing the curvilinear area into a number
of four-sided figures, as in \Fig[Figure]{10}, one side of
which only is curvilinear, and embracing so small an
arc that it may, without visible error, be considered
as rectilinear. The mathematical method begins with
the same principle, investigating upon this supposition,
not the sum of these rectilinear areas, but the
limit towards which this sum approaches, as the subdivision
is rendered more minute. This limit is shown
to be that of which we are in search, since it is proved
that the error diminishes without limit, as the subdivision
is indefinitely continued.

We now leave our reader to any elementary work
which may fall in his way, having done our best to
place before him those considerations, something
equivalent to which he must turn over in his mind before
he can understand the subject. The method so
generally followed in our elementary works, of leading
the student at once into the mechanical processes
of the science, postponing entirely all other considerations,
is to many students a source of obscurity at
least, if not an absolute impediment to their progress;
since they cannot imagine what is the object of that
which they are required to do. That they shall understand
everything contained in these treatises, on
the first or second reading, we cannot promise; but
that the want of illustration and the preponderance of
\emph{technical} reasoning are the great causes of the difficulties
which students experience, is the opinion of many
\index{Advice for studying the Calculus}%
\index{Approximate solutions in the Integral Calculus}%
\index{Rough methods of solution in the Integral Calculus}%
who have had experience in teaching this subject.
\PageSep{134}
%[Blank page]
\PageSep{135}
\BackMatter

\Section[Bibliography of Standard Text-books and Works of Reference on the Calculus]
{Brief Bibliography.\protect\footnotemark}

\footnotetext{The information given regarding the works mentioned in this list is designed
  to enable the reader to select the books which are best suited to his
  needs and his purse. Where the titles do not sufficiently indicate the character
  of the books, a note or extract from the Preface has been added. The
  American prices have been supplied by Messrs.\ Lemcke \&~Buechner, 812~Broadway,
  New~York, through whom the purchases, especially of the foreign
  books, may be conveniently made.---\Ed.}

\BibSect{Standard Text-books and Treatises on
the Calculus.}

\BibSubsect{English.}

\begin{Book}
Perry, John: \Title{Calculus for Engineers.} Second edition, London
and New York: Edward Arnold. 1897. Price, \Price{7s. 6d.} (\$2.50).

\begin{Descrip}
Extract from Author's Preface: ``This book describes what has
for many years been the most important part of the regular course in
the Calculus for Mechanical and Electrical Engineering students at
the Finsbury Technical College. The students in October knew only
the most elementary mathematics, many of them did not know the
Binomial Theorem, or the definition of the sine of an angle. In July
they had not only done the work of this book, but their knowledge
was of a practical kind, ready for use in any such engineering problems
as I give here.''

Especially good in the character and number of practical examples
given.
\end{Descrip}
\end{Book}

\begin{Book}
Lamb, Horace: \Title{Infinitesimal Calculus.} New York: The Macmillan
Co. 1898. Price,~\$3.00.

\begin{Descrip}
Extract from Author's Preface: ``This book attempts to teach
those portions of the Calculus which are of primary importance in
the application to such subjects as Physics and Engineering\dots.
Stress is laid on fundamental principles\dots. Considerable attention
has been paid to the logic of the subject.''
\end{Descrip}
\end{Book}
\PageSep{136}

\begin{Book}
Edwards, Joseph: \Title{An Elementary Treatise on the Differential
Calculus.} Second edition, revised. 8vo,~cloth. New York
and London: The Macmillan~Co. 1892. Price, \$3.50.---%
\Title{Differential Calculus for Beginners.} 8vo,~cloth. 1893.\Chg{ }{---}\Title{The
Integral Calculus for Beginners.} 8vo,~cloth. (Same Publishers.)
Price, \$1.10~each.
\end{Book}

\begin{Book}
Byerly, William E.: \Title{Elements of the Differential Calculus.} Boston:
Ginn \&~Co. Price, \$2.15.---\Title{Elements of the Integral
Calculus.} (Same Publishers.) Price,~\$2.15.
\end{Book}

\begin{Book}
Rice, J.~M., and Johnson, W.~W.: \Title{An Elementary Treatise on
the Differential Calculus Founded on the Method of Rates
or Fluxions.} New~York: John Wiley \&~Sons. 8vo. 1884.
Price, \$3.50. Abridged edition, 1889. Price,~\$1.50.
\end{Book}

\begin{Book}
Johnson, W.~W.: \Title{Elementary Treatise on the Integral Calculus
Founded on the Method of Rates or Fluxions.} 8vo,~cloth.
New~York: John Wiley \&~Sons. 1885. Price,~\$1.50.
\end{Book}

\begin{Book}
Greenhill, A.~G.: \Title{Differential and Integral Calculus.} With applications.
8vo,~cloth. Second edition. New~York and London:
The Macmillan~Co. 1891. Price, \Price{9s.}~(\$2.60).
\end{Book}

\begin{Book}
Price: \Title{Infinitesimal Calculus.} Four Vols. 1857--65. Out of
print and very scarce. Obtainable for about~\$27.00.
\end{Book}

\begin{Book}
Smith, William Benjamin: \Title{Infinitesimal Analysis.} Vol.~I., Elementary:
Real Variables. New~York and London: The Macmillan~Co.
1898. Price,~\$3.25.

\begin{Descrip}
``The aim has been, within a prescribed expense of time and
energy to penetrate as far as possible, and in as many directions, into
the subject in hand,---that the student should attain as wide knowledge
of the matter, as full comprehension of the methods, and as clear
consciousness of the spirit and power of analysis as the nature of the
case would admit.''---From Author's Preface.
\end{Descrip}
\end{Book}

\begin{Book}
Todhunter, Isaac: \Title{A Treatise on the Differential Calculus.} London
and New~York: The Macmillan~Co. Price, \Price{10s. 6d.}
(\$2.60). \Title{A Treatise on the Integral Calculus.} (Same publishers.)
Price, \Price{10s. 6d.} (\$2\Chg{ }{.}60).

\begin{Descrip}
Todhunter's text-books were, until recently, the most widely used
in England. His works on the Calculus still retain their standard
character, as general manuals.
\end{Descrip}
\end{Book}
\PageSep{137}

\begin{Book}
Williamson: \Title{Differential and Integral Calculus.} London and
New~York: Longmans, Green, \&~Co. 1872--1874. Two~Vols.
Price, \$3.50~each.
\end{Book}

\begin{Book}
De~Morgan, Augustus: \Title{Differential and Integral Calculus.} London:
Society for the Diffusion of Useful Knowledge. 1842.
Out of print. About~\$6.40.

\begin{Descrip}
The most extensive and complete work in English. ``The object
has been to contain within the prescribed limits, the whole of the
students' course from the confines of elementary algebra and trigonometry,
to the entrance of the highest works on mathematical physics''
(Author's Preface). Few examples. In typography, and general
arrangement of material, inferior to the best recent works. Valuable
for collateral study, and for its philosophical spirit.
\end{Descrip}
\end{Book}


\BibSubsect{French.}

\begin{Book}
Sturm: \Title{Cours d'analyse de l'École Polytechnique.} 10.~édition,
revue et corrigé par E.~Prouhet, et augmentée de~la théorie
élémentaire des fonctions elliptiques, par H.~Laurent. 2~volumes
in---8. Paris: Gauthier-Villars et~fils. 1895. Bound,
16~fr.\ 50~c. \$4.95.

\begin{Descrip}
One of the most widely used of text-books. First published in
1857. The new tenth edition has been thoroughly revised and brought
down to date. The exercises, while not numerous, are sufficient, those
which accompany the additions and complementary chapters of M.~De~Saint
Germain having been taken from the Collection of M.~Tisserand,
mentioned below.
\end{Descrip}
\end{Book}

\begin{Book}
Duhamel: \Title{Éléments de calcul infinitésimal.} 4.~Edition, revue et
annotée par J.~Bertrand. 2~volumes in---8; avec planches.
Paris: Gauthier-Villars et~fils. 1886. 15~fr. \$4.50.

\begin{Descrip}
The first edition was published between 1840 and 1841. ``Cordially
recommended to teachers and students'' by De~Morgan. Duhamel
paid great attention to the philosophy and logic of the mathematical
sciences, and the student may also be referred in this connexion to
his \Title{Méthodes dans les sciences de raisonnement}. 5~volumes. Paris:
Gauthier-Villars et~fils. Price, 25.50~francs. \$7.65.
\end{Descrip}
\end{Book}

\begin{Book}
Lacroix, S.-F.: \Title{Traité élémentaire de calcul différentiel et de
calcul intégral.} 9.~Edition, revue et augmentée de notes par
Hermite et Serret. 2~vols. Paris: Gauthier-Villars et~fils.
1881. 15~fr. \$4.50.

\begin{Descrip}
A very old work. The first edition was published in 1797. It was
the standard treatise during the early part of the century, and has
been kept revised by competent hands.
\end{Descrip}
\end{Book}
\PageSep{138}

\begin{Book}
Appell, P.; \Title{Éléments d'analyse mathématique.} À l'usage des
ingénieurs et dés physiciens. Cours professé à l'École Centrale
des Arts et Manufactures. 1~vol.\ in---8, 720~pages, avec
figures, cartonné à l'anglaise. Paris: Georges Carré \&~C.
Naud. 1899. Price, 24~francs.\ \$7.20.
\end{Book}

\begin{Book}
Boussinesq, J.: \Title{Cours d'analyse infinitésimal.} À l'usage des
personnes qui étudient cette science en vue de ses applications
mécaniques et physiques, 2~vols., grand in\Chg{-}{---}8, avec figures.
Tome~I\@. Calcul différentiel. Paris, 1887. 17~fr.\ (\$5.10).
Tome~II\@. Calcul intégral. Paris: Gauthier-Villars et~fils.
1890. 23~fr.\ 50~c.\ (\$7.05).
\end{Book}

\begin{Book}
Hermite, Ch.: \Title{Cours d'analyse de l'École Polytechnique.} 2~vols.
Vol.~I\@. Paris: Gauthier-Villars et~fils. 1897.

\begin{Descrip}
A new edition of Vol.~I. is in preparation (1899). Vol.~II. has not
yet appeared.
\end{Descrip}
\end{Book}

\begin{Book}
Jordan, Camille: \Title{Cours d'analyse de l'École Polytechnique.} 3~volumes.
2.~édition. Paris: Gauthier-Villars et~fils. 1893--1898.
51~fr.\ \$14.70.

\begin{Descrip}
Very comprehensive on the theoretical side. Enters deeply into
the metaphysical aspects of the subject.
\end{Descrip}
\end{Book}

\begin{Book}
Laurent, H.: \Title{Traité d'analyse.} 7~vols in---8. Paris: Gauthier-Villars
et~fils. 1885--1891. 73~fr.\ \$21.90.

\begin{Descrip}
The most extensive existing treatise on the Calculus. A general
handbook and work of reference for the results contained in the
more special works and memoirs.
\end{Descrip}
\end{Book}

\begin{Book}
Picard, Émile: \Title{Traité d'analyse.} 4~volumes grand in\Chg{-}{---}8. Paris:
Gauthier-Villars et~fils. 1891. 15~fr.\ each. Vols.~I.--III.,
\$14.40. Vol.~IV. has not yet appeared.

\begin{Descrip}
An advanced treatise on the Integral Calculus and the theory of
differential equations. Presupposes a knowledge of the Differential
Calculus.
\end{Descrip}
\end{Book}

\begin{Book}
Serret, J.-A.: \Title{Cours de calcul différentiel et intégral.} 4.~edition,
augmentée d'une note sur les fonctions elliptiques, par
Ch.~Hermite. 2~forts volumes in---8. Paris: Gauthier-Villars
et~fils. 1894. 25~fr.\ \$7.50.

\begin{Descrip}
A good German translation of this work by Axel Harnack has
passed through its second edition (Leipsic: Teubner, 1885 and 1897).
\end{Descrip}
\end{Book}
\PageSep{139}

\begin{Book}
Hoüel, J.: \Title{Cours de calcul infinitésimal.} 4~beaux volumes grand
in---8, avec figures. Paris: Gauthier-Villars et~fils. 1878--1879--1880--1881.
50~fr.\ \$15.00.
\end{Book}

\begin{Book}
Bertrand, J.: \Title{Traité de calcul différentiel et de calcul intégral.}
(1)~Calcul différentiel. Paris: Gauthier-Villars et fils. 1864.
Scarce. About \$48.00. (2)~Calcul intégral (Intégrales définies
et indéfinies). Paris, 1870. Scarce. About \$24.00.
\end{Book}

\begin{Book}
Boucharlat, J.-L.: \Title{Éléments de calcul différentiel et de calcul
intégral.} 9.~édition, revue et annotée par H.~Laurent. Paris:
Gauthier-Villars et~fils. 1891. 8~fr.\ \$2.40.
\end{Book}

\begin{Book}
Moigno: \Title{Leçons de calcul différentiel et de calcul intégral\Typo{,}{.}} 2~vols.,
Paris, 1840--1844. Scarce. About \$9.60.
\end{Book}

\begin{Book}
Navier: \Title{Leçons d'analyse de l'École Polytechnique.} Paris, 1840.
2nd~ed. 1856. Out of print. About \$3.60.

\begin{Descrip}
An able and practical work. Very popular in its day. The typical
course of the \textit{École Polytechnique}, and the basis of several of the treatises
that followed, including that of Sturm. Also much used in its
German translation.
\end{Descrip}
\end{Book}

\begin{Book}
Cournot: \Title{Théorie des fonctions et du calcul infinitésimal.} 2~vols.
Paris, 1841. 2nd~ed. 1856--1858. Out of print, and
scarce. About \$3.00.

\begin{Descrip}
The first edition (1841) was ``cordially recommended to teachers
and students'' by De~Morgan. Cournot was especially strong on the
philosophical side. He examined the foundations of many sciences
and developed original views on the theory of knowledge, which are
little known but have been largely drawn from by other philosophers.
\end{Descrip}
\end{Book}

\begin{Book}
Cauchy, A.: \Title{\OE{}uvres complètes.} Tome~III: \Title{Cours d'analyse
de l'École Polytechnique.} Tome~IV: \Title{Résumé des leçons
données à l'École Polytechnique sur le calcul infinitésimal.
Leçons sur le calcul différentiel.} Tome~V: \Title{Leçons sur les
applications du calcul infinitésimal à la géométrie.} Paris:
Gauthier-Villars et~fils, 1885--1897. 25~fr.\ each. \$9.50~each.

\begin{Descrip}
The works of Cauchy, as well as those of Lagrange, which follow,
are mentioned for their high historical and educational importance.
\end{Descrip}
\end{Book}

\begin{Book}
Lagrange, J.~L.: \Title{\OE{}uvres complètes.} Tome~IX: \Title{Théorie des fonctions
analytiques.} Tome~X.: \Title{Leçons sur le calcul des fonctions.}
\PageSep{140}
Paris: Gauthier-Villars et~fils, 1881--1884. 18~fr.\ per
volume. \$5.40 per~volume.

\begin{Descrip}
``The same power of abstraction and facility of treatment which
signalise these works are nowhere to be met with in the prior or subsequent
history of the subject. In addition, they are replete with the
profoundest aperçus into the history of the development of analytical
truths,---aperçus which could have come only from a man who combined
superior creative endowment with exact and comprehensive
knowledge of the facts. In the remarks woven into the body of the
text will be found what is virtually a detailed history of the subject,
and one which is not to be had elsewhere, least of all in diffuse histories
of mathematics. The student, thus, not only learns in these
works how to think, but also discovers how people actually have
thought, and what are the ways which human instinct and reason
have pursued in the different individuals who have participated in
the elaboration of the science.''---(E.~Dühring.)
\end{Descrip}
\end{Book}

\begin{Book}
Euler, L.:

\begin{Descrip}
The Latin treatises of Euler are also to be mentioned in this connexion,
for the benefit of those who wish to pursue the history of the
text-book making of this subject to its fountain-head. They are the
\Title{Differential Calculus} (St.~Petersburg, 1755), the \Title{Integral Calculus} (3~vols.,
St.~Petersburg, 1768--1770), and the \Title{Introduction to the Infinitesimal
Analysis} (2~vols., Lausanne, 1748). Of the last-mentioned work
an old French translation by Labey exists (Paris: Gauthier-Villars),
and a new German translation (of Vol.~I. only) by Maser (Berlin:
Julius Springer, 1885). Of the first-mentioned treatises on the Calculus
proper there exist two old German translations, which are not
difficult to obtain.
\end{Descrip}
\end{Book}


\BibSubsect{German.}

\begin{Book}
Harnack, Dr.\ Axel: \Title{Elemente der Differential- und Integralrechnung.}
Zur Einführung in das Studium dargestellt. Leipzig:
Teubner, 1881. M.~7.60. Bound, \$2.80. (English translation.
London: Williams \&~Norgate. 1891.)
\end{Book}

\begin{Book}
Junker, Dr.\ Friedrich: \Title{Höhere Analysis.} I.~\Title{Differentialrechnung.}
Mit 63~Figuren. II.~\Title{Integralrechnung.} Leipzig:
G.~J. Göschen'sche Verlagshandlung. 1898--1899. 80~pf.\ each.
30~cents each.

\begin{Descrip}
These books are marvellously cheap, and very concise. They
contain no examples. Pocket-size.
\end{Descrip}
\end{Book}

\begin{Book}
Autenheimer, F.: \Title{Elementarbuch der Differential- und Integralrechnung
mit zahlreichen Anwendungen aus der Analysis,
Geometrie, Mechanik, Physik etc.} Für höhere Lehranstalten
\PageSep{141}
und den Selbstunterricht. 4te~verbesserte Auflage. Weimar:
Bernhard Friedrich Voigt. 1895.

\begin{Descrip}
As indicated by its title, this book is specially rich in practical
applications.
\end{Descrip}
\end{Book}

\begin{Book}
Stegemann: \Title{Grundriss der Differential- und Integralrechnung},
8te~Auflage, herausgegeben von Kiepert. Hannover: Helwing,
1897. Two volumes, 26~marks. Two volumes, bound,
\$8.50.

\begin{Descrip}
This work was highly recommended by Prof.\ Felix Klein at the
Evanston Colloquium in~1893.
\end{Descrip}
\end{Book}

\begin{Book}
Schlömilch: \Title{Compendium der höheren Analysis.} Fifth edition,
1881. Two volumes, \$6.80.

\begin{Descrip}
Schlömilch's text-books have been very successful. The present
work was long the standard manual.
\end{Descrip}
\end{Book}

\begin{Book}
Stolz, Dr.\ Otto: \Title{Grundzüge der Differential- und Integralrechnung.}
In 2~Theilen. I.~Theil. Reelle Veränderliche und
Functionen. (460~S.) 1893. M.~8. II.~Complexe Veränderliche
und Functionen. (338~S.) Leipzig: Teubner. 1896.
M.~8. Two volumes, \$6.00.

\begin{Descrip}
A supplementary 3rd part entitled \Title{Die Lehre von den Doppelintegralen}
has just been published (1899). Based on the works of J.~Tannery,
Peano, and Dini.
\end{Descrip}
\end{Book}

\begin{Book}
Lipschitz, R.: \Title{Lehrbuch der Analysis.} 1877--1880. Two volumes,
bound, \$12.30.

\begin{Descrip}
Specially good on the theoretical side.
\end{Descrip}
\end{Book}


\BibSubsect{Collections of Examples and Illustrations.}

\begin{Book}
Byerly, W.~E.: \Title{Problems in Differential Calculus.} Supplementary
to a Treatise on Differential Calculus. Boston: Ginn \&~Co.
75~cents.
\end{Book}

\begin{Book}
Gregory: \Title{Examples on the Differential and Integral Calculus.}
1841. Second edition. 1846. Out of print. About \$6.40.
\end{Book}

\begin{Book}
Frenet: \Title{Recueil d'exercises sur le calcul infinitésimal.} 5.~édition,
augmentée d'un appendice, par H.~Laurent. Paris:
Gauthier-Villars et~fils. 1891. 8~fr.\ \$2.40.
\end{Book}
\PageSep{142}

\begin{Book}
Tisserand, F.: \Title{Recueil complémentaire d'exercises sur le calcul
infinitésimal.} Second edition. Paris: Gauthier-Villars et~fils
1896.

\begin{Descrip}
Complementary to Frenet.
\end{Descrip}
\end{Book}

\begin{Book}
Laisant, C.~A.: \Title{Recueil de problèmes de mathématiques.} Tome~VII\@.
Calcul infinitésimal et calcul des fonctions. Mécanique.
Astronomie. (Announced for publication.) Paris: Gauthier-Villars
et~fils.
\end{Book}

\begin{Book}
Schlömilch, Dr.\ Oscar: \Title{Uebungsbuch zum Studium der höheren
Analysis.} I.~Theil. Aufgaben aus der Differentialrechnung.
4te~Auflage. (336~S.) 1887. M.~6. II.~Aufgaben aus der
Integralrechnung. 3te~Auflage. (384~S.) Leipzig: Teubner,
1882. M.~7.60. Both volumes, bound, \$7.60.
\end{Book}

\begin{Book}
Sohncke, L.~A.: \Title{Sammlung von Aufgaben aus der Differential-
und Integralrechnung.} Herausgegeben von Heis. Two volumes,
in---8. Bound, \$3.00.
\end{Book}

\begin{Book}
Fuhrmann, Dr.\ Arwed: \Title{Anwendungen der Infinitesimalrechnung
in den Naturwissenschaften, im Hochbau und in der
Technik.} Lehrbuch und Aufgabensammlung. In sechs Theilen,
von denen jeder ein selbstständiges Ganzes bildet. Theil~I.
Naturwissenschaftliche Anwendungen der Differentialrechnung.
Theil~II. Naturwissenschaftliche Anwendungen der
Integralrechnung. Berlin: Verlag von Ernst \&~Korn. 1888--1890.
Vol.~I., Cloth, \$1.35. Vol.~II., Cloth, \$2.20.
\end{Book}
\PageSep{143}


\printindex

\iffalse
%INDEX.

Accelerated motion 57, 60

Accelerating force 62

Advice for studying the Calculus 132, 133

Angle, unit employed in measuring an#Angle 51

Approximate solutions in the Integral Calculus 132, 133

Arc and its chord, a continuously decreasing|EtSeq#Arc 7, 39 % et seq.

Archimedes 127

Astronomical ephemeris 76

Calculus, notation of 25

Calculus, notation of|EtSeq 79 % et seq.

Circle, equation of|EtSeq 31 % et seq.

Circle cut by straight line, investigated|EtSeq 31 % et seq.

Coefficients, differential|EtSeq 22 % et seq.

Coefficients, differential 38, 55, 82, 88, 96, 100, 112

Complete Differential Coefficients 96

Constants 14

Contiguous values 112

Continuous quantities|EtSeq 7 % et seq.

Continuous quantities 53

Co-ordinates 30

Curve, magnified 40

Curvilinear areas, determination of|EtSeq 124 % et seq.

Density, continuously varying|EtSeq 130 % et seq.

Derivatives 19, 21, 22

Derived Functions|EtSeq 19 % et seq.

Derived Functions 21

Differences
  arithmetical 4
  of increments 26
  calculus of 89

Differential coefficients|EtSeq 22 % et seq.,

Differential coefficients 38, 55, 82
  as the index of the change of a function 112
  of higher orders 88

Differentials
  partial|EtSeq 78 % et seq.;
  total|EtSeq 78 % et seq.

Differentiation
  of the common functions 85, 86
  successive|EtSeq 88 % et seq.;
  implicit|EtSeq 94 % et seq.;
  of complicated functions|EtSeq 100 % et seq.

Direct function 97

Direction 36

Equality 4

Equations, solution of 77

Equidistant values 104

Euler 27, 124

Errors, in the valuation of quantities 75, 84

Explicit functions 107

Falling bodies 56

Finite differences|EtSeq 88 % et seq.

Fluxions 11, 60, 112

Force 61-63

Functions
  definition of|EtSeq 14 % et seq.;
  derived|EtSeq 19 % et seq.,
  derived 21
  direct and indirect 97
  implicit and explicit 107, 108
  inverse|EtSeq 102 % et seq.;
  of several variables|EtSeq 78 % et seq.;
  recapitulation of results in the theory of 74

Generally@\emph{Generally}, the word 16

Implicit
  differentiation|EtSeq 94 % et seq.;
  function 107, 108

Impulse 60

Increase without limit|EtSeq 5, 65 % et seq.

Increment 16, 113
\PageSep{144}

Independent variables 106

Indirect function 97

Indivisibles
  method of|EtSeq 127 % et seq.;
  notion of, in mechanics|EtSeq 129 % et seq.

Infinite@\emph{Infinite}, the word#Infinite 128

Infinitely small, the notion of#Infinitely 12, 49, 59, 83

Infinitely small, the notion of|EtSeq#Infinitely 38 % et seq.,

Infinity, orders of|EtSeq 42 % et seq.

Integral Calculus 73
  notation of 119

Integral Calculus|EtSeq 115 % et seq.

Integrals
  definition of|EtSeq 119 % et seq.;
  relations between differential coefficients and 121
  indefinite 122, 123

Intersections, limit of|EtSeq 46 % et seq.

Inverse functions|EtSeq 102 % et seq.

Iron bar continually varying in density, weight of|EtSeq#Iron 130 % et seq.

Ladder against wall|EtSeq 45 % et seq.

Lagrange 124

Laplace 124

Leibnitz 11, 13, 38, 42, 48, 59, 60, 83, 123, 124, 128, 129

Limit of intersections|EtSeq 46 % et seq.

Limits|EtSeq 26 % et seq.

Limiting ratios|EtSeq 65 % et seq.

Limiting ratios 81

Logarithms 20, 38, 86, 87

Logarithms|EtSeq 112 % et seq.

Magnified curve 40

Motion
  accelerated 60
  simple harmonic 57

Newton 11, 60

Notation
  of the Differential Calculus 25
  of the Differential Calculus|EtSeq 79 % et seq.
  of the Integral Calculus 119

Orders, differential coefficients of higher 88

Orders of infinity|EtSeq 42 % et seq.

Parabola, the#Parabola 30, 127

Parabola, the|EtSeq#Parabola 124 % et seq.

Partial
  differentials|EtSeq 78 % et seq.;
  differential coefficients 96

Point@\emph{Point}, the word#Point 129

Points, the number of, in a straight line 129

Polygon 38

Proportion|EtSeq 2 % et seq.

Quantities, continuous|EtSeq 7 % et seq.

Quantities, continuous 53

Ratio
  defined|EtSeq 2 % et seq.;
  of two increments 87

Ratios, limiting|EtSeq 65 % et seq.

Ratios, limiting 81

Rough methods of solution in the Integral Calculus 132, 133

Series|EtSeq 15, 24 % et seq.

Signs|EtSeq 31 % et seq.

Simple harmonic motion 57

Sines 87

Singular values 16

Small, has no precise meaning 12

Specific gravity, continuously varying|EtSeq 130 % et seq.

Successive differentiation|EtSeq 88 % et seq.

Sun's longitude 76

Tangent 37, 38, 40

Taylor's Theorem|EtSeq 15, 19 % et seq.

Time, idea of#Time 4

Time, idea of|EtSeq#Time 110 % et seq.

Total
  differential coefficient 100
  differentials|EtSeq 78 % et seq.;
  variations 95

Transit instrument 84

Uniformly accelerated 57, 60

Values
  contiguous 112
  equidistant 104

Variables
  independent and dependent 14, 15, 106
  functions of several|EtSeq 78 % et seq.

Variations, total#Variations 95

Velocity
  linear|EtSeq 52 % et seq.
  linear 111
  angular 59

Weight of an iron bar of which the density varies from point to point|EtSeq#Iron 130 % et seq.
\fi
\PageSep{145}

\iffalse
%[** TN: Catalog text has been (lightly) proofread, but not marked up in LaTeX]

CATALOGUE OF PUBLICATIONS
OF THE
OPEN COURT PUBLISHING CO.

COPE, E. D.

THE PRIMARY FACTORS OF ORGANIC EVOLUTION.

121 cuts. Pp. xvi, 547. Cloth, \$2.00 net (10s.).


MULLER, F. MAX.

THREE INTRODUCTORY LECTURES ON THE SCIENCE OF
THOUGHT.
128 pages. Cloth, 75c (3s. 6d.).

THREE LECTURES ON THE SCIENCE OF LANGUAGE.
112 pages. 2nd Edition. Cloth, 75c (3s. 6d.).


ROMANES, GEORGE JOHN.

DARWIN AND AFTER DARWIN.

An Exposition of the Darwinian Theory and a Discussion of Post-Darwinian
Questions. Three Vols., \$4.00 net. Singly, as follows:

1. The Darwinian Theory. 460 pages. 125 illustrations. Cloth, \$2.00.

2. Post-Darwinian Questions. Heredity and Utility. Pp. 338. \$1.50.

3. Post-Darwinian Questions. Isolation and Physiological Selection.
Pp. 181. \$1.00.

AN EXAMINATION OF WEISMANNISM.
236 pages. Cloth, \$1.00 net.

THOUGHTS ON RELIGION.
Edited by Charles Gore, M. A., Canon of Westminster. Third Edition,
Pages, 184. Cloth, gilt top, \$1.25 net.


SHUTE. DR. D. KERFOOT.

FIRST BOOK IN ORGANIC EVOLUTION.

Colored plates, and numerous diagrams. (In Preparation.)


MACH, ERNST.

THE SCIENCE OF MECHANICS.

A Critical and Historical Exposition of its Principles. Translated
by T. J. McCormack. 250 cuts. 534 pages. 1/2 in., gilt top. \$2.50 (12s. 6d.).

POPULAR SCIENTIFIC LECTURES.
Third Edition. 415 pages. 59 cuts. Cloth, gilt top. Net, \$1.50 (7s. 6d.).

THE ANALYSIS OF THE SENSATIONS.
Pp. 208. 37 cuts. Cloth, \$1.25 net (6s. 6d.).


LAGRANGE. J. L.

LECTURES ON ELEMENTARY MATHEMATICS.
With portrait of the author. Pp. 172. Price, \$1.00 net (5s.).


DE MORGAN, AUGUSTUS.

ON THE STUDY AND DIFFICULTIES OF MATHEMATICS.
New Reprint edition with notes. Pp. viii+288. Cloth, \$1.25 net (5s.).

ELEMENTARY ILLUSTRATIONS OF THE DIFFERENTIAL AND
INTEGRAL CALCULUS.
New reprint edition. Price, \$1.00 (5s.).


SCHUBERT, HERMANN.

MATHEMATICAL ESSAYS AND RECREATIONS.
Pp. 149. Cuts, 37. Cloth, 75c net (3s. 6d.).


HUC AND GABET, MM.

TRAVELS IN TARTARY, THIBET AND CHINA.
(1844--1846.) Translated from the French by W. Hazlitt. Illustrated
with 100 engravings on wood. 2 vols. Pp. 28+660. Cl., \$2.00 (10s.).
% --- Page

CARUS, PAUL.

THE ETHICAL PROBLEM.
Second edition, revised and greatly enlarged. 351 pages. Cloth, \$1.25
(6s. 6d.).

FUNDAMENTAL PROBLEMS.
Second edition, enlarged and revised. 372 pp. Cl., \$1.50 (7s. 6d.).

HOMILIES OF SCIENCE.
317 pages. Cloth, Gilt Top, \$1.50 (7s. 6d.).

THE IDEA OF GOD.
Fourth edition. 32 pages. Paper, 15c (9d.).

THE SOUL OF MAN.
With 152 cuts and diagrams. 458 pages. Cloth, \$3.00 (15s.).

TRUTH IN FICTION. Twelve Tales with a Moral.
White and gold binding, gilt edges. Pp. 111. \$1.00 (5s.).

THE RELIGION OF SCIENCE.
Second, extra edition. Pp. 103. Price, 50c net (2s. 6d.).

PRIMER OF PHILOSOPHY.
240 pages. Second Edition. Cloth, \$1.00 (5s.).

THE GOSPEL OF BUDDHA. According to Old Records.
Fifth Edition. Pp. 275. Cloth, \$1.00 (5s.). In German, \$1.25 (6s. 6d.).

BUDDHISM AND ITS CHRISTIAN CRITICS.
Pages, 311. Cloth, \$1.25 (6s. 6d.).

KARMA. A Story of Early Buddhism.
Illustrated by Japanese artists. Crêpe paper, 75c (3s. 6d.).

NIRVANA: A Story of Buddhist Psychology.
Japanese edition, like Karma. \$1.00 (4s. 6d.).

LAO-TZE'S TAO TEH-KING.
Chinese-English. With introduction, transliteration. Notes, etc. Pp.
360. Cloth, \$3.00 (15s.).


CORNILL, CARL HEINRICH.

THE PROPHETS OF ISRAEL.
Popular Sketches from Old Testament History. Pp., 200. Cloth, \$1.00
net (5s.).

HISTORY OF THE PEOPLE OF ISRAEL.
Pp. vi + 325. Cloth, \$1.50 (7s. 6d.).


POWELL, J. W.

TRUTH AND ERROR; or, the Science of Intellection.
Pp. 423. Cloth, \$1.75 (7s. 6d.).


RIBOT. TH.

THE PSYCHOLOGY OF ATTENTION.

THE DISEASES OF PERSONALITY.

THE DISEASES OF THE WILL.
Authorised translations. Cloth, 75 cents each (3s. 6d.). Full set, cloth,
\$1.75 net (9s.).

EVOLUTION OF GENERAL IDEAS.
Pp. 231. Cloth, \$1.25 net (6s. 6d.).


WAGNER, RICHARD.

A PILGRIMAGE TO BEETHOVEN.
A Story. With portrait of Beethoven. Pp. 40. Boards, 50c net (2s. 6d.)


HUTCHINSON, WOODS.

THE GOSPEL ACCORDING TO DARWIN.
Pp. xii+241. Price, \$1.50 (6s.).


FREYTAG, GUSTAV.

THE LOST MANUSCRIPT. A Novel.

2 vols. 953 pages. Extra cloth, \$4.00 (21s). One vol., cl., \$1.00 (5s.).

MARTIN LUTHER.
Illustrated. Pp.130. Cloth, \$1.00 net (5s.).
% --- Page

TRUMBULL, M. M.

THE FREE TRADE STRUGGLE IN ENGLAND.
Second Edition. 296 pages. Cloth, 75c (3s. 6d.).

WHEELBARROW: Articles and Discussions on the Labor Question
With portrait of the author. 303 pages. Cloth, \$1.00 (5s.).


GOETHE AND SCHILLER'S XENIONS.
Translated by Paul Carus. Album form. Pp. 162. Cl., \$1.00 (5s.).


OLDENBERG. H.

ANCIENT INDIA: ITS LANGUAGE AND RELIGIONS.
Pp. 100. Cloth, 50c net (2s. 6d.).


CONWAY, DR. MONCURE DANIEL.

SOLOMON, AND SOLOMONIC LITERATURE.
Pp. 243. Cloth, \$1.50 net (6s. 6d.).


GARBE, RICHARD.

THE REDEMPTION OF THE BRAHMAN. A Tale of Hindu Life.
Laid paper. Gilt top. 96 pages. Price, 75c (3s. 6d.).

THE PHILOSOPHY OF ANCIENT INDIA.
Pp. 89. Cloth, 50c net (2s. 6d.).


HUEPPE. FERDINAND.

THE PRINCIPLES OF BACTERIOLOGY.
28 Woodcuts. Pp. x+467. Price, \$1.75 net (9s.).


LÉVY-BRUHL, PROF. L.

HISTORY OF MODERN PHILOSOPHY IN FRANCE.
With 23 Photogravure and Half-Tone Portraits of the Chief French
Philosophers. Handsomely bound. (In the press.)


TOPINARD, DR. PAUL.

SCIENCE AND FAITH, or Man as an Animal and Man as a Member
of Society.
Pp. 374. Cloth, \$1.50 net (7s. 6d.).


BINET, ALFRED.

THE PSYCHOLOGY OF REASONING.
Pp. 193. Cloth, 75c (3s. 6d.).

THE PSYCHIC LIFE OF MICRO-ORGANISMS.
Pp. 135. Cloth, 75 cents.

ON DOUBLE CONSCIOUSNESS.
See No. 8, Religion of Science Library.


THE OPEN COURT.
A Monthly Magazine Devoted to the Science of Religion, the Religion of
Science, and the Extension of the Religious Parliament Idea.
Terms: \$1.00 a year; 5s. 6d. to foreign countries in the Postal Union.
Single Copies, 10 cents (6d.).

THE MONIST.
A Quarterly Magazine of Philosophy and Science.
Per copy, 50 cents; Yearly, \$2.00. In England and all countries in
U.P.U. per copy, 2s. 6d.: Yearly, 9s. 6d.


CHICAGO:

THE OPEN COURT PUBLISHING CO.,
Monon Building, 324 Dearborn St.
LONDON: Kegan Paul, Trench, Trübner & Company.
% --- Page


The Religion of Science Library,

A collection of bi-monthly publications, most of which are reprints of
books published by The Open Court Publishing Company. Yearly, \$1.50.
Separate copies according to prices quoted. The books are printed upon
good paper, from large type.

The Religion of Science Library, by its extraordinarily reasonable price,
will place a large number of valuable books within the reach of all readers.

The following have already appeared in the series:

No. 1. The Religion of Science. By Paul Carus. 25c (1s. 6d).

2. Three Introductory Lectures on the Science of Thought. By F. Max
Müller. 25c (1s. 6d.).

3. Three Lectures on the Science of Language. F. Max Müller. 25c (1s. 6d.)

4. The Diseases of Personality. By Th. Ribot. 25c (1s. 6d.).

5. The Psychology of Attention. By Th. Ribot. 25c (1s. 6d.).

6. The Psychic Life of Micro-Organisms. By Alfred Binet. 25c (1s. 6d.)

7. The Nature of the State. By Paul Carus. 15c (9d.).

8. On Double Consciousness. By Alfred Binet. 15c (9d.).

9. Fundamental Problems. By Paul Carus. 50c (2s. 6d.).

10. The Diseases of the Will. By Th. Ribot. 25c (1s. 6d.).

11. The Origin of Language. By Ludwig Noire. 15c (9d.).

12. The Free Trade Struggle in England. M. M. Trumbull. 25c (1s. 6d.).

13. Wheelbarrow on the Labor Question. By M. M. Trumbull. 35c (2s.).

14. The Gospel of Buddha By Paul Carus. 35c (2s.).

15. The Primer of Philosophy. By Paul Carus. 25c (1s. 6d.).

16. On Memory, and The Specific Energies of the Nervous System. By Prof.
Ewald Hering. 15c (9d.).

17. The Redemption of the Brahman. A Tale of Hindu Life. By Richard
Garbe. 25c (1s. 6d.).

18. An Examination of Weismannism. By G. J. Romanes. 35c (2s.).

19. On Germinal Selection. By August Weismann. 25c (1s. 6d.).

20. Lovers Three Thousand Years Ago. By T. A. Goodwin. (Out of print.)

21. Popular Scientific Lectures. By Ernst Mach. 50c (2s. 6d.).

22. Ancient India: Its Language and Religions. By H. Oldenberg. 25c
(1s. 6d.).

23. The Prophets of Israel. By Prof. C. H. Cornill. 25c (1s. 6d.).

24. Homilies of Science. By Paul Carus. 35c (2s.).

25. Thoughts on Religion. By G. J. Romanes. 50c (2s. 6d.).

26. The Philosophy of Ancient India. By Prof. Richard Garbe. 25c (1s. 6d.).

27. Martin Luther. By Gustav Freytag. 25c (1s. 6d.).

28. English Secularism. By George Jacob Holyoake. 25c (1s. 6d.).

29. On Orthogenesis. By Th. Eimer. 25c (1s. 6d.).

30. Chinese Philosophy. By Paul Carus. 25c (1s. 6d.).

31. The Lost Manuscript. By Gustav Freytag. 60c (3s.).

32. A Mechanico-Physiological Theory of Organic Evolution. By Carl von
Naegeli. 15c (9d.).

33. Chinese Fiction. By Dr. George T. Candlin. 15c (9d.).

34. Mathematical Essays and Recreations. By H. Schubert. 25c (1s. 6d.).

35. The Ethical Problem. By Paul Carus. 50c (2s. 6d.).

36. Buddhism and Its Christian Critics. By Paul Carus. 50c (2s. 6d.).

37. Psychology for Beginners. By Hiram M. Stanley. 20c (1s.).

38. Discourse on Method. By Descartes. 25c (1s. 6d.).

THE OPEN COURT PUBLISHING CO.,

CHICAGO: 324 Dearborn Street.
LONDON: Kegan Paul, Trench, Trübner & Company.
\fi

%%%%%%%%%%%%%%%%%%%%%%%%% GUTENBERG LICENSE %%%%%%%%%%%%%%%%%%%%%%%%%%
\PGLicense
\begin{PGtext}
End of the Project Gutenberg EBook of Elementary Illustrations of the
Differential and Integral Calculus, by Augustus De Morgan

*** END OF THIS PROJECT GUTENBERG EBOOK DIFFERENTIAL, INTEGRAL CALCULUS ***

***** This file should be named 39041-pdf.pdf or 39041-pdf.zip *****
This and all associated files of various formats will be found in:
        http://www.gutenberg.org/3/9/0/4/39041/

Produced by Andrew D. Hwang.

Updated editions will replace the previous one--the old editions
will be renamed.

Creating the works from public domain print editions means that no
one owns a United States copyright in these works, so the Foundation
(and you!) can copy and distribute it in the United States without
permission and without paying copyright royalties.  Special rules,
set forth in the General Terms of Use part of this license, apply to
copying and distributing Project Gutenberg-tm electronic works to
protect the PROJECT GUTENBERG-tm concept and trademark.  Project
Gutenberg is a registered trademark, and may not be used if you
charge for the eBooks, unless you receive specific permission.  If you
do not charge anything for copies of this eBook, complying with the
rules is very easy.  You may use this eBook for nearly any purpose
such as creation of derivative works, reports, performances and
research.  They may be modified and printed and given away--you may do
practically ANYTHING with public domain eBooks.  Redistribution is
subject to the trademark license, especially commercial
redistribution.



*** START: FULL LICENSE ***

THE FULL PROJECT GUTENBERG LICENSE
PLEASE READ THIS BEFORE YOU DISTRIBUTE OR USE THIS WORK

To protect the Project Gutenberg-tm mission of promoting the free
distribution of electronic works, by using or distributing this work
(or any other work associated in any way with the phrase "Project
Gutenberg"), you agree to comply with all the terms of the Full Project
Gutenberg-tm License (available with this file or online at
http://gutenberg.net/license).


Section 1.  General Terms of Use and Redistributing Project Gutenberg-tm
electronic works

1.A.  By reading or using any part of this Project Gutenberg-tm
electronic work, you indicate that you have read, understand, agree to
and accept all the terms of this license and intellectual property
(trademark/copyright) agreement.  If you do not agree to abide by all
the terms of this agreement, you must cease using and return or destroy
all copies of Project Gutenberg-tm electronic works in your possession.
If you paid a fee for obtaining a copy of or access to a Project
Gutenberg-tm electronic work and you do not agree to be bound by the
terms of this agreement, you may obtain a refund from the person or
entity to whom you paid the fee as set forth in paragraph 1.E.8.

1.B.  "Project Gutenberg" is a registered trademark.  It may only be
used on or associated in any way with an electronic work by people who
agree to be bound by the terms of this agreement.  There are a few
things that you can do with most Project Gutenberg-tm electronic works
even without complying with the full terms of this agreement.  See
paragraph 1.C below.  There are a lot of things you can do with Project
Gutenberg-tm electronic works if you follow the terms of this agreement
and help preserve free future access to Project Gutenberg-tm electronic
works.  See paragraph 1.E below.

1.C.  The Project Gutenberg Literary Archive Foundation ("the Foundation"
or PGLAF), owns a compilation copyright in the collection of Project
Gutenberg-tm electronic works.  Nearly all the individual works in the
collection are in the public domain in the United States.  If an
individual work is in the public domain in the United States and you are
located in the United States, we do not claim a right to prevent you from
copying, distributing, performing, displaying or creating derivative
works based on the work as long as all references to Project Gutenberg
are removed.  Of course, we hope that you will support the Project
Gutenberg-tm mission of promoting free access to electronic works by
freely sharing Project Gutenberg-tm works in compliance with the terms of
this agreement for keeping the Project Gutenberg-tm name associated with
the work.  You can easily comply with the terms of this agreement by
keeping this work in the same format with its attached full Project
Gutenberg-tm License when you share it without charge with others.

1.D.  The copyright laws of the place where you are located also govern
what you can do with this work.  Copyright laws in most countries are in
a constant state of change.  If you are outside the United States, check
the laws of your country in addition to the terms of this agreement
before downloading, copying, displaying, performing, distributing or
creating derivative works based on this work or any other Project
Gutenberg-tm work.  The Foundation makes no representations concerning
the copyright status of any work in any country outside the United
States.

1.E.  Unless you have removed all references to Project Gutenberg:

1.E.1.  The following sentence, with active links to, or other immediate
access to, the full Project Gutenberg-tm License must appear prominently
whenever any copy of a Project Gutenberg-tm work (any work on which the
phrase "Project Gutenberg" appears, or with which the phrase "Project
Gutenberg" is associated) is accessed, displayed, performed, viewed,
copied or distributed:

This eBook is for the use of anyone anywhere at no cost and with
almost no restrictions whatsoever.  You may copy it, give it away or
re-use it under the terms of the Project Gutenberg License included
with this eBook or online at www.gutenberg.net

1.E.2.  If an individual Project Gutenberg-tm electronic work is derived
from the public domain (does not contain a notice indicating that it is
posted with permission of the copyright holder), the work can be copied
and distributed to anyone in the United States without paying any fees
or charges.  If you are redistributing or providing access to a work
with the phrase "Project Gutenberg" associated with or appearing on the
work, you must comply either with the requirements of paragraphs 1.E.1
through 1.E.7 or obtain permission for the use of the work and the
Project Gutenberg-tm trademark as set forth in paragraphs 1.E.8 or
1.E.9.

1.E.3.  If an individual Project Gutenberg-tm electronic work is posted
with the permission of the copyright holder, your use and distribution
must comply with both paragraphs 1.E.1 through 1.E.7 and any additional
terms imposed by the copyright holder.  Additional terms will be linked
to the Project Gutenberg-tm License for all works posted with the
permission of the copyright holder found at the beginning of this work.

1.E.4.  Do not unlink or detach or remove the full Project Gutenberg-tm
License terms from this work, or any files containing a part of this
work or any other work associated with Project Gutenberg-tm.

1.E.5.  Do not copy, display, perform, distribute or redistribute this
electronic work, or any part of this electronic work, without
prominently displaying the sentence set forth in paragraph 1.E.1 with
active links or immediate access to the full terms of the Project
Gutenberg-tm License.

1.E.6.  You may convert to and distribute this work in any binary,
compressed, marked up, nonproprietary or proprietary form, including any
word processing or hypertext form.  However, if you provide access to or
distribute copies of a Project Gutenberg-tm work in a format other than
"Plain Vanilla ASCII" or other format used in the official version
posted on the official Project Gutenberg-tm web site (www.gutenberg.net),
you must, at no additional cost, fee or expense to the user, provide a
copy, a means of exporting a copy, or a means of obtaining a copy upon
request, of the work in its original "Plain Vanilla ASCII" or other
form.  Any alternate format must include the full Project Gutenberg-tm
License as specified in paragraph 1.E.1.

1.E.7.  Do not charge a fee for access to, viewing, displaying,
performing, copying or distributing any Project Gutenberg-tm works
unless you comply with paragraph 1.E.8 or 1.E.9.

1.E.8.  You may charge a reasonable fee for copies of or providing
access to or distributing Project Gutenberg-tm electronic works provided
that

- You pay a royalty fee of 20% of the gross profits you derive from
     the use of Project Gutenberg-tm works calculated using the method
     you already use to calculate your applicable taxes.  The fee is
     owed to the owner of the Project Gutenberg-tm trademark, but he
     has agreed to donate royalties under this paragraph to the
     Project Gutenberg Literary Archive Foundation.  Royalty payments
     must be paid within 60 days following each date on which you
     prepare (or are legally required to prepare) your periodic tax
     returns.  Royalty payments should be clearly marked as such and
     sent to the Project Gutenberg Literary Archive Foundation at the
     address specified in Section 4, "Information about donations to
     the Project Gutenberg Literary Archive Foundation."

- You provide a full refund of any money paid by a user who notifies
     you in writing (or by e-mail) within 30 days of receipt that s/he
     does not agree to the terms of the full Project Gutenberg-tm
     License.  You must require such a user to return or
     destroy all copies of the works possessed in a physical medium
     and discontinue all use of and all access to other copies of
     Project Gutenberg-tm works.

- You provide, in accordance with paragraph 1.F.3, a full refund of any
     money paid for a work or a replacement copy, if a defect in the
     electronic work is discovered and reported to you within 90 days
     of receipt of the work.

- You comply with all other terms of this agreement for free
     distribution of Project Gutenberg-tm works.

1.E.9.  If you wish to charge a fee or distribute a Project Gutenberg-tm
electronic work or group of works on different terms than are set
forth in this agreement, you must obtain permission in writing from
both the Project Gutenberg Literary Archive Foundation and Michael
Hart, the owner of the Project Gutenberg-tm trademark.  Contact the
Foundation as set forth in Section 3 below.

1.F.

1.F.1.  Project Gutenberg volunteers and employees expend considerable
effort to identify, do copyright research on, transcribe and proofread
public domain works in creating the Project Gutenberg-tm
collection.  Despite these efforts, Project Gutenberg-tm electronic
works, and the medium on which they may be stored, may contain
"Defects," such as, but not limited to, incomplete, inaccurate or
corrupt data, transcription errors, a copyright or other intellectual
property infringement, a defective or damaged disk or other medium, a
computer virus, or computer codes that damage or cannot be read by
your equipment.

1.F.2.  LIMITED WARRANTY, DISCLAIMER OF DAMAGES - Except for the "Right
of Replacement or Refund" described in paragraph 1.F.3, the Project
Gutenberg Literary Archive Foundation, the owner of the Project
Gutenberg-tm trademark, and any other party distributing a Project
Gutenberg-tm electronic work under this agreement, disclaim all
liability to you for damages, costs and expenses, including legal
fees.  YOU AGREE THAT YOU HAVE NO REMEDIES FOR NEGLIGENCE, STRICT
LIABILITY, BREACH OF WARRANTY OR BREACH OF CONTRACT EXCEPT THOSE
PROVIDED IN PARAGRAPH 1.F.3.  YOU AGREE THAT THE FOUNDATION, THE
TRADEMARK OWNER, AND ANY DISTRIBUTOR UNDER THIS AGREEMENT WILL NOT BE
LIABLE TO YOU FOR ACTUAL, DIRECT, INDIRECT, CONSEQUENTIAL, PUNITIVE OR
INCIDENTAL DAMAGES EVEN IF YOU GIVE NOTICE OF THE POSSIBILITY OF SUCH
DAMAGE.

1.F.3.  LIMITED RIGHT OF REPLACEMENT OR REFUND - If you discover a
defect in this electronic work within 90 days of receiving it, you can
receive a refund of the money (if any) you paid for it by sending a
written explanation to the person you received the work from.  If you
received the work on a physical medium, you must return the medium with
your written explanation.  The person or entity that provided you with
the defective work may elect to provide a replacement copy in lieu of a
refund.  If you received the work electronically, the person or entity
providing it to you may choose to give you a second opportunity to
receive the work electronically in lieu of a refund.  If the second copy
is also defective, you may demand a refund in writing without further
opportunities to fix the problem.

1.F.4.  Except for the limited right of replacement or refund set forth
in paragraph 1.F.3, this work is provided to you 'AS-IS' WITH NO OTHER
WARRANTIES OF ANY KIND, EXPRESS OR IMPLIED, INCLUDING BUT NOT LIMITED TO
WARRANTIES OF MERCHANTIBILITY OR FITNESS FOR ANY PURPOSE.

1.F.5.  Some states do not allow disclaimers of certain implied
warranties or the exclusion or limitation of certain types of damages.
If any disclaimer or limitation set forth in this agreement violates the
law of the state applicable to this agreement, the agreement shall be
interpreted to make the maximum disclaimer or limitation permitted by
the applicable state law.  The invalidity or unenforceability of any
provision of this agreement shall not void the remaining provisions.

1.F.6.  INDEMNITY - You agree to indemnify and hold the Foundation, the
trademark owner, any agent or employee of the Foundation, anyone
providing copies of Project Gutenberg-tm electronic works in accordance
with this agreement, and any volunteers associated with the production,
promotion and distribution of Project Gutenberg-tm electronic works,
harmless from all liability, costs and expenses, including legal fees,
that arise directly or indirectly from any of the following which you do
or cause to occur: (a) distribution of this or any Project Gutenberg-tm
work, (b) alteration, modification, or additions or deletions to any
Project Gutenberg-tm work, and (c) any Defect you cause.


Section  2.  Information about the Mission of Project Gutenberg-tm

Project Gutenberg-tm is synonymous with the free distribution of
electronic works in formats readable by the widest variety of computers
including obsolete, old, middle-aged and new computers.  It exists
because of the efforts of hundreds of volunteers and donations from
people in all walks of life.

Volunteers and financial support to provide volunteers with the
assistance they need are critical to reaching Project Gutenberg-tm's
goals and ensuring that the Project Gutenberg-tm collection will
remain freely available for generations to come.  In 2001, the Project
Gutenberg Literary Archive Foundation was created to provide a secure
and permanent future for Project Gutenberg-tm and future generations.
To learn more about the Project Gutenberg Literary Archive Foundation
and how your efforts and donations can help, see Sections 3 and 4
and the Foundation web page at http://www.pglaf.org.


Section 3.  Information about the Project Gutenberg Literary Archive
Foundation

The Project Gutenberg Literary Archive Foundation is a non profit
501(c)(3) educational corporation organized under the laws of the
state of Mississippi and granted tax exempt status by the Internal
Revenue Service.  The Foundation's EIN or federal tax identification
number is 64-6221541.  Its 501(c)(3) letter is posted at
http://pglaf.org/fundraising.  Contributions to the Project Gutenberg
Literary Archive Foundation are tax deductible to the full extent
permitted by U.S. federal laws and your state's laws.

The Foundation's principal office is located at 4557 Melan Dr. S.
Fairbanks, AK, 99712., but its volunteers and employees are scattered
throughout numerous locations.  Its business office is located at
809 North 1500 West, Salt Lake City, UT 84116, (801) 596-1887, email
business@pglaf.org.  Email contact links and up to date contact
information can be found at the Foundation's web site and official
page at http://pglaf.org

For additional contact information:
     Dr. Gregory B. Newby
     Chief Executive and Director
     gbnewby@pglaf.org


Section 4.  Information about Donations to the Project Gutenberg
Literary Archive Foundation

Project Gutenberg-tm depends upon and cannot survive without wide
spread public support and donations to carry out its mission of
increasing the number of public domain and licensed works that can be
freely distributed in machine readable form accessible by the widest
array of equipment including outdated equipment.  Many small donations
($1 to $5,000) are particularly important to maintaining tax exempt
status with the IRS.

The Foundation is committed to complying with the laws regulating
charities and charitable donations in all 50 states of the United
States.  Compliance requirements are not uniform and it takes a
considerable effort, much paperwork and many fees to meet and keep up
with these requirements.  We do not solicit donations in locations
where we have not received written confirmation of compliance.  To
SEND DONATIONS or determine the status of compliance for any
particular state visit http://pglaf.org

While we cannot and do not solicit contributions from states where we
have not met the solicitation requirements, we know of no prohibition
against accepting unsolicited donations from donors in such states who
approach us with offers to donate.

International donations are gratefully accepted, but we cannot make
any statements concerning tax treatment of donations received from
outside the United States.  U.S. laws alone swamp our small staff.

Please check the Project Gutenberg Web pages for current donation
methods and addresses.  Donations are accepted in a number of other
ways including including checks, online payments and credit card
donations.  To donate, please visit: http://pglaf.org/donate


Section 5.  General Information About Project Gutenberg-tm electronic
works.

Professor Michael S. Hart is the originator of the Project Gutenberg-tm
concept of a library of electronic works that could be freely shared
with anyone.  For thirty years, he produced and distributed Project
Gutenberg-tm eBooks with only a loose network of volunteer support.


Project Gutenberg-tm eBooks are often created from several printed
editions, all of which are confirmed as Public Domain in the U.S.
unless a copyright notice is included.  Thus, we do not necessarily
keep eBooks in compliance with any particular paper edition.


Most people start at our Web site which has the main PG search facility:

     http://www.gutenberg.net

This Web site includes information about Project Gutenberg-tm,
including how to make donations to the Project Gutenberg Literary
Archive Foundation, how to help produce our new eBooks, and how to
subscribe to our email newsletter to hear about new eBooks.
\end{PGtext}

% %%%%%%%%%%%%%%%%%%%%%%%%%%%%%%%%%%%%%%%%%%%%%%%%%%%%%%%%%%%%%%%%%%%%%%% %
%                                                                         %
% End of the Project Gutenberg EBook of Elementary Illustrations of the   %
% Differential and Integral Calculus, by Augustus De Morgan               %
%                                                                         %
% *** END OF THIS PROJECT GUTENBERG EBOOK DIFFERENTIAL, INTEGRAL CALCULUS ***
%                                                                         %
% ***** This file should be named 39041-t.tex or 39041-t.zip *****        %
% This and all associated files of various formats will be found in:      %
%         http://www.gutenberg.org/3/9/0/4/39041/                         %
%                                                                         %
% %%%%%%%%%%%%%%%%%%%%%%%%%%%%%%%%%%%%%%%%%%%%%%%%%%%%%%%%%%%%%%%%%%%%%%% %

\end{document}
###
@ControlwordReplace = (
  ['\\th', 'th'],
  ['\\ie', 'i.e.'],
  ['\\eg', 'e.g.']
  );

@ControlwordArguments = (
  ['\\Signature', 1, 1, '', ' ', 1, 1, '', ' '],
  ['\\Figure', 0, 0, '', '', 1, 0, '<GRAPHIC>', ''],
  ['\\BookMark', 1, 0, '', '', 1, 0, '', ''],
  ['\\First', 1, 1, '', ''],
  ['\\Section', 0, 0, '', '', 1, 1, '', ''],
  ['\\SubSectHead', 1, 1, '', ''],
  ['\\Subsection', 0, 0, '', '', 1, 1, '', ''],
  ['\\Fig', 1, 1, 'Fig. ', ''],
  ['\\Pagelabel', 1, 0, '', ''],
  ['\\Pageref', 1, 1, 'p. ', ''],
  ['\\Pagerefs', 1, 1, 'pp. ', ', ', 1, 1, '', ''],
  ['\\Eq', 1, 1, '', ''],
  ['\\Typo', 1, 0, '', '', 1, 1, '', ''],
  ['\\Add', 1, 1, '', ''],
  ['\\Chg', 1, 0, '', '', 1, 1, '', '']
  );

$PageSeparator = qr/^\\PageSep/;
$CustomClean = 'print "\\nCustom cleaning in progress...";
my $cline = 0;
 while ($cline <= $#file) {
   $file[$cline] =~ s/--------[^\n]*\n//; # strip page separators
   $cline++
 }
 print "done\\n";';
###
This is pdfTeX, Version 3.1415926-1.40.10 (TeX Live 2009/Debian) (format=pdflatex 2011.9.6)  3 MAR 2012 15:22
entering extended mode
 %&-line parsing enabled.
**39041-t.tex
(./39041-t.tex
LaTeX2e <2009/09/24>
Babel <v3.8l> and hyphenation patterns for english, usenglishmax, dumylang, noh
yphenation, farsi, arabic, croatian, bulgarian, ukrainian, russian, czech, slov
ak, danish, dutch, finnish, french, basque, ngerman, german, german-x-2009-06-1
9, ngerman-x-2009-06-19, ibycus, monogreek, greek, ancientgreek, hungarian, san
skrit, italian, latin, latvian, lithuanian, mongolian2a, mongolian, bokmal, nyn
orsk, romanian, irish, coptic, serbian, turkish, welsh, esperanto, uppersorbian
, estonian, indonesian, interlingua, icelandic, kurmanji, slovenian, polish, po
rtuguese, spanish, galician, catalan, swedish, ukenglish, pinyin, loaded.
(/usr/share/texmf-texlive/tex/latex/base/book.cls
Document Class: book 2007/10/19 v1.4h Standard LaTeX document class
(/usr/share/texmf-texlive/tex/latex/base/bk12.clo
File: bk12.clo 2007/10/19 v1.4h Standard LaTeX file (size option)
)
\c@part=\count79
\c@chapter=\count80
\c@section=\count81
\c@subsection=\count82
\c@subsubsection=\count83
\c@paragraph=\count84
\c@subparagraph=\count85
\c@figure=\count86
\c@table=\count87
\abovecaptionskip=\skip41
\belowcaptionskip=\skip42
\bibindent=\dimen102
) (/usr/share/texmf-texlive/tex/latex/base/inputenc.sty
Package: inputenc 2008/03/30 v1.1d Input encoding file
\inpenc@prehook=\toks14
\inpenc@posthook=\toks15
(/usr/share/texmf-texlive/tex/latex/base/latin1.def
File: latin1.def 2008/03/30 v1.1d Input encoding file
)) (/usr/share/texmf-texlive/tex/latex/base/ifthen.sty
Package: ifthen 2001/05/26 v1.1c Standard LaTeX ifthen package (DPC)
) (/usr/share/texmf-texlive/tex/latex/amsmath/amsmath.sty
Package: amsmath 2000/07/18 v2.13 AMS math features
\@mathmargin=\skip43
For additional information on amsmath, use the `?' option.
(/usr/share/texmf-texlive/tex/latex/amsmath/amstext.sty
Package: amstext 2000/06/29 v2.01
(/usr/share/texmf-texlive/tex/latex/amsmath/amsgen.sty
File: amsgen.sty 1999/11/30 v2.0
\@emptytoks=\toks16
\ex@=\dimen103
)) (/usr/share/texmf-texlive/tex/latex/amsmath/amsbsy.sty
Package: amsbsy 1999/11/29 v1.2d
\pmbraise@=\dimen104
) (/usr/share/texmf-texlive/tex/latex/amsmath/amsopn.sty
Package: amsopn 1999/12/14 v2.01 operator names
)
\inf@bad=\count88
LaTeX Info: Redefining \frac on input line 211.
\uproot@=\count89
\leftroot@=\count90
LaTeX Info: Redefining \overline on input line 307.
\classnum@=\count91
\DOTSCASE@=\count92
LaTeX Info: Redefining \ldots on input line 379.
LaTeX Info: Redefining \dots on input line 382.
LaTeX Info: Redefining \cdots on input line 467.
\Mathstrutbox@=\box26
\strutbox@=\box27
\big@size=\dimen105
LaTeX Font Info:    Redeclaring font encoding OML on input line 567.
LaTeX Font Info:    Redeclaring font encoding OMS on input line 568.
\macc@depth=\count93
\c@MaxMatrixCols=\count94
\dotsspace@=\muskip10
\c@parentequation=\count95
\dspbrk@lvl=\count96
\tag@help=\toks17
\row@=\count97
\column@=\count98
\maxfields@=\count99
\andhelp@=\toks18
\eqnshift@=\dimen106
\alignsep@=\dimen107
\tagshift@=\dimen108
\tagwidth@=\dimen109
\totwidth@=\dimen110
\lineht@=\dimen111
\@envbody=\toks19
\multlinegap=\skip44
\multlinetaggap=\skip45
\mathdisplay@stack=\toks20
LaTeX Info: Redefining \[ on input line 2666.
LaTeX Info: Redefining \] on input line 2667.
) (/usr/share/texmf-texlive/tex/latex/amsfonts/amssymb.sty
Package: amssymb 2009/06/22 v3.00
(/usr/share/texmf-texlive/tex/latex/amsfonts/amsfonts.sty
Package: amsfonts 2009/06/22 v3.00 Basic AMSFonts support
\symAMSa=\mathgroup4
\symAMSb=\mathgroup5
LaTeX Font Info:    Overwriting math alphabet `\mathfrak' in version `bold'
(Font)                  U/euf/m/n --> U/euf/b/n on input line 96.
)) (/usr/share/texmf-texlive/tex/latex/base/alltt.sty
Package: alltt 1997/06/16 v2.0g defines alltt environment
) (/usr/share/texmf-texlive/tex/latex/tools/indentfirst.sty
Package: indentfirst 1995/11/23 v1.03 Indent first paragraph (DPC)
) (/usr/share/texmf-texlive/tex/latex/yfonts/yfonts.sty
Package: yfonts 2003/01/08 v1.3 (WaS)
) (/usr/share/texmf-texlive/tex/latex/footmisc/footmisc.sty
Package: footmisc 2009/09/15 v5.5a a miscellany of footnote facilities
\FN@temptoken=\toks21
\footnotemargin=\dimen112
\c@pp@next@reset=\count100
\c@@fnserial=\count101
Package footmisc Info: Declaring symbol style bringhurst on input line 855.
Package footmisc Info: Declaring symbol style chicago on input line 863.
Package footmisc Info: Declaring symbol style wiley on input line 872.
Package footmisc Info: Declaring symbol style lamport-robust on input line 883.

Package footmisc Info: Declaring symbol style lamport* on input line 903.
Package footmisc Info: Declaring symbol style lamport*-robust on input line 924
.
) (/usr/share/texmf-texlive/tex/latex/was/icomma.sty
Package: icomma 2002/03/10 v2.0 (WaS)
) (/usr/share/texmf-texlive/tex/latex/tools/multicol.sty
Package: multicol 2008/12/05 v1.6h multicolumn formatting (FMi)
\c@tracingmulticols=\count102
\mult@box=\box28
\multicol@leftmargin=\dimen113
\c@unbalance=\count103
\c@collectmore=\count104
\doublecol@number=\count105
\multicoltolerance=\count106
\multicolpretolerance=\count107
\full@width=\dimen114
\page@free=\dimen115
\premulticols=\dimen116
\postmulticols=\dimen117
\multicolsep=\skip46
\multicolbaselineskip=\skip47
\partial@page=\box29
\last@line=\box30
\mult@rightbox=\box31
\mult@grightbox=\box32
\mult@gfirstbox=\box33
\mult@firstbox=\box34
\@tempa=\box35
\@tempa=\box36
\@tempa=\box37
\@tempa=\box38
\@tempa=\box39
\@tempa=\box40
\@tempa=\box41
\@tempa=\box42
\@tempa=\box43
\@tempa=\box44
\@tempa=\box45
\@tempa=\box46
\@tempa=\box47
\@tempa=\box48
\@tempa=\box49
\@tempa=\box50
\@tempa=\box51
\c@columnbadness=\count108
\c@finalcolumnbadness=\count109
\last@try=\dimen118
\multicolovershoot=\dimen119
\multicolundershoot=\dimen120
\mult@nat@firstbox=\box52
\colbreak@box=\box53
) (/usr/share/texmf-texlive/tex/latex/base/makeidx.sty
Package: makeidx 2000/03/29 v1.0m Standard LaTeX package
) (/usr/share/texmf-texlive/tex/latex/tools/array.sty
Package: array 2008/09/09 v2.4c Tabular extension package (FMi)
\col@sep=\dimen121
\extrarowheight=\dimen122
\NC@list=\toks22
\extratabsurround=\skip48
\backup@length=\skip49
) (/usr/share/texmf-texlive/tex/latex/graphics/graphicx.sty
Package: graphicx 1999/02/16 v1.0f Enhanced LaTeX Graphics (DPC,SPQR)
(/usr/share/texmf-texlive/tex/latex/graphics/keyval.sty
Package: keyval 1999/03/16 v1.13 key=value parser (DPC)
\KV@toks@=\toks23
) (/usr/share/texmf-texlive/tex/latex/graphics/graphics.sty
Package: graphics 2009/02/05 v1.0o Standard LaTeX Graphics (DPC,SPQR)
(/usr/share/texmf-texlive/tex/latex/graphics/trig.sty
Package: trig 1999/03/16 v1.09 sin cos tan (DPC)
) (/etc/texmf/tex/latex/config/graphics.cfg
File: graphics.cfg 2009/08/28 v1.8 graphics configuration of TeX Live
)
Package graphics Info: Driver file: pdftex.def on input line 91.
(/usr/share/texmf-texlive/tex/latex/pdftex-def/pdftex.def
File: pdftex.def 2009/08/25 v0.04m Graphics/color for pdfTeX
\Gread@gobject=\count110
))
\Gin@req@height=\dimen123
\Gin@req@width=\dimen124
) (/usr/share/texmf-texlive/tex/latex/tools/calc.sty
Package: calc 2007/08/22 v4.3 Infix arithmetic (KKT,FJ)
\calc@Acount=\count111
\calc@Bcount=\count112
\calc@Adimen=\dimen125
\calc@Bdimen=\dimen126
\calc@Askip=\skip50
\calc@Bskip=\skip51
LaTeX Info: Redefining \setlength on input line 76.
LaTeX Info: Redefining \addtolength on input line 77.
\calc@Ccount=\count113
\calc@Cskip=\skip52
) (/usr/share/texmf-texlive/tex/latex/fancyhdr/fancyhdr.sty
\fancy@headwidth=\skip53
\f@ncyO@elh=\skip54
\f@ncyO@erh=\skip55
\f@ncyO@olh=\skip56
\f@ncyO@orh=\skip57
\f@ncyO@elf=\skip58
\f@ncyO@erf=\skip59
\f@ncyO@olf=\skip60
\f@ncyO@orf=\skip61
) (/usr/share/texmf-texlive/tex/latex/geometry/geometry.sty
Package: geometry 2008/12/21 v4.2 Page Geometry
(/usr/share/texmf-texlive/tex/generic/oberdiek/ifpdf.sty
Package: ifpdf 2009/04/10 v2.0 Provides the ifpdf switch (HO)
Package ifpdf Info: pdfTeX in pdf mode detected.
) (/usr/share/texmf-texlive/tex/generic/oberdiek/ifvtex.sty
Package: ifvtex 2008/11/04 v1.4 Switches for detecting VTeX and its modes (HO)
Package ifvtex Info: VTeX not detected.
)
\Gm@cnth=\count114
\Gm@cntv=\count115
\c@Gm@tempcnt=\count116
\Gm@bindingoffset=\dimen127
\Gm@wd@mp=\dimen128
\Gm@odd@mp=\dimen129
\Gm@even@mp=\dimen130
\Gm@dimlist=\toks24
(/usr/share/texmf-texlive/tex/xelatex/xetexconfig/geometry.cfg)) (/usr/share/te
xmf-texlive/tex/latex/hyperref/hyperref.sty
Package: hyperref 2009/10/09 v6.79a Hypertext links for LaTeX
(/usr/share/texmf-texlive/tex/generic/ifxetex/ifxetex.sty
Package: ifxetex 2009/01/23 v0.5 Provides ifxetex conditional
) (/usr/share/texmf-texlive/tex/latex/oberdiek/hycolor.sty
Package: hycolor 2009/10/02 v1.5 Code for color options of hyperref/bookmark (H
O)
(/usr/share/texmf-texlive/tex/latex/oberdiek/xcolor-patch.sty
Package: xcolor-patch 2009/10/02 xcolor patch
))
\@linkdim=\dimen131
\Hy@linkcounter=\count117
\Hy@pagecounter=\count118
(/usr/share/texmf-texlive/tex/latex/hyperref/pd1enc.def
File: pd1enc.def 2009/10/09 v6.79a Hyperref: PDFDocEncoding definition (HO)
) (/usr/share/texmf-texlive/tex/generic/oberdiek/etexcmds.sty
Package: etexcmds 2007/12/12 v1.2 Prefix for e-TeX command names (HO)
(/usr/share/texmf-texlive/tex/generic/oberdiek/infwarerr.sty
Package: infwarerr 2007/09/09 v1.2 Providing info/warning/message (HO)
)
Package etexcmds Info: Could not find \expanded.
(etexcmds)             That can mean that you are not using pdfTeX 1.50 or
(etexcmds)             that some package has redefined \expanded.
(etexcmds)             In the latter case, load this package earlier.
) (/etc/texmf/tex/latex/config/hyperref.cfg
File: hyperref.cfg 2002/06/06 v1.2 hyperref configuration of TeXLive
) (/usr/share/texmf-texlive/tex/latex/oberdiek/kvoptions.sty
Package: kvoptions 2009/08/13 v3.4 Keyval support for LaTeX options (HO)
(/usr/share/texmf-texlive/tex/generic/oberdiek/kvsetkeys.sty
Package: kvsetkeys 2009/07/30 v1.5 Key value parser with default handler suppor
t (HO)
))
Package hyperref Info: Option `hyperfootnotes' set `false' on input line 2864.
Package hyperref Info: Option `bookmarks' set `true' on input line 2864.
Package hyperref Info: Option `linktocpage' set `false' on input line 2864.
Package hyperref Info: Option `pdfdisplaydoctitle' set `true' on input line 286
4.
Package hyperref Info: Option `pdfpagelabels' set `true' on input line 2864.
Package hyperref Info: Option `bookmarksopen' set `true' on input line 2864.
Package hyperref Info: Option `colorlinks' set `true' on input line 2864.
Package hyperref Info: Hyper figures OFF on input line 2975.
Package hyperref Info: Link nesting OFF on input line 2980.
Package hyperref Info: Hyper index ON on input line 2983.
Package hyperref Info: Plain pages OFF on input line 2990.
Package hyperref Info: Backreferencing OFF on input line 2995.
Implicit mode ON; LaTeX internals redefined
Package hyperref Info: Bookmarks ON on input line 3191.
(/usr/share/texmf-texlive/tex/latex/ltxmisc/url.sty
\Urlmuskip=\muskip11
Package: url 2006/04/12  ver 3.3  Verb mode for urls, etc.
)
LaTeX Info: Redefining \url on input line 3428.
(/usr/share/texmf-texlive/tex/generic/oberdiek/bitset.sty
Package: bitset 2007/09/28 v1.0 Data type bit set (HO)
(/usr/share/texmf-texlive/tex/generic/oberdiek/intcalc.sty
Package: intcalc 2007/09/27 v1.1 Expandable integer calculations (HO)
) (/usr/share/texmf-texlive/tex/generic/oberdiek/bigintcalc.sty
Package: bigintcalc 2007/11/11 v1.1 Expandable big integer calculations (HO)
(/usr/share/texmf-texlive/tex/generic/oberdiek/pdftexcmds.sty
Package: pdftexcmds 2009/09/23 v0.6 LuaTeX support for pdfTeX utility functions
 (HO)
(/usr/share/texmf-texlive/tex/generic/oberdiek/ifluatex.sty
Package: ifluatex 2009/04/17 v1.2 Provides the ifluatex switch (HO)
Package ifluatex Info: LuaTeX not detected.
) (/usr/share/texmf-texlive/tex/generic/oberdiek/ltxcmds.sty
Package: ltxcmds 2009/08/05 v1.0 Some LaTeX kernel commands for general use (HO
)
)
Package pdftexcmds Info: LuaTeX not detected.
Package pdftexcmds Info: \pdf@primitive is available.
Package pdftexcmds Info: \pdf@ifprimitive is available.
)))
\Fld@menulength=\count119
\Field@Width=\dimen132
\Fld@charsize=\dimen133
\Field@toks=\toks25
Package hyperref Info: Hyper figures OFF on input line 4377.
Package hyperref Info: Link nesting OFF on input line 4382.
Package hyperref Info: Hyper index ON on input line 4385.
Package hyperref Info: backreferencing OFF on input line 4392.
Package hyperref Info: Link coloring ON on input line 4395.
Package hyperref Info: Link coloring with OCG OFF on input line 4402.
Package hyperref Info: PDF/A mode OFF on input line 4407.
(/usr/share/texmf-texlive/tex/generic/oberdiek/atbegshi.sty
Package: atbegshi 2008/07/31 v1.9 At begin shipout hook (HO)
)
\Hy@abspage=\count120
\c@Item=\count121
)
*hyperref using driver hpdftex*
(/usr/share/texmf-texlive/tex/latex/hyperref/hpdftex.def
File: hpdftex.def 2009/10/09 v6.79a Hyperref driver for pdfTeX
\Fld@listcount=\count122
)
\TmpLen=\skip62
\c@SecNo=\count123
\@indexfile=\write3
\openout3 = `39041-t.idx'.

Writing index file 39041-t.idx
\ParIndent=\skip63
(./39041-t.aux)
\openout1 = `39041-t.aux'.

LaTeX Font Info:    Checking defaults for OML/cmm/m/it on input line 605.
LaTeX Font Info:    ... okay on input line 605.
LaTeX Font Info:    Checking defaults for T1/cmr/m/n on input line 605.
LaTeX Font Info:    ... okay on input line 605.
LaTeX Font Info:    Checking defaults for OT1/cmr/m/n on input line 605.
LaTeX Font Info:    ... okay on input line 605.
LaTeX Font Info:    Checking defaults for OMS/cmsy/m/n on input line 605.
LaTeX Font Info:    ... okay on input line 605.
LaTeX Font Info:    Checking defaults for OMX/cmex/m/n on input line 605.
LaTeX Font Info:    ... okay on input line 605.
LaTeX Font Info:    Checking defaults for U/cmr/m/n on input line 605.
LaTeX Font Info:    ... okay on input line 605.
LaTeX Font Info:    Checking defaults for LY/yfrak/m/n on input line 605.
LaTeX Font Info:    ... okay on input line 605.
LaTeX Font Info:    Checking defaults for LYG/ygoth/m/n on input line 605.
LaTeX Font Info:    ... okay on input line 605.
LaTeX Font Info:    Checking defaults for PD1/pdf/m/n on input line 605.
LaTeX Font Info:    ... okay on input line 605.
(/usr/share/texmf/tex/context/base/supp-pdf.mkii
[Loading MPS to PDF converter (version 2006.09.02).]
\scratchcounter=\count124
\scratchdimen=\dimen134
\scratchbox=\box54
\nofMPsegments=\count125
\nofMParguments=\count126
\everyMPshowfont=\toks26
\MPscratchCnt=\count127
\MPscratchDim=\dimen135
\MPnumerator=\count128
\everyMPtoPDFconversion=\toks27
)
*geometry auto-detecting driver*
*geometry detected driver: pdftex*
-------------------- Geometry parameters
paper: class default
landscape: --
twocolumn: --
twoside: true
asymmetric: --
h-parts: 9.03375pt, 307.14749pt, 9.03375pt
v-parts: 1.26749pt, 466.58623pt, 1.90128pt
hmarginratio: 1:1
vmarginratio: 2:3
lines: --
heightrounded: --
bindingoffset: 0.0pt
truedimen: --
includehead: true
includefoot: true
includemp: --
driver: pdftex
-------------------- Page layout dimensions and switches
\paperwidth  325.215pt
\paperheight 469.75499pt
\textwidth  307.14749pt
\textheight 404.71243pt
\oddsidemargin  -63.23624pt
\evensidemargin -63.23624pt
\topmargin  -71.0025pt
\headheight 12.0pt
\headsep    19.8738pt
\footskip   30.0pt
\marginparwidth 98.0pt
\marginparsep   7.0pt
\columnsep  10.0pt
\skip\footins  10.8pt plus 4.0pt minus 2.0pt
\hoffset 0.0pt
\voffset 0.0pt
\mag 1000
\@twosidetrue \@mparswitchtrue 
(1in=72.27pt, 1cm=28.45pt)
-----------------------
(/usr/share/texmf-texlive/tex/latex/graphics/color.sty
Package: color 2005/11/14 v1.0j Standard LaTeX Color (DPC)
(/etc/texmf/tex/latex/config/color.cfg
File: color.cfg 2007/01/18 v1.5 color configuration of teTeX/TeXLive
)
Package color Info: Driver file: pdftex.def on input line 130.
)
Package hyperref Info: Link coloring ON on input line 605.
(/usr/share/texmf-texlive/tex/latex/hyperref/nameref.sty
Package: nameref 2007/05/29 v2.31 Cross-referencing by name of section
(/usr/share/texmf-texlive/tex/latex/oberdiek/refcount.sty
Package: refcount 2008/08/11 v3.1 Data extraction from references (HO)
)
\c@section@level=\count129
)
LaTeX Info: Redefining \ref on input line 605.
LaTeX Info: Redefining \pageref on input line 605.
(./39041-t.out) (./39041-t.out)
\@outlinefile=\write4
\openout4 = `39041-t.out'.

\AtBeginShipoutBox=\box55

Overfull \hbox (3.10696pt too wide) in paragraph at lines 621--621
[]\OT1/cmtt/m/n/8 Title: Elementary Illustrations of the Differential and Integ
ral Calculus[] 
 []


Overfull \hbox (20.10721pt too wide) in paragraph at lines 631--631
[]\OT1/cmtt/m/n/8 *** START OF THIS PROJECT GUTENBERG EBOOK DIFFERENTIAL, INTEG
RAL CALCULUS ***[] 
 []

LaTeX Font Info:    Try loading font information for U+msa on input line 633.
(/usr/share/texmf-texlive/tex/latex/amsfonts/umsa.fd
File: umsa.fd 2009/06/22 v3.00 AMS symbols A
)
LaTeX Font Info:    Try loading font information for U+msb on input line 633.
(/usr/share/texmf-texlive/tex/latex/amsfonts/umsb.fd
File: umsb.fd 2009/06/22 v3.00 AMS symbols B
) [1

{/var/lib/texmf/fonts/map/pdftex/updmap/pdftex.map}] [2]
Underfull \hbox (badness 1127) in paragraph at lines 669--675
\OT1/cmr/m/n/12 LECTURES ON EL-E-MEN-TARY MATH-E-MAT-ICS. By
 []

[1


] [2] [3

] [4] [5] [6] (./39041-t.toc)
\tf@toc=\write5
\openout5 = `39041-t.toc'.

[7


] [8] [9] [1


] [2] [3]
LaTeX Font Info:    Try loading font information for OMS+cmr on input line 999.

(/usr/share/texmf-texlive/tex/latex/base/omscmr.fd
File: omscmr.fd 1999/05/25 v2.5h Standard LaTeX font definitions
)
LaTeX Font Info:    Font shape `OMS/cmr/m/n' in size <8> not available
(Font)              Font shape `OMS/cmsy/m/n' tried instead on input line 999.
LaTeX Font Info:    Font shape `OMS/cmr/m/n' in size <7> not available
(Font)              Font shape `OMS/cmsy/m/n' tried instead on input line 1001.

[4] [5] [6] [7] <./images/fig1.png, id=244, 239.2137pt x 128.6406pt>
File: ./images/fig1.png Graphic file (type png)
<use ./images/fig1.png> [8] [9 <./images/fig1.png (PNG copy)>] [10]
File: ./images/fig1.png Graphic file (type png)
<use ./images/fig1.png> [11] [12] [13] [14] [15] [16] [17] [18] [19] [20] [21] 
[22] [23] [24] [25] [26] [27] [28] [29] [30] [31] [32] <./images/fig2.png, id=3
95, 234.8775pt x 121.4136pt>
File: ./images/fig2.png Graphic file (type png)
<use ./images/fig2.png> [33] [34] [35 <./images/fig2.png>] <./images/fig3.png, 
id=417, 235.6002pt x 199.2243pt>
File: ./images/fig3.png Graphic file (type png)
<use ./images/fig3.png> [36] [37 <./images/fig3.png>] [38] [39] [40]
File: ./images/fig2.png Graphic file (type png)
<use ./images/fig2.png> [41] [42] [43] [44] <./images/fig4.png, id=471, 235.600
2pt x 94.4328pt>
File: ./images/fig4.png Graphic file (type png)
<use ./images/fig4.png> [45] [46 <./images/fig4.png>] [47] [48] [49] [50] [51] 
<./images/fig5.png, id=545, 241.3818pt x 97.5645pt>
File: ./images/fig5.png Graphic file (type png)
<use ./images/fig5.png>
File: ./images/fig5.png Graphic file (type png)
<use ./images/fig5.png> [52 <./images/fig5.png (PNG copy)>] [53] [54] <./images
/fig6.png, id=582, 237.7683pt x 118.5228pt>
File: ./images/fig6.png Graphic file (type png)
<use ./images/fig6.png> <./images/fig7.png, id=583, 239.9364pt x 138.7584pt>
File: ./images/fig7.png Graphic file (type png)
<use ./images/fig7.png> [55] [56 <./images/fig6.png (PNG copy)>] [57 <./images/
fig7.png (PNG copy)>] [58]
File: ./images/fig7.png Graphic file (type png)
<use ./images/fig7.png> [59] [60] [61] [62] [63] [64] [65] [66] <./images/fig8.
png, id=654, 238.491pt x 121.4136pt>
File: ./images/fig8.png Graphic file (type png)
<use ./images/fig8.png> [67] [68 <./images/fig8.png (PNG copy)>] [69] [70] [71]

Overfull \hbox (0.70534pt too wide) detected at line 3196
\OT1/cmr/m/n/12 [1 + 2 + 3 + [] + (\OML/cmm/m/it/12 n \OMS/cmsy/m/n/12 ^^@ \OT1
/cmr/m/n/12 1) + \OML/cmm/m/it/12 n\OT1/cmr/m/n/12 ]\OML/cmm/m/it/12 v[]t[] \OT
1/cmr/m/n/12 = \OML/cmm/m/it/12 n \OMS/cmsy/m/n/12 ^^A [] \OML/cmm/m/it/12 v[]t
[] \OT1/cmr/m/n/12 = []\OML/cmm/m/it/12 :
 []

[72] <./images/fig9.png, id=697, 241.3818pt x 49.8663pt>
File: ./images/fig9.png Graphic file (type png)
<use ./images/fig9.png> [73] [74 <./images/fig9.png (PNG copy)>] [75] [76] [77]
[78] [79] [80] [81] [82] [83] [84] [85] [86] [87] [88]
Overfull \hbox (0.92978pt too wide) detected at line 3767
[]
 []

[89] [90] [91] [92] [93] [94] [95] [96] [97] [98] [99] [100] [101] [102] [103] 
[104] [105] [106] [107] [108] [109] [110] [111] [112] [113] [114] [115]
Overfull \hbox (5.08812pt too wide) in paragraph at lines 4726--4726
[] 
 []

[116] [117] [118] [119] [120] [121] [122] [123] [124] [125] [126] [127] [128] [
129] [130] [131] [132] [133] [134] [135] [136] [137] [138] [139] [140] [141] [1
42] [143] [144] [145] <./images/fig10.png, id=1178, 239.2137pt x 172.7253pt>
File: ./images/fig10.png Graphic file (type png)
<use ./images/fig10.png> [146] [147 <./images/fig10.png (PNG copy)>] [148] [149
] [150] [151] [152] <./images/fig11.png, id=1226, 240.6591pt x 44.8074pt>
File: ./images/fig11.png Graphic file (type png)
<use ./images/fig11.png> [153] [154 <./images/fig11.png (PNG copy)>] [155] [156
]
LaTeX Font Info:    Font shape `OMS/cmr/m/n' in size <10> not available
(Font)              Font shape `OMS/cmsy/m/n' tried instead on input line 6101.

[157


] [158] [159]
Overfull \hbox (1.97104pt too wide) in paragraph at lines 6243--6249
\OT1/cmr/m/n/12 Sturm: \OT1/cmr/m/it/12 Cours d'analyse de l'[]Ecole Poly-tech-
nique. \OT1/cmr/m/n/12 10. ^^Sedition,
 []

[160] [161] [162] [163]
Underfull \hbox (badness 1931) in paragraph at lines 6413--6420
\OT1/cmr/m/it/12 de l'[]Ecole Poly-tech-nique. \OT1/cmr/m/n/12 Tome IV: \OT1/cm
r/m/it/12 R^^Sesum^^Se des
 []

[164] [165] [166]
Underfull \hbox (badness 1881) in paragraph at lines 6565--6569
\OT1/cmr/m/n/12 Frenet: \OT1/cmr/m/it/12 Re-cueil d'exercises sur le cal-cul in
-finit^^Sesimal.
 []

[167]
Overfull \hbox (1.74687pt too wide) in paragraph at lines 6597--6601
\OT1/cmr/m/n/12 Sohncke, L. A.: \OT1/cmr/m/it/12 Samm-lung von Auf-gaben aus de
r Differential-
 []

[168] (./39041-t.ind [169] [170

] [171] [172] [173])
Overfull \hbox (11.60709pt too wide) in paragraph at lines 7277--7277
[]\OT1/cmtt/m/n/8 *** END OF THIS PROJECT GUTENBERG EBOOK DIFFERENTIAL, INTEGRA
L CALCULUS ***[] 
 []

[1


]
Overfull \hbox (3.10696pt too wide) in paragraph at lines 7344--7344
[]\OT1/cmtt/m/n/8 1.C.  The Project Gutenberg Literary Archive Foundation ("the
 Foundation"[] 
 []


Overfull \hbox (3.10696pt too wide) in paragraph at lines 7349--7349
[]\OT1/cmtt/m/n/8 located in the United States, we do not claim a right to prev
ent you from[] 
 []


Overfull \hbox (3.10696pt too wide) in paragraph at lines 7354--7354
[]\OT1/cmtt/m/n/8 freely sharing Project Gutenberg-tm works in compliance with 
the terms of[] 
 []

[2] [3]
Overfull \hbox (3.10696pt too wide) in paragraph at lines 7417--7417
[]\OT1/cmtt/m/n/8 posted on the official Project Gutenberg-tm web site (www.gut
enberg.net),[] 
 []

[4] [5] [6] [7] [8] [9] [10] (./39041-t.aux)

 *File List*
    book.cls    2007/10/19 v1.4h Standard LaTeX document class
    bk12.clo    2007/10/19 v1.4h Standard LaTeX file (size option)
inputenc.sty    2008/03/30 v1.1d Input encoding file
  latin1.def    2008/03/30 v1.1d Input encoding file
  ifthen.sty    2001/05/26 v1.1c Standard LaTeX ifthen package (DPC)
 amsmath.sty    2000/07/18 v2.13 AMS math features
 amstext.sty    2000/06/29 v2.01
  amsgen.sty    1999/11/30 v2.0
  amsbsy.sty    1999/11/29 v1.2d
  amsopn.sty    1999/12/14 v2.01 operator names
 amssymb.sty    2009/06/22 v3.00
amsfonts.sty    2009/06/22 v3.00 Basic AMSFonts support
   alltt.sty    1997/06/16 v2.0g defines alltt environment
indentfirst.sty    1995/11/23 v1.03 Indent first paragraph (DPC)
  yfonts.sty    2003/01/08 v1.3 (WaS)
footmisc.sty    2009/09/15 v5.5a a miscellany of footnote facilities
  icomma.sty    2002/03/10 v2.0 (WaS)
multicol.sty    2008/12/05 v1.6h multicolumn formatting (FMi)
 makeidx.sty    2000/03/29 v1.0m Standard LaTeX package
   array.sty    2008/09/09 v2.4c Tabular extension package (FMi)
graphicx.sty    1999/02/16 v1.0f Enhanced LaTeX Graphics (DPC,SPQR)
  keyval.sty    1999/03/16 v1.13 key=value parser (DPC)
graphics.sty    2009/02/05 v1.0o Standard LaTeX Graphics (DPC,SPQR)
    trig.sty    1999/03/16 v1.09 sin cos tan (DPC)
graphics.cfg    2009/08/28 v1.8 graphics configuration of TeX Live
  pdftex.def    2009/08/25 v0.04m Graphics/color for pdfTeX
    calc.sty    2007/08/22 v4.3 Infix arithmetic (KKT,FJ)
fancyhdr.sty    
geometry.sty    2008/12/21 v4.2 Page Geometry
   ifpdf.sty    2009/04/10 v2.0 Provides the ifpdf switch (HO)
  ifvtex.sty    2008/11/04 v1.4 Switches for detecting VTeX and its modes (HO)
geometry.cfg
hyperref.sty    2009/10/09 v6.79a Hypertext links for LaTeX
 ifxetex.sty    2009/01/23 v0.5 Provides ifxetex conditional
 hycolor.sty    2009/10/02 v1.5 Code for color options of hyperref/bookmark (HO
)
xcolor-patch.sty    2009/10/02 xcolor patch
  pd1enc.def    2009/10/09 v6.79a Hyperref: PDFDocEncoding definition (HO)
etexcmds.sty    2007/12/12 v1.2 Prefix for e-TeX command names (HO)
infwarerr.sty    2007/09/09 v1.2 Providing info/warning/message (HO)
hyperref.cfg    2002/06/06 v1.2 hyperref configuration of TeXLive
kvoptions.sty    2009/08/13 v3.4 Keyval support for LaTeX options (HO)
kvsetkeys.sty    2009/07/30 v1.5 Key value parser with default handler support 
(HO)
     url.sty    2006/04/12  ver 3.3  Verb mode for urls, etc.
  bitset.sty    2007/09/28 v1.0 Data type bit set (HO)
 intcalc.sty    2007/09/27 v1.1 Expandable integer calculations (HO)
bigintcalc.sty    2007/11/11 v1.1 Expandable big integer calculations (HO)
pdftexcmds.sty    2009/09/23 v0.6 LuaTeX support for pdfTeX utility functions (
HO)
ifluatex.sty    2009/04/17 v1.2 Provides the ifluatex switch (HO)
 ltxcmds.sty    2009/08/05 v1.0 Some LaTeX kernel commands for general use (HO)

atbegshi.sty    2008/07/31 v1.9 At begin shipout hook (HO)
 hpdftex.def    2009/10/09 v6.79a Hyperref driver for pdfTeX
supp-pdf.mkii
   color.sty    2005/11/14 v1.0j Standard LaTeX Color (DPC)
   color.cfg    2007/01/18 v1.5 color configuration of teTeX/TeXLive
 nameref.sty    2007/05/29 v2.31 Cross-referencing by name of section
refcount.sty    2008/08/11 v3.1 Data extraction from references (HO)
 39041-t.out
 39041-t.out
    umsa.fd    2009/06/22 v3.00 AMS symbols A
    umsb.fd    2009/06/22 v3.00 AMS symbols B
  omscmr.fd    1999/05/25 v2.5h Standard LaTeX font definitions
./images/fig1.png
./images/fig1.png
./images/fig2.png
./images/fig3.png
./images/fig2.png
./images/fig4.png
./images/fig5.png
./images/fig5.png
./images/fig6.png
./images/fig7.png
./images/fig7.png
./images/fig8.png
./images/fig9.png
./images/fig10.png
./images/fig11.png
 39041-t.ind
 ***********

 ) 
Here is how much of TeX's memory you used:
 7120 strings out of 493848
 95599 string characters out of 1152824
 205472 words of memory out of 3000000
 9724 multiletter control sequences out of 15000+50000
 18054 words of font info for 69 fonts, out of 3000000 for 9000
 714 hyphenation exceptions out of 8191
 37i,18n,43p,298b,640s stack positions out of 5000i,500n,10000p,200000b,50000s
</usr/share/texmf-texlive/fonts/type1/public/amsfonts/cm/cmcsc10.pfb></usr/sh
are/texmf-texlive/fonts/type1/public/amsfonts/cm/cmex10.pfb></usr/share/texmf-t
exlive/fonts/type1/public/amsfonts/cm/cmmi10.pfb></usr/share/texmf-texlive/font
s/type1/public/amsfonts/cm/cmmi12.pfb></usr/share/texmf-texlive/fonts/type1/pub
lic/amsfonts/cm/cmmi6.pfb></usr/share/texmf-texlive/fonts/type1/public/amsfonts
/cm/cmmi7.pfb></usr/share/texmf-texlive/fonts/type1/public/amsfonts/cm/cmmi8.pf
b></usr/share/texmf-texlive/fonts/type1/public/amsfonts/cm/cmr10.pfb></usr/shar
e/texmf-texlive/fonts/type1/public/amsfonts/cm/cmr12.pfb></usr/share/texmf-texl
ive/fonts/type1/public/amsfonts/cm/cmr6.pfb></usr/share/texmf-texlive/fonts/typ
e1/public/amsfonts/cm/cmr7.pfb></usr/share/texmf-texlive/fonts/type1/public/ams
fonts/cm/cmr8.pfb></usr/share/texmf-texlive/fonts/type1/public/amsfonts/cm/cmsy
10.pfb></usr/share/texmf-texlive/fonts/type1/public/amsfonts/cm/cmsy6.pfb></usr
/share/texmf-texlive/fonts/type1/public/amsfonts/cm/cmsy7.pfb></usr/share/texmf
-texlive/fonts/type1/public/amsfonts/cm/cmsy8.pfb></usr/share/texmf-texlive/fon
ts/type1/public/amsfonts/cm/cmti10.pfb></usr/share/texmf-texlive/fonts/type1/pu
blic/amsfonts/cm/cmti12.pfb></usr/share/texmf-texlive/fonts/type1/public/amsfon
ts/cm/cmtt10.pfb></usr/share/texmf-texlive/fonts/type1/public/amsfonts/cm/cmtt8
.pfb></usr/share/texmf-texlive/fonts/type1/public/amsfonts/symbols/msam10.pfb>
Output written on 39041-t.pdf (194 pages, 919386 bytes).
PDF statistics:
 1752 PDF objects out of 2073 (max. 8388607)
 523 named destinations out of 1000 (max. 500000)
 120 words of extra memory for PDF output out of 10000 (max. 10000000)

